\documentclass[11pt,a4paper]{article}
\usepackage[utf8]{inputenc}
\usepackage[T1]{fontenc}
\usepackage[top=2.5cm,bottom=2.5cm,left=2.5cm,right=2.5cm]{geometry}
\usepackage{amsmath,amssymb,amsthm,mathtools}
\usepackage[numbers,sort&compress]{natbib}
\usepackage{hyperref}
\usepackage{enumitem}
\usepackage{booktabs}
\usepackage{tikz}
\usetikzlibrary{arrows.meta,positioning,shapes.geometric,fit,decorations.pathreplacing,calc}
\usepackage{xcolor}
\usepackage{float}
\DeclareMathOperator*{\argmin}{arg\,min}
\DeclareMathOperator*{\argmax}{arg\,max}


\theoremstyle{plain}
\newtheorem{theorem}{Theorem}[section]
\newtheorem{proposition}[theorem]{Proposition}
\newtheorem{lemma}[theorem]{Lemma}
\theoremstyle{definition}
\newtheorem{definition}[theorem]{Definition}
\newtheorem{example}[theorem]{Example}
\theoremstyle{remark}
\newtheorem{remark}[theorem]{Remark}

\definecolor{br1}{RGB}{31,119,180}
\definecolor{br2}{RGB}{255,127,14}
\definecolor{br3}{RGB}{44,160,44}
\definecolor{br4}{RGB}{214,39,40}
\definecolor{br5}{RGB}{148,103,189}
\definecolor{br6}{RGB}{140,86,75}
\hypersetup{colorlinks=true,linkcolor=br1,citecolor=br3,urlcolor=br4}
\begin{document}

\title{\textbf{Stochastic Degradation-Aware Optimal Fleet Management with Application to Electric Vehicles}\\[6pt]
{\normalsize A Comprehensive Reference Document}}
\author{L.\ Fiaschi \\ \small ETH Z\"urich \\ \small \texttt{lfiaschi@ethz.ch}}
\date{May 2026}

\maketitle
\begin{abstract}
This is a project scoping where I collect all the needed information to write an article.
The goal is to develop a general purpose model and its algorithmic solver to address the management of a fleet during a certain horizon.
Each member of the fleet is described as a collection of components.
Each day, a member can be sent to serve a mission (route), to maintenance (partial or full repair of each component), or left idle.
The service of a mission damages the components of the fleet member.
Each component has its own stochastic degradation model, and the actual degradation changes according to the mission served.
We show a concrete application to the management of electric-vehicle fleets, such as buses and cars.
\end{abstract}
\tableofcontents
\newpage

\section{Overview of the problem}
\label{sec:overview}

We consider a fleet of $F$ members, each composed of $L$ components subject to stochastic degradation.
The planning horizon spans $T = H_1 + H_2$ discrete time steps (days).
The first $H_1$ days form the \emph{transitory phase}, which starts from the current state of the fleet.
The last $H_2$ days form the \emph{operating phase}, which is the schedule to repeat.
At each step, every fleet member performs at most one activity: it serves one of $M$ available missions or spends the day in the depot.
A member with neither activity is idle.
Serving a mission damages the components of the assigned member according to a stochastic degradation process that depends on both the mission and the component.
On a depot day, each component may independently undergo imperfect repair (which partially reduces accumulated damage), full replacement (which resets the component to a like-new state), or no intervention at all.

The degradation state of component $\ell$ of fleet member $i$ at time $k$ is described by a random variable $D_{i\ell k}$, whose distribution belongs to a parametric family specified per component.
A component is considered to have failed when $D_{i\ell k}$ exceeds a threshold $\tau$; the schedule must guarantee that the probability of this event remains below $\varepsilon$ at every time step.
To ensure that the schedule can be repeated indefinitely, a \emph{loop constraint} compares each component at the end of the horizon with the same component at the end of the transitory phase.
Every statistic that certifies the reliability of the component must be no larger at the end of the horizon.
So the operating phase leaves the fleet in at least as good a condition as it found it.

Table~\ref{tab:notation} collects the notation used throughout this document.
All variables are general and independent of any specific degradation model.

\begin{table}[htbp]
\centering
\caption{Notation and variable definitions.}
\label{tab:notation}
\small
\begin{tabular}{@{}lp{10.5cm}@{}}
\toprule
\textbf{Symbol} & \textbf{Description} \\
\midrule
\multicolumn{2}{@{}l}{\textit{Index sets}} \\
$i \in \{1,\ldots,F\}$        & Fleet member \\
$j \in \{0,1,\ldots,M\}$      & Activity ($j=0$: depot day; $j \ge 1$: mission; a member with no activity is idle) \\
$k \in \{1,\ldots,T\}$        & Time step (day) \\
$\ell \in \{1,\ldots,L\}$     & Component \\
\midrule
\multicolumn{2}{@{}l}{\textit{Scalar and fleet-level parameters}} \\
$F$             & Number of fleet members \\
$H_1$, $H_2$    & Lengths of the transitory phase (days $1,\ldots,H_1$) and of the operating phase (days $H_1+1,\ldots,T$) \\
$T$             & Horizon length, $T = H_1 + H_2$ \\
$M$             & Number of missions (must satisfy $F > M$) \\
$L$             & Number of components per fleet member \\
$\tau$          & Degradation failure threshold, taken uniform across components ($\tau = 1$ in the Palmgren--Miner normalisation) \\
$\varepsilon$   & Reliability threshold: $\Pr(D_{i\ell k} > \tau) \le \varepsilon$ for all $i,\ell,k$ \\
\midrule
\multicolumn{2}{@{}l}{\textit{Cost functions}} \\
$C_M$           & Depot-day cost: $\{x_{i0k}\}_{i,k} \to \mathbb{R}_{\ge 0}$ \\
$C_R$           & Repair cost: $\{z_{i\ell k}\}_{i,\ell,k} \to \mathbb{R}_{\ge 0}$ \\
$C_{\mathrm{rep}}$ & Replacement cost: $\{x^r_{i\ell k}\}_{i,\ell,k} \to \mathbb{R}_{\ge 0}$ \\
$C_D$           & Damage regularisation cost: $u \to \mathbb{R}_{\ge 0}$ \\
\midrule
\multicolumn{2}{@{}l}{\textit{Component-level parameters}} \\
$\rho_{i\ell}$      & Repair efficiency $\in (0,1]$ \\
\midrule
\multicolumn{2}{@{}l}{\textit{Initial and reference values}} \\
$D^{\mathrm{start}}_{i\ell}$       & Initial degradation state of component $\ell$ of member $i$ (known distribution; must satisfy $\Pr(D^{\mathrm{start}}_{i\ell} > \tau) \le \varepsilon$) \\
$D^{\mathrm{new}}_{i\ell}$ & Post-replacement degradation state (known distribution; typically near-zero damage) \\
\midrule
\multicolumn{2}{@{}l}{\textit{Stochastic state}} \\
$D_{i\ell k}$  & Random variable: degradation of component $\ell$ of member $i$ at time $k$ \\
$\Delta D_{i\ell k}^{(j)}$ & Random damage increment to component $\ell$ of member $i$ when mission $j$ is served on day $k$ \\
$\xi_{i\ell k}$ & Tracked state of component $\ell$ of member $i$ on day $k$: the statistics of $D_{i\ell k}$ that certify its reliability and, under $\text{ARD}_1$, their anchors (Section~\ref{sec:loop_constraint}) \\
\midrule
\multicolumn{2}{@{}l}{\textit{Decision variables}} \\
$x_{ijk} \in \{0,1\}$         & $1$ if member $i$ is assigned to activity $j$ on day $k$ \\
$x^m_{i\ell k} \in \{0,1\}$  & $1$ if component $\ell$ of member $i$ undergoes imperfect repair on day $k$ \\
$x^r_{i\ell k} \in \{0,1\}$   & $1$ if component $\ell$ of member $i$ is replaced on day $k$ \\
$z_{i\ell k} \ge 0$            & Degradation removed by imperfect repair (in expected-damage units) \\
$u \ge 0$                      & Damage regularisation variable \\
\bottomrule
\end{tabular}
\end{table}

\subsection{Scheduling and maintenance variables}

The schedule is encoded by three families of binary decision variables.
The \emph{assignment variable} $x_{ijk}$ determines which activity fleet member~$i$ performs on day~$k$: serving mission~$j \ge 1$, or spending the day in the depot ($j=0$).
Each member performs at most one activity per day,
\begin{equation}\label{eq:one_activity}
  \sum_{j=0}^{M} x_{ijk} \;\le\; 1 \qquad \forall\, i,\,k,
\end{equation}
and each mission must be covered by exactly one member,
\begin{equation}\label{eq:mission_coverage}
  \sum_{i=1}^{F} x_{ijk} \;=\; 1 \qquad \forall\, j \ge 1,\,k.
\end{equation}
A member whose assignment row is all zeros on day~$k$ is \emph{idle}: it serves no mission and is not in the depot.
Since $F > M$, exactly $M$ members serve missions at every step, and the other $F - M$ members are in the depot or idle.
Because~\eqref{eq:one_activity} is an inequality, the choice between a depot day and an idle day is free.
An idle day costs nothing, and the depot-day cost~$C_M$ prices this choice (Section~\ref{sec:cost_function}).
Idle days and depot days add no damage, because a component accumulates damage only on the days on which its member serves a mission (Section~\ref{sec:damage_dynamics}).
A depot day must be chosen explicitly, and repairs and replacements need one.

On a depot day ($x_{i0k}=1$), the intervention on each component is controlled independently by two additional binary variables.
The \emph{repair variable} $x^m_{i\ell k}$ indicates whether component~$\ell$ undergoes imperfect repair, and the \emph{replacement variable} $x^r_{i\ell k}$ indicates full replacement.
Both require a depot day and are mutually exclusive:
\begin{equation}\label{eq:maint_compat}
  x^m_{i\ell k} + x^r_{i\ell k} \;\le\; x_{i0k}
  \qquad \forall\, i,\,\ell,\,k.
\end{equation}
Several components of one member may receive interventions on the same depot day.
A component on a depot day with $x^m_{i\ell k} = x^r_{i\ell k} = 0$ receives no intervention (stays idle); its degradation state carries over unchanged.
So, on a depot day, each component is repaired, replaced, or rests without an intervention, and exactly one of the three cases holds.

\begin{remark}[One row instead of three]\label{rem:one_depot_row}
Row~\eqref{eq:maint_compat} has the same binary solutions as the three rows $x^m_{i\ell k} \le x_{i0k}$, $x^r_{i\ell k} \le x_{i0k}$ and $x^m_{i\ell k} + x^r_{i\ell k} \le 1$.
Its linear relaxation is tighter: the point $x_{i0k} = x^m_{i\ell k} = x^r_{i\ell k} = 0.5$ satisfies the three rows but not~\eqref{eq:maint_compat}.
With $x^m_{i\ell k}, x^r_{i\ell k} \ge 0$ and $x_{i0k} \le 1$, row~\eqref{eq:maint_compat} describes the convex hull of the four binary solutions $(x_{i0k}, x^m_{i\ell k}, x^r_{i\ell k}) \in \{(0,0,0), (1,0,0), (1,1,0), (1,0,1)\}$.
\end{remark}

\subsection{Damage dynamics}
\label{sec:damage_dynamics}

The degradation of each component evolves in discrete time according to the activity assigned to the fleet member that carries it.
We write $\Delta D_{i\ell k}^{(j)}$ for the \emph{random damage increment} sustained by component~$\ell$ of member~$i$ when mission~$j$ is served on day~$k$.
The distribution of each increment is assumed known (from the component's degradation model) and may depend on~$i$,~$\ell$,~$j$, and~$k$.
On days when the fleet member is idle or in the depot, no damage is accumulated.

At each time step, exactly one of the following mutually exclusive cases applies to every component $(i,\ell)$:

\begin{enumerate}[label=(\roman*)]
  \item \textbf{Mission.} Member~$i$ serves mission~$j \ge 1$ ($x_{ijk}=1$).
    The degradation grows by a mission- and component-dependent random increment:
    \begin{equation}\label{eq:dyn_mission}
      D_{i\ell k} \;=\; D_{i\ell,k-1} \;+\; \Delta D_{i\ell k}^{(j)}.
    \end{equation}

  \item \textbf{Imperfect repair.} The member is on a depot day ($x_{i0k}=1$) and component~$\ell$ undergoes imperfect repair ($x^m_{i\ell k}=1$).
    The repair reduces the accumulated damage by a fraction determined by the repair efficiency~$\rho_{i\ell}$ and the specific maintenance policy in use.
    Abstractly:
    \begin{equation}\label{eq:dyn_repair}
      D_{i\ell k} \;=\; \mathcal{R}\!\bigl(D_{i\ell,k-1},\;\rho_{i\ell}\bigr),
    \end{equation}
    where $\mathcal{R}$ denotes the repair operator.
    The formulation offers two repair operators, chosen per component.
    $\text{ARD}_\infty$ (arithmetic reduction of deterioration with infinite memory) removes the fraction~$\rho_{i\ell}$ of the whole damage.
    $\text{ARD}_1$ (arithmetic reduction of deterioration of order~1) removes the fraction~$\rho_{i\ell}$ of the damage accumulated since the last intervention, repair or replacement (Section~\ref{sec:repair_connection}).
    See Section \ref{sec:maint_strategies} for more details.
    The degradation removed by the repair action defines the repair cost variable:
    \begin{equation}\label{eq:z_def}
      z_{i\ell k} \;=\; \mathbb{E}\bigl[D_{i\ell,k-1}\bigr] - \mathbb{E}\bigl[D_{i\ell k}\bigr],
      \qquad \text{when } x^m_{i\ell k}=1.
    \end{equation}
    On days without repair the equality is not imposed and $z_{i\ell k}$ is a free non-negative variable driven to zero by the cost minimisation; equivalently, $z_{i\ell k} = 0$ when $x^m_{i\ell k} = 0$.

  \item \textbf{Full replacement.} The member is on a depot day and component~$\ell$ is replaced ($x^r_{i\ell k}=1$).
    The degradation is reset:
    \begin{equation}\label{eq:dyn_replace}
      D_{i\ell k} \;=\; D^{\mathrm{new}}_{i\ell}.
    \end{equation}

  \item \textbf{Idle day, or depot day with no intervention on this component.}
    The degradation is unchanged:
    \begin{equation}\label{eq:dyn_idle}
      D_{i\ell k} \;=\; D_{i\ell,k-1}.
    \end{equation}
\end{enumerate}

Cases~(ii) and~(iii) are mutually exclusive: a component cannot be both repaired and replaced on the same day.
Both require a depot day ($x_{i0k}=1$).
Since exactly one case applies at each step, the dynamics can be written as a single equation:
\begin{equation}\label{eq:dyn_full}
  D_{i\ell k} \;=\;
    D_{i\ell,k-1}\bigl(1 - x^m_{i\ell k}\bigr)\bigl(1 - x^r_{i\ell k}\bigr) + \sum_{j=1}^{M} x_{ijk}\,\Delta D_{i\ell k}^{(j)}
    \;+\; \mathcal{R}\!\bigl(D_{i\ell,k-1},\,\rho_{i\ell}\bigr)\,x^m_{i\ell k}
    \;+\; D^{\mathrm{new}}_{i\ell}\,x^r_{i\ell k},
\end{equation}
with initial condition $D_{i\ell,0} = D^{\mathrm{start}}_{i\ell}$.

\paragraph{Horizon and phases.}
We write the horizon as $H = [H_1, H_2]$, with $T = H_1 + H_2$ days in two phases.
A single value $H$ means $H_1 = H_2 = H$.
The transitory phase, days $1,\ldots,H_1$, starts from the current state of the fleet.
The operating phase, days $H_1+1,\ldots,T$, is the schedule to repeat without end.
The two phases may have different lengths and different increment laws.
The increment laws of the operating phase repeat with period~$H_2$.
In the \emph{endless schedule}, the transitory phase runs once and the operating phase then repeats: day $T + nH_2 + h$, for $n \ge 0$ and $1 \le h \le H_2$, carries the decisions and the increment laws of day $H_1 + h$.
When the transitory phase has no increment laws of its own, its days read the laws of the operating phase cyclically.
For example, with $H = [2, 3]$, days $1,\ldots,5$ read the operating laws $1, 2, 1, 2, 3$; with separate transitory laws, days $1, 2$ read them and days $3, 4, 5$ read the operating laws $1, 2, 3$.
With a single value $H$ and no transitory laws, both phases read the same $H$ laws, one phase after the other.

\paragraph{Degradation models of the formulation.}
Each component has one of five degradation models: Gaussian, inverse Gaussian, Wiener, Gamma or rainflow.
Section~\ref{sec:loop_per_model} lists the statistics that each model tracks and its reliability row.
The rainflow model is the distribution-free pipeline of Section~\ref{sec:cons_surrogates}.
Both the Gaussian and the Wiener model use Gaussian increments.
The Gaussian model tracks a variance that a repair reduces.
The Wiener model fits components whose uncertainty does not depend on the route and is not reduced by a repair.
When in doubt, the Gaussian model is the choice.
The rainflow model makes no distributional assumption: its certificate holds for every increment law with the measured summaries, at the price of a looser bound than an exact distribution function.

\paragraph{State after a replacement.}
A replacement installs a new part.
Its damage $D^{\mathrm{new}}_{i\ell}$ is described by the same statistics as the initial damage: a known damage when its variance is $0$, a random one otherwise.
The default is a new part with zero damage.
On a replacement day every tracked statistic takes its replacement value, and so does every anchor (Section~\ref{sec:repair_connection}).
The replacement state must satisfy every row that holds on a replacement day, in the form built at the component.
A linear form of a reliability row is stricter than the exact row (Section~\ref{sec:linear_forms}).
So a replacement state that passes only the exact row would make every replacement of the component infeasible; an input check prevents this (Section~\ref{sec:check_reasons}).

\subsection{Stochastic loop constraint}
\label{sec:loop_constraint}

For the schedule to be repeatable indefinitely, the reliability constraint $\Pr(D_{i\ell k} > \tau) \le \varepsilon$ must hold on every day of the endless schedule, not only on the $T$ days of the horizon.
So the fleet must leave the operating phase in no worse condition than it entered it.
A first formalisation compares the probabilities of failure at the start and at the end of the operating phase, component by component:
\begin{equation}\label{eq:loop_prob}
  \Pr\!\bigl(D_{i\ell T} > \tau\bigr)
  \;\le\;
  \Pr\!\bigl(D_{i\ell H_1} > \tau\bigr)
  \qquad \forall\, i,\,\ell.
\end{equation}
Together with the reliability constraint enforced at every step, condition~\eqref{eq:loop_prob} ensures that the operating phase does not consume more reliability margin than the transitory phase created.
A single-threshold comparison of two survival functions is, however, a lossy summary of the state, and by itself it does not guarantee that the inequality persists when the schedule is repeated: the non-invariance that Section~\ref{sec:loop_routeI} exhibits for bound-based certificates arises already for the exact probabilities whenever the state distribution carries more than one free parameter---for a Gaussian state, a variance--mean trade-off can satisfy~\eqref{eq:loop_prob} over one period while the failure probability $\Phi\bigl((\mu-\tau)/\sigma\bigr)$ creeps above~$\varepsilon$ at a later one.
Concatenating copies of the schedule is guaranteed to keep the failure probability below~$\varepsilon$ whenever the comparison is made on a quantity that bounds the failure probability monotonically and propagates monotonically under the dynamics~\eqref{eq:dyn_full}.
One such comparison is the dominance requirement $D_{i\ell T} \le_{\mathrm{st}} D_{i\ell H_1}$, which chains under $\text{ARD}_\infty$ because the dynamics is then a coordinatewise non-decreasing map of the state and independent randomness (Proposition~\ref{prop:st_closure}).
The formulation compares a few statistics of the state instead, one by one.

\paragraph{Tracked state.}
Each degradation model certifies the reliability of a component through a few statistics of its damage~$D_{i\ell k}$.
They are the mean for the inverse Gaussian and Gamma models, the mean and the variance for the Gaussian and Wiener models, and, for the rainflow model, the mean and, except for the Markov bound, one statistic of the chosen tail bound (Sections~\ref{sec:cons_surrogates} and~\ref{sec:loop_per_model}).
Under the repair law $\text{ARD}_1$, each statistic also has an \emph{anchor}: its value just after the last intervention (Section~\ref{sec:repair_connection}).
The \emph{tracked state} $\xi_{i\ell k}$ collects the statistics and their anchors on day~$k$.
The loop constraint is one row per coordinate of the tracked state, anchors included:
\begin{equation}\label{eq:loop}
  \xi_{i\ell T} \;\le\; \xi_{i\ell H_1}
  \qquad \forall\, i,\,\ell \quad \text{(componentwise)}.
\end{equation}
For a Gaussian component with mean $\mu$ and variance $v$, for example, the rows are $\mu_T \le \mu_{H_1}$ and $v_T \le v_{H_1}$, plus the same rows for the two anchors under $\text{ARD}_1$.

\subsubsection{One-sided encoding of the dynamics}
\label{sec:encoding}

Each day, every tracked statistic of a component follows one rule: a mission adds an increment, a repair applies the repair map of Section~\ref{sec:repair_connection}, and a replacement resets the statistic (Section~\ref{sec:damage_dynamics}).
The formulation writes each rule as a \emph{one-sided} row: an inequality ($\ge$), not an equality.
The solver may choose a value above the result of the rule, never a value below it.
So a tracked statistic is an upper bound on the value that the rule gives.
Example: with $\mu_{i\ell,k-1} = 0.3$ and a mission increment of $0.2$, the rule gives $0.5$; the rows accept $\mu_{i\ell k} = 0.5$ and also $\mu_{i\ell k} = 0.6$.
Replacement rows and anchor updates are equalities.
The rows have two forms, big-M rows and indicator constraints, with the same feasible set at binary points (Section~\ref{sec:impl_encoding}).

Fix a schedule, that is, values of every binary, and one component.
The following objects are used here and in the rest of this section.
\begin{itemize}[nosep]
  \item The \emph{full tracked state} $\xi_k$ collects the values on day~$k$ of every tracked statistic of the component and of every anchor.
    $\xi_0$ holds the initial values and the starting anchors.
  \item The \emph{day map} $\pi_k$ gives, for each coordinate, the value that the rule of day~$k$ assigns: $s + \delta s_k$ and $g$ on a day without intervention; $a\,s + (1-a)\,g$ for both $s$ and $g$ on a repair day; the replacement value $s^{\mathrm{new}}$ for both on a replacement day.
  \item The \emph{reference trajectory} is $\zeta_0 = \xi_0$, $\zeta_k = \pi_k(\zeta_{k-1})$: the rows with every inequality replaced by an equality.
  \item The \emph{true statistics} $\theta_k$ are the true values on day~$k$ of the tracked statistics (the mean, the variance, the Hoeffding budget or the cumulant generating function of the damage $D_{i\ell k}$) and of their anchors (their values just after the last intervention).
  \item The \emph{certificate set} $\mathcal{C}_{i\ell}$ is the set of values of the tracked statistics at which the reliability row of the component holds in its exact form, side rows included.
    Anchors do not enter it.
\end{itemize}

For every model of the formulation, three facts hold (Section~\ref{sec:loop_per_model} states them model by model):
\begin{enumerate}[nosep,label=(M\arabic*)]
  \item Each repair factor $a$ lies in $[0,1]$, and each increment $\delta s_k$ is non-negative.
  \item On day~$0$ the true statistics are at most the initial values and starting anchors, $\theta_0 \le \xi_0$.
    On every later day the true statistic is at most the day map applied to the true statistics of the day before: $\theta_k \le \pi_k(\theta_{k-1})$.
    Every map holds with equality, except the map of the cumulant generating function of the Chernoff bound, which is an upper bound (Proposition~\ref{prop:rf_maps}).
  \item If values $s \in \mathcal{C}_{i\ell}$ are at least the true statistics of day~$k$, then $\Pr(D_{i\ell k} > \tau) \le \varepsilon$.
\end{enumerate}
The set $\mathcal{C}_{i\ell}$ is closed, and it is a down-set: if $s \in \mathcal{C}_{i\ell}$ and $0 \le s' \le s$ componentwise, then $s' \in \mathcal{C}_{i\ell}$.
The reason is that every reliability row, written in chance-bound form, reads ``a continuous function of the tracked statistics, non-decreasing in each of them, is at most a constant'' (Table~\ref{tab:chance_rows}).

\begin{table}[htbp]
\centering
\caption{Reliability rows of the formulation in chance-bound form.
$\kappa_G = \Phi^{-1}(1-\varepsilon)$, $\kappa_C = \sqrt{(1-\varepsilon)/\varepsilon}$, $\kappa = \ln(1/\varepsilon)$, $c = \kappa b_\ell/3$.}
\label{tab:chance_rows}
\small
\begin{tabular}{@{}lll@{}}
\toprule
\textbf{Model or bound} & \textbf{Row in chance-bound form} & \textbf{Source} \\
\midrule
Gaussian, Wiener      & $\mu + \kappa_G\sqrt{v} \le \tau$ & Section~\ref{sec:loop_per_model} \\
Inverse Gaussian      & $\mu \le \tau/\bar q$ & Theorem~\ref{thm:ig_cap} \\
Gamma                 & $\mu \le A_1(\beta^*\tau,\varepsilon)/\beta^*$ & Proposition~\ref{prop:gamma_safe} \\
Rainflow, Markov      & $\mu \le \varepsilon\,\tau$ & Table~\ref{tab:five_bounds} \\
Rainflow, Cantelli    & $\mu + \kappa_C\sqrt{v} \le \tau$ & Table~\ref{tab:five_bounds} \\
Rainflow, Hoeffding   & $\mu + \sqrt{\kappa R/2} \le \tau$ & Table~\ref{tab:five_bounds} \\
Rainflow, Bernstein   & $\mu + c + \sqrt{c^2 + 2\kappa v} \le \tau$ & Table~\ref{tab:five_bounds} \\
Rainflow, Chernoff    & $K \le s\,\tau + \ln\varepsilon$ & Table~\ref{tab:five_bounds} \\
\bottomrule
\end{tabular}
\end{table}

\begin{lemma}[Monotone day maps]\label{lem:monotone}
Each day map $\pi_k$ is non-decreasing in each coordinate of the full tracked state, anchors included.
\end{lemma}
\begin{proof}
Let $\xi \le \xi'$ componentwise.
Each coordinate of $\pi_k$ has one of the forms $s + \delta s_k$, $g$, $s$, $a\,s$, $a\,s + (1-a)\,g$ or $s^{\mathrm{new}}$.
By (M1), $a \ge 0$ and $1 - a \ge 0$.
Therefore each coordinate of $\pi_k(\xi)$ is at most the same coordinate of $\pi_k(\xi')$.
\end{proof}

\begin{proposition}[Domination]\label{prop:domination}
Every feasible solution satisfies $\xi_k \ge \pi_k(\xi_{k-1})$ and $\theta_k \le \zeta_k \le \xi_k$ for $k = 1, \ldots, T$.
\end{proposition}
\begin{proof}
At binary points the rows of day~$k$ read: $s_k \ge \pi_k(\xi_{k-1})_s$ for every tracked statistic $s$, with equality on a replacement day; $g_k = s_k \ge \pi_k(\xi_{k-1})_g$ on an intervention day; $g_k = g_{k-1} = \pi_k(\xi_{k-1})_g$ on any other day.
So $\xi_k \ge \pi_k(\xi_{k-1})$.
Let $\theta_{k-1} \le \zeta_{k-1} \le \xi_{k-1}$ (true for $k = 1$ by (M2), since $\zeta_0 = \xi_0$).
Then, by (M2) and Lemma~\ref{lem:monotone},
\[
  \theta_k \le \pi_k(\theta_{k-1}) \le \pi_k(\zeta_{k-1}) = \zeta_k \le \pi_k(\xi_{k-1}) \le \xi_k .
\]
Therefore $\theta_k \le \zeta_k \le \xi_k$ for every $k$, by induction.
\end{proof}

On every day $k = 1, \ldots, T$ the tracked values satisfy the reliability row in the form chosen for the component.
Every form lies inside $\mathcal{C}_{i\ell}$ (Proposition~\ref{prop:tangent_safe}), so $\xi_k \in \mathcal{C}_{i\ell}$.
By Proposition~\ref{prop:domination}, $\theta_k \le \xi_k$, so (M3) gives $\Pr(D_{i\ell k} > \tau) \le \varepsilon$ on the days of the horizon.
Theorem~\ref{thm:perpetual} below extends this to every day of every later repetition of the operating phase.

\paragraph{Meaning of the reported values.}
The output of a solve reports the tracked values $\xi_k$, not the reference values.
By Proposition~\ref{prop:domination}, each reported statistic and anchor is an upper bound on its reference value, and so on its true statistic.
The costs computed from them are upper bounds on the costs computed from the reference values: the repair cost (Proposition~\ref{prop:repair_cost}) and the damage penalty (Section~\ref{sec:cost_function}).
Against the true costs, Section~\ref{sec:cost_function} states one exception, for a starting anchor above the true one.
The minimisation removes slack only where slack costs something.
A larger mean raises the damage penalty $C_D\,u$ where $u$ is binding, and raises the repair cost $C_R\,z$ of every later repair.
A larger variance, Hoeffding budget or cumulant generating function tightens the reliability row and so restricts the schedule.
Elsewhere the solver may return any value above the reference value.
With $C_D = 0$ the damage penalty is absent, and the reported states may carry more slack.
Slack never affects safety (Theorem~\ref{thm:perpetual}).

\subsubsection{Perpetual safety}
\label{sec:perpetual}

In the endless schedule (Section~\ref{sec:damage_dynamics}), day $T + nH_2 + h$ carries the binaries and the increment laws of day $H_1 + h$.
So every repetition applies the same day maps $\pi_{H_1+1}, \ldots, \pi_T$.
We write $\Pi_h = \pi_{H_1+h} \circ \cdots \circ \pi_{H_1+1}$ for the first $h$ days of the operating phase and $\Pi = \Pi_{H_2}$ for the whole phase.

\begin{theorem}[Perpetual safety]\label{thm:perpetual}
Let a solution satisfy every row of the problem.
Then, for every component and every day $k \ge 1$ of the endless schedule, the true statistics $\theta_k$ and the reference values $\zeta_k$ are at most a tracked value of the horizon that lies in $\mathcal{C}_{i\ell}$.
Therefore $\zeta_k \in \mathcal{C}_{i\ell}$ and $\Pr(D_{i\ell k} > \tau) \le \varepsilon$ on every day $k \ge 1$.
\end{theorem}
\begin{proof}
Fix a component.
Let $\xi_k$ be its full tracked state, and $\theta_k$ and $\zeta_k$ the true statistics and the reference trajectory of the endless schedule.
By (M2) and Lemma~\ref{lem:monotone}, the induction of Proposition~\ref{prop:domination} gives $\theta_k \le \zeta_k$ on every day $k \ge 1$.
\begin{itemize}[nosep]
  \item Days $k \le T$: $\zeta_k \le \xi_k$ (Proposition~\ref{prop:domination}) and $\xi_k \in \mathcal{C}_{i\ell}$ (Section~\ref{sec:encoding}).
  \item By Proposition~\ref{prop:domination} and Lemma~\ref{lem:monotone}, $\xi_{H_1+h} \ge \pi_{H_1+h}(\xi_{H_1+h-1}) \ge \cdots \ge \Pi_h(\xi_{H_1})$ for $h = 1, \ldots, H_2$.
    With~\eqref{eq:loop}, $\Pi(\xi_{H_1}) \le \xi_T \le \xi_{H_1}$.
  \item Let $\Pi^n(\xi_{H_1}) \le \xi_{H_1}$ (true for $n = 0$).
    Then $\Pi^{n+1}(\xi_{H_1}) = \Pi(\Pi^n(\xi_{H_1})) \le \Pi(\xi_{H_1}) \le \xi_{H_1}$ by Lemma~\ref{lem:monotone}.
    So $\Pi^n(\xi_{H_1}) \le \xi_{H_1}$ for every $n \ge 0$.
  \item Day $k = H_1 + n H_2 + h$ with $n \ge 0$ and $1 \le h \le H_2$:
    \[
      \zeta_k = \Pi_h\bigl(\Pi^n(\zeta_{H_1})\bigr)
      \le \Pi_h\bigl(\Pi^n(\xi_{H_1})\bigr)
      \le \Pi_h(\xi_{H_1})
      \le \xi_{H_1+h} \in \mathcal{C}_{i\ell}.
    \]
\end{itemize}
Therefore, on every day $k \ge 1$, $\theta_k \le \zeta_k \le s$ for a tracked value $s \in \mathcal{C}_{i\ell}$.
Since $\mathcal{C}_{i\ell}$ is a down-set, $\zeta_k \in \mathcal{C}_{i\ell}$, and (M3) gives $\Pr(D_{i\ell k} > \tau) \le \varepsilon$.
\end{proof}

The theorem holds for every model and both repair laws, because it uses only the facts (M1)--(M3) and Lemma~\ref{lem:monotone}.
It holds in both encodings and in every form of the reliability rows.
It needs no relation between $H_1$ and $H_2$: the transitory phase may be shorter or longer than the operating phase, because $H_1$ only marks the day at which the repeated block starts.
Slack in the tracked values does not weaken it: the proof compares the reference trajectory with the tracked values, whatever slack they carry.

\paragraph{Why the anchors enter the loop.}
The third step of the proof needs $\Pi(\xi_{H_1}) \le \xi_{H_1}$ in every coordinate.
Under $\text{ARD}_1$ a larger anchor gives a weaker repair, so an anchor that grows over one repetition weakens the repairs of the next one.
Example: $\rho = 0.5$; the transitory phase ends at mean $0.4$ and anchor $0.1$; the operating phase is a mission ($+0.1$), a repair, and a mission ($+0.05$).
The phase ends at mean $0.35$ and anchor $0.3$, because the repair gives $0.5 \cdot 0.5 + 0.5 \cdot 0.1 = 0.3$.
The row on the mean holds ($0.35 \le 0.4$), and the row on the anchor fails ($0.3 > 0.1$).
Each repetition raises the mean after the repair by $0.075$, so the mean grows without bound: the third repetition already peaks at $0.525$, above the checked peak $0.5$.
A loop row on the mean alone would accept this unsafe schedule.
For the same reason, the loop constraint compares the tracked statistics one by one and never a scalar function of them, such as a bound on the failure probability (Section~\ref{sec:loop_routeI} gives an example).
The loop constraint states an order between tracked statistics; it is not a stochastic order between damage laws.

\subsubsection{What the one-sided rows accept}
\label{sec:acceptance}

Call a schedule \emph{safe} for a component when the reference trajectory of the endless schedule lies in $\mathcal{C}_{i\ell}$ and in the natural ranges of the statistics on every day.
The natural range of a statistic is an upper bound of its value at every point that satisfies the reliability row (Section~\ref{sec:impl_encoding}).
The one-sided rows relate to safety as follows.
\begin{itemize}[nosep]
  \item \emph{They never accept an unsafe schedule} (Theorem~\ref{thm:perpetual}).
  \item \emph{They accept every schedule that equality rows accept}, and more.
    Equality rows force $\xi_k = \zeta_k$; Proposition~\ref{prop:acceptance}~(a) shows that such a schedule satisfies the one-sided rows too.
    The first example below shows a safe schedule that only the one-sided rows accept.
  \item \emph{They may reject a safe schedule}, as equality rows may.
    Proposition~\ref{prop:acceptance} (both parts) excludes this for a component in the exact form with one tracked statistic, the law $\text{ARD}_\infty$ and no replacement on day~$H_1$.
    Otherwise a rejection can occur in three cases: when tracked statistics move in opposite directions over the first repetition (the second example below); when the component is replaced on day~$H_1$, because the first repetition then starts from the fixed replacement state; or when a linear form is stricter than the exact row.
\end{itemize}

\begin{proposition}[Acceptance]\label{prop:acceptance}
Fix a schedule and a component, use the exact reliability form, and let the schedule be safe for the component.
Let $\bar\zeta = \zeta_{H_1}$.
\begin{enumerate}[nosep,label=(\alph*)]
  \item If $\Pi(\bar\zeta) \le \bar\zeta$, the rows of the component, the loop row included, accept the schedule with $\xi_k = \zeta_k$ for every~$k$.
  \item If $\Pi(\bar\zeta) \ge \bar\zeta$, the component has no anchor (law $\text{ARD}_\infty$), and the component is not replaced on day~$H_1$, the rows of the component accept the schedule.
\end{enumerate}
\end{proposition}
\begin{proof}
(a) The values $\xi_k = \zeta_k$ satisfy every row of day~$k$ with equality at binary points.
They lie in $\mathcal{C}_{i\ell}$ because the schedule is safe.
The loop row reads $\zeta_T = \Pi(\bar\zeta) \le \bar\zeta = \zeta_{H_1}$.

(b) Let $y_n = \Pi^n(\bar\zeta)$, the reference state at the end of the $n$-th repetition.
By Lemma~\ref{lem:monotone} and $\Pi(\bar\zeta) \ge \bar\zeta$, $y_{n+1} = \Pi(y_n) \ge \Pi(y_{n-1}) = y_n$.
The sequence is non-decreasing and lies in the natural ranges, so it converges to some $y^*$ with $y^* \ge \bar\zeta$.
Each day map is affine or constant, so $\Pi$ and its prefixes $\Pi_h$ are continuous: $\Pi(y^*) = \lim_n y_{n+1} = y^*$ and $\Pi_h(y^*) = \lim_n \Pi_h(y_n)$.
Each $\Pi_h(y_n)$ is a reference value of a later day, so it lies in the closed set $\mathcal{C}_{i\ell}$; hence $\Pi_h(y^*) \in \mathcal{C}_{i\ell}$ for $h = 1, \ldots, H_2$, and $y^* = \Pi_{H_2}(y^*) \in \mathcal{C}_{i\ell}$.
Take $\xi_k = \zeta_k$ for $k < H_1$, $\xi_{H_1} = y^*$ and $\xi_{H_1+h} = \Pi_h(y^*)$.
Day~$H_1$ is not a replacement day and the component has no anchor, so $\xi_{H_1} = y^* \ge \pi_{H_1}(\xi_{H_1-1})$ satisfies the rows of day~$H_1$.
The later days satisfy their rows with equality, and the loop row reads $\xi_T = \Pi(y^*) = y^* = \xi_{H_1}$.
\end{proof}

A scalar state is always comparable with its image.
So a component in the exact form with one tracked statistic under $\text{ARD}_\infty$, not replaced on day~$H_1$, accepts every safe schedule.
Under $\text{ARD}_1$ the anchor is a second tracked statistic, and its value on day~$H_1$ is fixed by the anchor rows; part~(b) does not apply, and acceptance of every safe schedule is not proven.

\begin{example}[One statistic, accepted only by the one-sided rows]
One tracked mean, law $\text{ARD}_\infty$, $\rho = 0.5$.
The operating phase is a mission ($\delta\mu = 0.2$) followed by a repair, so $\Pi(y) = 0.5\,(y + 0.2) = 0.5\,y + 0.1$, and $\Pi(y) \le y$ if and only if $y \ge 0.2$.
Let $\bar\zeta = 0.1$.
Then $\Pi(0.1) = 0.15 > 0.1$, so equality rows reject the schedule.
The one-sided rows accept it with $\xi_{H_1} = 0.2$: the tracked cycle is $0.2, 0.4, 0.2$, and the reliability row is checked at $0.4$.
The reference trajectory is $0.1, 0.3, 0.15, 0.35, 0.175, \ldots$; its start-of-cycle values rise to $0.2$ and its peaks rise to $0.4$.
So the schedule is safe whenever $0.4 \in \mathcal{C}_{i\ell}$.
The price is in the reported cost: each repair is charged $\rho \cdot 0.4 = 0.2$, while the reference removal rises from $0.15$ toward $0.2$.
\end{example}

\begin{example}[Two statistics that move in opposite directions, rejected]
Cantelli bound, $\varepsilon = 0.1$ ($\kappa_C^2 = 9$), $\tau = 1$, $\rho = 0.5$, law $\text{ARD}_\infty$.
The operating phase is a mission, which adds $0.2$ to the mean and $0.009$ to the variance, followed by a repair.
So $\Pi(\mu, v) = (0.5\,\mu + 0.1,\; 0.25\,v + 0.00225)$, with fixed point $(0.2, 0.003)$.
Let $\bar\zeta = (0.5, 0)$.
\begin{itemize}[nosep]
  \item Reference: over the repetitions the mean falls to $0.2$ and the variance rises to $0.003$.
    The first peak is $(0.7, 0.009)$, with $9 \cdot 0.009 = 0.081 \le 0.3^2 = 0.09$.
    Every later peak has a mean of at most $0.55$ and a variance below $0.012$, with $9 \cdot 0.012 = 0.108 \le 0.45^2 = 0.2025$.
    The schedule is safe.
  \item One-sided rows: the loop row needs $\xi_{H_1} = y \ge (0.5, 0)$ with $\Pi(y) \le y$, so $y_v \ge 0.003$.
    The tracked peak is then at least $(0.7, 0.012)$, where $9 \cdot 0.012 = 0.108 > 0.09$.
    Rejected.
  \item Equality rows: $\Pi(\bar\zeta) = (0.35, 0.00225)$ has a larger variance than $\bar\zeta$.
    Rejected as well.
\end{itemize}
\end{example}

\subsubsection{Replacements required by the loop constraint}
\label{sec:loop_replacements}

A component may be replaced on any depot day of its member, any number of times, and a replacement costs $C_{\mathrm{rep},i\ell}$ and nothing else (Section~\ref{sec:cost_function}).
The loop constraint can require at least one replacement in the operating phase, for two reasons.
The first is the Wiener variance, which falls only by replacement (Section~\ref{sec:loop_per_model}).
The second is the repair law $\text{ARD}_1$.

\begin{proposition}[Replacement under $\text{ARD}_1$]\label{prop:ard1_replacement}
Let a component follow $\text{ARD}_1$ with $\rho < 1$, and let its member serve, in the operating phase, at least one mission with a positive mean increment for the component.
Then every schedule that satisfies the loop constraint~\eqref{eq:loop} replaces the component at least once in the operating phase.
\end{proposition}
\begin{proof}
Suppose that the component is not replaced in the operating phase.
Let $\mu_k$ be its tracked mean, $g_k$ the anchor of the mean, and $\delta\mu_k$ the mean increment of day~$k$.
Then $\mu_k \ge g_k$ on every day: with equality on a repair day, and on other days because $\mu$ does not decrease while $g$ stays constant, with $g_0 \le \mu_0$.
On a repair day~$k$ of the operating phase,
\[
  g_k = \mu_k \;\ge\; (1-\rho)\,\mu_{k-1} + \rho\, g_{k-1} \;=\; g_{k-1} + (1-\rho)\,(\mu_{k-1} - g_{k-1}).
\]
So $g_T \ge g_{H_1} + (1-\rho)\sum (\mu_{k-1} - g_{k-1})$, the sum over the repair days of the operating phase, with non-negative terms.
The loop row $g_T \le g_{H_1}$ and $\rho < 1$ force every term to~$0$.
Then every repair day gives $\mu_k \ge \mu_{k-1}$, every other day gives $\mu_k \ge \mu_{k-1} + \delta\mu_k$, and
\[
  \mu_T \;\ge\; \mu_{H_1} + \sum_{k=H_1+1}^{T} \delta\mu_k \;>\; \mu_{H_1},
\]
which violates the loop row $\mu_T \le \mu_{H_1}$.
\end{proof}

The requirement is not an artefact of the formulation.
Without a replacement, each repair leaves the share $1-\rho$ of the damage since the last intervention in place.
So the anchor rises with every repair, and the mean grows without bound over the repetitions of the operating phase.
Example: $\rho = 0.5$, and a mission adds $0.2$ to the mean between two repairs; each repair leaves the mean $0.1$ higher than the previous one.
With $\rho = 1$ a repair returns the state to the anchor, and no replacement is needed.
When replacements are not allowed, the member of such a component cannot serve, in the operating phase, a mission with a positive mean increment for the component.
If the fleet cannot cover the missions without such a mission, the problem is infeasible.

\subsection{Cost function}
\label{sec:cost_function}

The objective is to minimise the total cost of operating the fleet over the planning horizon.
Four terms contribute to the cost:
\begin{align}\label{eq:objective}
\min \quad
  C_M \left(\{x_{i0k}\}_{i,k}\right)
  \;+\; C_{R} \left(\{z_{i\ell k}\}_{i,\ell,k}\right)
  \;+\; C_{\mathrm{rep}}\, (\{x^r_{i\ell k}\}_{i,\ell,k}) +\; C_D\,(u).
\end{align}

All four cost functions are assumed to be non-negative and non-decreasing in each of their arguments: scheduling more depot days, removing more degradation, replacing more components, or tolerating higher aggregate damage never decreases the cost.

The first term charges for every depot day scheduled across the fleet.
It is paid once per member and depot day ($x_{i0k} = 1$), whatever the number of components repaired or replaced on that day.
It prices the visit to the depot, for example the member out of service and the set-up of the work.
An idle day costs nothing.
The second term accounts for imperfect repair: whenever a component is repaired, the cost grows with the amount of degradation removed~$z_{i\ell k}$.
The third term charges for every full replacement.

The last term is a damage regularisation penalty.
The variable~$u$ captures the worst-case aggregate expected damage across all fleet members and time steps:
\begin{equation}\label{eq:u_def}
  u \;=\; \max_{i,k}\; \phi\!\Bigl(\mathbb{E}[D_{i1k}],\;\ldots,\;\mathbb{E}[D_{iLk}]\Bigr),
\end{equation}
where $\phi\colon \mathbb{R}^L \to \mathbb{R}_{\ge 0}$ is a non-decreasing aggregation function.
The formulation offers two penalty types.
The default type takes $\phi$ as the sum of the expected component damages, which gives one linear row per member and day.
The quadratic type takes $\phi$ as the sum of their squares, which gives one convex quadratic row per member and day.
The quadratic type is used only with the exact form of the reliability rows, because the linear forms of those rows (Section~\ref{sec:linear_forms}) serve to keep the whole model linear.
The cost controls the trade-off between minimising operational costs and spreading damage evenly across the fleet.
The term prices the worst member-level damage, so the minimisation spreads damage across the fleet.
It also removes slack from the tracked means where $u$ is binding (Section~\ref{sec:encoding}).
Since the tracked means satisfy $\mu_{i\ell k} \ge \zeta^{\mu}_{i\ell k} \ge 0$ (Proposition~\ref{prop:domination}), both penalty types computed from them are upper bounds on the same quantity computed from the reference values.
The coefficient satisfies $C_D \ge 0$, because a negative $C_D$ would reward damage.

A variable that feeds only a cost term with a zero coefficient is not built, and neither are the rows that define it.
For example, with $C_D = 0$ the variable $u$ and its rows are dropped.
At a component whose repair-cost coefficient is zero, the variables $z_{i\ell k}$ and their rows are dropped.

\paragraph{Objective modes.}
Let $c_k$ be the cost of day~$k$ (the depot-day, repair and replacement terms of that day), $J_{\mathrm{init}} = \sum_{k \le H_1} c_k$ the cost of the transitory phase and $J_{\mathrm{op}} = \sum_{k > H_1} c_k$ the cost of the operating phase.
The formulation offers three objectives.
\begin{itemize}[nosep]
  \item \emph{Total}, $J_{\mathrm{init}} + J_{\mathrm{op}} + C_D\,u$: the cost of the $T$ modelled days.
  \item \emph{Operating average}, $J_{\mathrm{op}}/H_2 + C_D\,u$: the cost per day of the operating phase.
    It is the long-run cost per day of the endless schedule, in which the cost of the transitory phase spreads over infinitely many days.
    It is comparable across values of~$H_2$, whereas a raw objective over $T$ days is not.
    The transitory phase has no cost in this objective, but it is still constrained: the solver chooses its schedule for feasibility only.
  \item \emph{Evaluation total}, $J_{\mathrm{init}} + \frac{E - H_1}{H_2}\,J_{\mathrm{op}} + C_D\,u$: the projected cost over an evaluation horizon of $E$ days, the transitory phase once and then $(E - H_1)/H_2$ repetitions of the operating phase, a fraction of a repetition counted pro rata.
    It also prices the transitory phase, so it compares solves with different~$H_1$.
    It needs $E > H_1$: for a smaller~$E$ the operating cost would have a negative weight.
\end{itemize}
The damage penalty is added once in every mode, because $u$ is the maximum over the whole horizon.
Example: $H_1 = 5$, $H_2 = 10$, $E = 365$.
Then the evaluation total weights the operating cost by $36$, and the operating average by $1/10$.
In the operating average, the repair-cost variable is not built on the transitory days, because their cost does not enter the objective; the repair cost of those days is computed after the solve, from the tracked mean and its anchor.
The objective value is computed from the tracked values, so it is an upper bound on the cost of the schedule computed from the reference values (Section~\ref{sec:encoding}).

\paragraph{Repair cost.}
On a repair day, $z_{i\ell k}$ is the expected degradation that the repair removes (Eq.~\eqref{eq:z_def}).
The formulation writes it with the tracked mean $\mu$ and its anchor $g^{\mu}$, with $g^{\mu} \equiv 0$ under $\text{ARD}_\infty$:
\begin{equation}\label{eq:z_row}
  z_{i\ell k} \;\ge\; \rho\,\bigl(\mu_{i\ell,k-1} - g^{\mu}_{i\ell,k-1}\bigr) - \rho\,\tau\,\bigl(1 - x^m_{i\ell k}\bigr),
  \qquad z_{i\ell k} \ge 0 .
\end{equation}
On a repair day the first row reads $z_{i\ell k} \ge \rho\,(\mu_{i\ell,k-1} - g^{\mu}_{i\ell,k-1})$, the amount that the repair map removes from the mean.
On any other day its right-hand side is at most $0$, because $0 \le \mu_{i\ell,k-1} - g^{\mu}_{i\ell,k-1} \le \mu_{i\ell,k-1} \le \tau$.
A replacement day has $x^m_{i\ell k} = 0$, so $z_{i\ell k} = 0$ at the optimum, and a replacement costs $C_{\mathrm{rep},i\ell}$ only.
The replaced part is discarded, so its damage does not change the price of the new part.
The variable $z_{i\ell k}$ appears in no row other than~\eqref{eq:z_row}, and in the objective with a positive coefficient.
The minimisation therefore drives it to the largest of its lower bounds.
The upper rows of a McCormick envelope, $z \le \rho\,(\mu - g^{\mu})$ and $z \le \rho\,\tau\,x^m$, would only remove points that are never optimal, so they are not built.
A solution that is not optimal (for example, the best one found at the time limit) may carry a larger~$z$; its reported cost is then larger, never smaller.

\begin{proposition}[What $z$ measures]\label{prop:repair_cost}
In every feasible solution, on every repair day~$k$ of a component, under both laws,
\[
  z_{i\ell k} \;\ge\; \rho\,\bigl(\mu_{k-1} - g^{\mu}_{k-1}\bigr) \;\ge\; \rho\,\bigl(\zeta^{\mu}_{k-1} - \zeta^{g}_{k-1}\bigr),
\]
where $\zeta^{\mu}$ and $\zeta^{g}$ are the mean and its anchor on the reference trajectory (Section~\ref{sec:encoding}).
When the starting anchor equals the true mean level after the last intervention before the horizon, the right-hand side is the expected degradation that the repair removes.
\end{proposition}
\begin{proof}
The first inequality is~\eqref{eq:z_row} with $x^m_{i\ell k} = 1$.
Under $\text{ARD}_\infty$, $g^{\mu} = \zeta^{g} = 0$, and $\mu_{k-1} \ge \zeta^{\mu}_{k-1}$ by Proposition~\ref{prop:domination}.
Under $\text{ARD}_1$, let $t < k$ be the last intervention day before~$k$, or $t = 0$ if there is none.
Then $g^{\mu}_{k-1} = g^{\mu}_t$ and $\zeta^{g}_{k-1} = \zeta^{g}_t$.
\begin{itemize}[nosep]
  \item If $t = 0$: $g^{\mu}_0 = \zeta^{g}_0$ (the starting anchor), so $g^{\mu}_{k-1} - \zeta^{g}_{k-1} = 0 \le \mu_{k-1} - \zeta^{\mu}_{k-1}$.
  \item If $t \ge 1$: $g^{\mu}_t = \mu_t$ and $\zeta^{g}_t = \zeta^{\mu}_t$.
    On the days $j = t+1, \ldots, k-1$ there is no intervention: $\mu_j \ge \mu_{j-1} + \delta\mu_j$ and $\zeta^{\mu}_j = \zeta^{\mu}_{j-1} + \delta\mu_j$.
    So $\mu_j - \zeta^{\mu}_j$ does not decrease, and $\mu_{k-1} - \zeta^{\mu}_{k-1} \ge \mu_t - \zeta^{\mu}_t = g^{\mu}_{k-1} - \zeta^{g}_{k-1}$.
\end{itemize}
Rearranging gives $\mu_{k-1} - g^{\mu}_{k-1} \ge \zeta^{\mu}_{k-1} - \zeta^{g}_{k-1}$.
The mean recursion is linear (Section~\ref{sec:repair_connection}); with an exact starting anchor, $\zeta^{\mu}$ and $\zeta^{g}$ are the true means, so $\rho\,(\zeta^{\mu}_{k-1} - \zeta^{g}_{k-1}) = \mathbb{E}[D^- - D^+]$.
\end{proof}

A starting anchor above the true one, which the default may be (Section~\ref{sec:repair_connection}), models a weaker repair.
For the first repair before any intervention of the horizon, the charged removal then follows the starting anchor and may be below the true removal.
From the first intervention on, the reference gap $\zeta^{\mu} - \zeta^{g}$ equals the true gap, the sum of the mean increments since that intervention, so the bound holds again.
Under both laws, then, the reported repair cost $\sum C_R\, z$ is an upper bound on the expected repair cost of the schedule, with this exception for a starting anchor above the true one.
Under $\text{ARD}_1$ the tracked mean and the tracked anchor are both upper bounds, and their difference is still an upper bound on the reference difference: the anchor carries the slack that the mean had at the last intervention, and the slack of the mean does not decrease until the next one.
The bound is exact when the tracked mean has no slack.

\subsection{Abstract problem formulation}

Collecting the objective~\eqref{eq:objective}, the scheduling constraints~\eqref{eq:one_activity}--\eqref{eq:maint_compat}, the degradation dynamics~\eqref{eq:dyn_full}, the reliability constraint, and the loop constraint~\eqref{eq:loop}, the fleet management problem reads:

\begin{subequations}\label{eq:abstract}
\begin{align}
\min_{x,\,x^m,\,x^r,\,z,\,u} \quad
  & C_M\!\bigl(\{x_{i0k}\}\bigr)
    + C_R\!\bigl(\{z_{i\ell k}\}\bigr)
    + C_{\mathrm{rep}}\!\bigl(\{x^r_{i\ell k}\}\bigr)
    + C_D(u)
    \label{eq:abs_obj} \\[6pt]
%
\text{s.t.} \quad
  & \sum_{j=0}^{M} x_{ijk} \le 1
    && \forall\, i,\,k,
    \label{eq:abs_assign} \\
  & \sum_{i=1}^{F} x_{ijk} = 1
    && \forall\, j \ge 1,\,k,
    \label{eq:abs_cover} \\
  & x^m_{i\ell k} + x^r_{i\ell k} \le x_{i0k}
    && \forall\, i,\,\ell,\,k,
    \label{eq:abs_maint} \\[6pt]
%
  & D_{i\ell k} = D_{i\ell,k-1}(1-x^m_{i\ell k})(1-x^r_{i\ell k})
    + \sum_{j=1}^{M} x_{ijk}\,\Delta D_{i\ell k}^{(j)} + \notag \\
  & \qquad\qquad + \mathcal{R}(D_{i\ell,k-1},\rho_{i\ell})\,x^m_{i\ell k}
    + D^{\mathrm{new}}_{i\ell}\,x^r_{i\ell k}
    && \forall\, i,\,\ell,\,k,
    \label{eq:abs_dyn} \\
  & D_{i\ell,0} = D^{\mathrm{start}}_{i\ell}
    && \forall\, i,\,\ell,
    \label{eq:abs_init} \\[6pt]
%
  & \Pr(D_{i\ell k} > \tau) \le \varepsilon
    && \forall\, i,\,\ell,\,k,
    \label{eq:abs_rel} \\
  & \xi_{i\ell T} \le \xi_{i\ell H_1} \quad \text{(componentwise)}
    && \forall\, i,\,\ell,
    \label{eq:abs_loop} \\[6pt]
%
  & u = \max_{i,k}\;\phi\!\bigl(\mathbb{E}[D_{i1k}],\ldots,\mathbb{E}[D_{iLk}]\bigr),
    \label{eq:abs_u} \\
  & x_{ijk},\, x^m_{i\ell k},\, x^r_{i\ell k} \in \{0,1\},\quad
    z_{i\ell k} \ge 0,\quad u \ge 0.
    \label{eq:abs_domain}
\end{align}
\end{subequations}
Here $\xi_{i\ell k}$ is the tracked state of Section~\ref{sec:loop_constraint}.
When $C_D = 0$, the variable $u$ and row~\eqref{eq:abs_u} are not built (Section~\ref{sec:cost_function}).
The constraints have three groups.
The fleet rows (assignment and depot days) do not depend on the degradation models.
The common rules of the dynamics (repair laws, one-sided encoding, replacement, linear forms of the reliability rows) hold for every model, and each model adds its tracked statistics, day maps and reliability row (Section~\ref{sec:loop_per_model}).
The repair cost, the damage regularisation and the loop constraint close the problem.
All rows hold for every day $k = 1, \ldots, T$ unless a section says otherwise.

Problem~\eqref{eq:abstract} is a stochastic combinatorial optimisation problem.
It makes no commitment to a specific degradation model (the distribution of~$\Delta D_{i\ell k}^{(j)}$), repair policy (the operator~$\mathcal{R}$), aggregation function~($\phi$), or cost structure (the functions $C_M, C_R, C_{\mathrm{rep}}, C_D$).
All of these are instantiated when the formulation is specialised to a concrete fleet (Section~\ref{sec:evs}).

\section{Literature review}

\subsection{Stochastic degradation and fatigue models}

Stochastic models of component degradation fall into two broad families.
\emph{Damage-accumulation models} track a cumulative damage variable~$X(t)$ whose distribution evolves over time; failure occurs when $X(t)$ exceeds a threshold.
\emph{Remaining-life models} characterise the number of load cycles a component can sustain before failure, relating applied stress to fatigue life through empirical laws.

\subsubsection{Damage-accumulation processes}

Table~\ref{tab:deg_models} summarises four widely used stochastic processes for modelling cumulative damage.
In each case $X(t)$ denotes the cumulative damage at time~$t$ with $X(0)=0$; all processes have independent, stationary increments (L\'evy processes).

\begin{table}[htbp]
\centering
\caption{Stochastic damage-accumulation models.
Parameters: $\alpha,\beta>0$ (Gamma shape rate / inverse scale), $\mu$ (drift), $\sigma^2$ (diffusion coefficient), $\lambda$ (IG shape rate), $\nu$ (Poisson intensity), $F_Y$ (jump-size distribution).}
\label{tab:deg_models}
\small
\begin{tabular}{@{}p{3cm}p{3.6cm}p{2.2cm}cp{3.5cm}@{}}
\toprule
\textbf{Model} & \textbf{Increment} $X(t)-X(s)$ & \textbf{Marginal} $X(t)$ & \textbf{Mon.} & \textbf{Closure under $+$} \\
\midrule
Gamma process \cite{AbdelHameed1975}
  & $\mathrm{Ga}\!\bigl(\alpha(t{-}s),\,\beta\bigr)$
  & $\mathrm{Ga}(\alpha t,\,\beta)$
  & Yes
  & Same rate~$\beta$ \\[4pt]
Wiener process \cite{WhitmoreSchenkelberg1997}
  & $\mathcal{N}\!\bigl(\mu(t{-}s),\,\sigma^2(t{-}s)\bigr)$
  & $\mathcal{N}(\mu t,\,\sigma^2 t)$
  & No
  & Always \\[4pt]
Inverse Gaussian process \cite{YeChen2014}
  & $\mathrm{IG}\!\bigl(\mu(t{-}s),\,\lambda(t{-}s)^2\bigr)$
  & $\mathrm{IG}(\mu t,\,\lambda t^2)$
  & Yes
  & Same $\mu^2/\lambda$ \\[4pt]
Compound Poisson \cite{EsaryMarshallProschan1973,Wang2022}
  & $\sum_{n=N(s)+1}^{N(t)} Y_n$
  & $\mathrm{CP}(\nu t,\,F_Y)$
  & If $Y{\ge}0$
  & Always \\
\bottomrule
\end{tabular}
\end{table}

\paragraph{Gamma process.}
Introduced by Abdel-Hameed~\cite{AbdelHameed1975} and surveyed in depth by van~Noortwijk~\cite{vanNoortwijk2009}, the Gamma process is the most widely used model for monotonic degradation (wear, corrosion, crack growth).
The process $\{X(t)\}$ has non-negative independent increments $X(t)-X(s) \sim \mathrm{Ga}(\alpha(t-s),\,\beta)$, where $\alpha>0$ is the shape rate and $\beta>0$ is the inverse scale (rate).
Its mean and variance at time~$t$ are $\mathbb{E}[X(t)] = \alpha t/\beta$ and $\mathrm{Var}[X(t)] = \alpha t/\beta^2$.
The Gamma family is closed under convolution when the rate parameter~$\beta$ is shared: if $X_1 \sim \mathrm{Ga}(\alpha_1,\beta)$ and $X_2 \sim \mathrm{Ga}(\alpha_2,\beta)$ are independent, then $X_1+X_2 \sim \mathrm{Ga}(\alpha_1+\alpha_2,\beta)$.

\paragraph{Wiener process (Brownian motion with drift).}
The Wiener degradation model~\cite{WhitmoreSchenkelberg1997} sets $X(t) = \mu t + \sigma W(t)$, where $W(t)$ is a standard Brownian motion, $\mu>0$ the drift rate, and $\sigma>0$ the diffusion coefficient.
Increments are normally distributed: $X(t)-X(s) \sim \mathcal{N}(\mu(t-s),\,\sigma^2(t-s))$.
The normal family is unconditionally closed under addition.
Unlike the other three models, the Wiener process is \emph{not} monotone: sample paths can decrease, making it suitable for degradation phenomena with short-term fluctuations.
A notable property is that the first-passage time to a fixed threshold~$\tau$ follows an Inverse Gaussian distribution: $T_\tau \sim \mathrm{IG}(\tau/\mu,\;\tau^2/\sigma^2)$.

\paragraph{Inverse Gaussian process.}
Proposed as a degradation model by Ye and Chen~\cite{YeChen2014}, the IG process has non-negative increments $X(t)-X(s) \sim \mathrm{IG}(\mu(t-s),\,\lambda(t-s)^2)$, with mean $\mu t$ and variance $\mu^3 t/\lambda$.
Its increments are heavier-tailed than those of the Gamma process, making it appropriate when occasional large damage jumps are observed.
The IG distribution is closed under convolution when the dispersion parameter $\mu^2/\lambda$ is shared across summands; however, the sum of two independent IG \emph{processes} with different parameters is not itself an IG process~\cite{YeChen2014}.

\paragraph{Compound Poisson process.}
The compound Poisson process $X(t) = \sum_{n=1}^{N(t)} Y_n$ models shock-based degradation~\cite{EsaryMarshallProschan1973}, where shocks arrive according to a Poisson process $N(t)$ with rate~$\nu$ and each shock causes a random damage~$Y_n$ drawn i.i.d.\ from a distribution~$F_Y$.
When $Y_n \ge 0$, the process is monotone.
The sum of independent compound Poisson processes is always a compound Poisson process.
This model is natural for discrete damage events (e.g.\ impacts, overloads) rather than continuous wear; see Wang~\cite{Wang2022} for a recent application to structural reliability.

\subsubsection{Discrete-time formulation}
\label{sec:discrete_deg}

The continuous-time processes of Table~\ref{tab:deg_models} are sampled at discrete steps in the fleet management setting of Section~\ref{sec:overview}: each day $k$, the damage of component~$(i,\ell)$ grows by a random increment $\Delta D_{i\ell k}^{(j)}$ whose distribution depends on the mission~$j$ served that day.
The cumulative damage after $k$~steps is therefore a random walk
\begin{equation}\label{eq:discrete_damage}
  D_{i\ell k} \;=\; D_{i\ell,0} \;+\; \sum_{t=1}^{k} \Delta D_{i\ell t},
\end{equation}
where the increments $\Delta D_{i\ell t}$ are independent but \emph{not} identically distributed (different missions may be served on different days).
Table~\ref{tab:discrete_models} restates the four models in this discrete-time setting.

\begin{table}[htbp]
\centering
\caption{Discrete-time specialisation of the damage-accumulation models.
Each row gives the distribution of a single-step increment~$\Delta D$ and the resulting marginal of $D_k = \sum_{t=1}^{k}\Delta D_t$ when all increments share the parameters required for closure.}
\label{tab:discrete_models}
\small
\begin{tabular}{@{}p{2.1cm}p{3.4cm}p{4.0cm}p{4.2cm}@{}}
\toprule
\textbf{Model} & \textbf{Increment} $\Delta D_t$ & \textbf{Marginal} $D_k$ (homogeneous) & \textbf{Closure condition} \\
\midrule
Gamma
  & $\mathrm{Ga}(\alpha_t,\,\beta)$
  & $\mathrm{Ga}\!\bigl(\sum_{t}\alpha_t,\;\beta\bigr)$
  & Same rate~$\beta$ \\[4pt]
Gaussian
  & $\mathcal{N}(\mu_t,\,\sigma_t^2)$
  & $\mathcal{N}\!\bigl(\sum_{t}\mu_t,\;\sum_{t}\sigma_t^2\bigr)$
  & Always \\[4pt]
Inv.\ Gaussian
  & $\mathrm{IG}(\mu_t,\,\lambda_t)$
  & $\mathrm{IG}\!\bigl(\sum_{t}\mu_t,\;\frac{(\sum_t \mu_t)^3}{ \sum_t(\mu_t^3/\lambda_t)}\bigr)$
  & Same $\mu^2/\lambda$ \\[4pt]
Compound Poisson
  & $\sum_{n=1}^{N_t} Y_n$
  & $\mathrm{CP}\!\bigl(\sum_t\nu_t,\,F_Y\bigr)$
  & Same $F_Y$ \\
\bottomrule
\end{tabular}
\end{table}

The closure conditions determine whether the marginal distribution of~$D_{i\ell k}$ remains within the same parametric family after summing increments from different missions.
When the condition holds, the distribution of~$D_{i\ell k}$ is available in closed form as a function of the cumulative parameters, which is essential for evaluating the reliability constraint $\Pr(D_{i\ell k} > \tau) \le \varepsilon$ efficiently.
When it does not hold---for instance, summing Gamma increments with different rate parameters---the marginal must be computed numerically or approximated.

In the fleet management problem of Section~\ref{sec:overview}, each day brings an increment with mission-dependent parameters.
Closure therefore requires that the parameter shared across missions is indeed constant for a given component: the rate~$\beta$ for Gamma, or the dispersion $\mu^2/\lambda$ for Inverse Gaussian.
The shape (Gamma) and mean/variance (Gaussian) parameters may vary freely across missions and accumulate additively; for the Inverse Gaussian the means accumulate additively and the cumulative shape is induced by the shared dispersion, $\lambda = (\sum_t \mu_t)^2\big/(\mu^2/\lambda)$.
For the Gaussian model, no restriction is needed: normal random variables are unconditionally closed under addition.

\subsubsection{Remaining-life models: the S--N curve and Palmgren--Miner rule}

The stress--life (S--N) approach, originating with W\"ohler's experiments in the 1850s, relates the stress amplitude~$\sigma_a$ to the number of cycles to failure~$N_f$.
The most common analytical form is Basquin's law~\cite{Basquin1910}:
\begin{equation}\label{eq:basquin}
  \sigma_a^{\,m}\, N_f \;=\; C,
\end{equation}
where $m>0$ is the fatigue exponent and $C>0$ is a material constant, both typically fitted from experimental data.
Equivalently, $N_f = C\,\sigma_a^{-m}$, which appears as a straight line on a log--log plot.

Under variable-amplitude loading, the Palmgren--Miner linear damage rule~\cite{Palmgren1924,Miner1945} defines the cumulative damage fraction as
\begin{equation}\label{eq:miner}
  D_{\mathrm{PM}} \;=\; \sum_{i} \frac{n_i}{N_{f,i}},
\end{equation}
where $n_i$ is the number of applied cycles at stress level~$\sigma_{a,i}$ and $N_{f,i}$ is the corresponding life from the S--N curve.
Failure is predicted when $D_{\mathrm{PM}} \ge 1$.
The rule assumes that damage accumulates linearly and independently of load ordering---an approximation that has been questioned experimentally (the critical damage sum $D_{\mathrm{PM}}^*$ at failure is scattered around~1 rather than being exactly~1) and extended stochastically by treating $D_{\mathrm{PM}}$ and/or $D_{\mathrm{PM}}^*$ as random variables~\cite{Shen1980,CastilloFernandezCanteli2009,Hilow2017,BohmEtAl2021,MunizCalvente2022}.

Two sources of randomness enter the Palmgren--Miner accumulation:
\begin{enumerate}[label=(\roman*)]
  \item \emph{S--N scatter.}
    At a given stress level~$\sigma_a$, the fatigue life~$N_f$ is not a deterministic constant but a random variable whose distribution reflects material variability, surface finish, and testing conditions.
    The increment $\Delta D_{\mathrm{PM}} = n/N_f$ therefore inherits this randomness.
  \item \emph{Random critical damage sum.}
    Experimental evidence shows that the critical value $D_{\mathrm{PM}}^*$ at failure is itself scattered (typically in the range $0.5$--$2.0$) rather than being exactly~1~\cite{Shen1980}.
    This can be modelled by replacing the deterministic threshold with a random one, or equivalently by normalising the damage by a random capacity.
\end{enumerate}
In the fleet management framework of Section~\ref{sec:overview}, source~(i) is the relevant one: the threshold is fixed at $\tau = 1$ and the randomness enters through the increments.
We therefore focus on the distributional assumptions for~$N_f$ and the resulting distribution of~$\Delta D_{\mathrm{PM}}$.
Table~\ref{tab:sn_scatter} summarises the three most common choices.

\begin{table}[htbp]
\centering
\caption{Distributional assumptions for S--N scatter and their effect on the Palmgren--Miner increment $\Delta D_{\mathrm{PM}} = n/N_f$.
In all cases, $\mu(\sigma_a)$ and the scale/shape parameters depend on the stress level through Basquin's law~\eqref{eq:basquin}.}
\label{tab:sn_scatter}
\small
\begin{tabular}{@{}p{2.1cm}p{3.4cm}p{3.5cm}p{1.2cm}p{3.0cm}@{}}
\toprule
\textbf{Model} & \textbf{Life} $N_f$ & \textbf{Increment} $\Delta D_{\mathrm{PM}}$ & \textbf{Mon.} & \textbf{Closure under $+$} \\
\midrule
Log-normal
  & $\ln N_f \sim \mathcal{N}\!\bigl(\mu,\,\sigma^2\bigr)$
  & $\mathrm{LogN}\!\bigl(\ln n {-} \mu,\,\sigma^2\bigr)$
  & Yes
  & No \\[4pt]
Weibull
  & $N_f \sim \mathrm{Wei}(k,\,\lambda)$
  & Inv.\ Weibull (Fr\'echet)
  & Yes
  & No \\[4pt]
Gumbel
  & $\ln N_f \sim \mathrm{Gumbel}(\mu,\,\beta)$
  & Log-Gumbel
  & Yes
  & No \\
\bottomrule
\end{tabular}
\end{table}

\paragraph{Log-normal.}
The most widely used assumption~\cite{MunizCalvente2022}: at a fixed stress level~$\sigma_a$, the fatigue life satisfies $\ln N_f \sim \mathcal{N}(\mu(\sigma_a),\,\sigma^2)$, where $\mu(\sigma_a) = \ln C - m\ln\sigma_a$ from Basquin's law and $\sigma$ captures the scatter.
Since $\ln(n/N_f) = \ln n - \ln N_f$, each increment is log-normally distributed: $\Delta D_{\mathrm{PM}} \sim \mathrm{LogN}(\ln n - \mu,\,\sigma^2)$, with mean $n\,e^{-\mu + \sigma^2/2}$ and variance $n^2 e^{-2\mu+\sigma^2}(e^{\sigma^2}-1)$.
The log-normal family is \emph{not} closed under summation: the sum of independent log-normals has no standard parametric form.
In practice, $D_{\mathrm{PM},k} = \sum_{t=1}^{k}\Delta D_{\mathrm{PM},t}$ is often approximated as log-normal via moment-matching (Fenton--Wilkinson method), but this is an approximation whose quality degrades as the number of summands grows or the coefficient of variation increases.

\paragraph{Weibull.}
Widely adopted in fatigue standards (e.g.\ Eurocode, ASTM): $N_f \sim \mathrm{Wei}(k,\,\lambda(\sigma_a))$ with shape~$k$ and scale~$\lambda$ depending on the stress level.
The reciprocal $1/N_f$ follows an inverse Weibull (Fr\'echet) distribution, so $\Delta D_{\mathrm{PM}} = n/N_f$ is a scaled Fr\'echet variable.
The sum of independent Fr\'echet variables has no closed form, making the marginal of~$D_{\mathrm{PM},k}$ intractable analytically.

\paragraph{Gumbel (extreme-value type~I).}
Castillo and Fern\'andez-Canteli~\cite{CastilloFernandezCanteli2009} proposed a compatible statistical model for the entire S--N field based on a Gumbel (minimum) distribution for~$\ln N_f$: $\ln N_f \sim \mathrm{Gumbel}(\mu(\sigma_a),\,\beta)$.
The increment $\Delta D_{\mathrm{PM}} = n/N_f$ then follows a log-Gumbel distribution.
As with the other two families, the log-Gumbel is not closed under summation.

\paragraph{Contrast with damage-accumulation models.}
A fundamental difference with the processes of Table~\ref{tab:deg_models} is that \emph{none} of the standard distributional assumptions for S--N scatter yields increments that are closed under summation.
The marginal distribution of $D_{\mathrm{PM},k}$ after $k$~loading blocks must therefore always be computed via convolution (analytically intractable in general), Monte Carlo simulation, or moment-based approximations.
In the fleet management formulation, this means that the reliability constraint $\Pr(D_{\mathrm{PM},k} \ge 1) \le \varepsilon$ cannot be evaluated in closed form from the cumulative parameters alone, in contrast to the Gamma or Gaussian models where the CDF of the marginal is directly available.
The upper bounds of Section~\ref{sec:stoch_orders} (Markov, Cantelli, Chernoff) remain applicable: because the mean and variance of~$D_{\mathrm{PM},k}$ are additive over increments regardless of the distributional family, the Cantelli bound provides a conservative but tractable constraint.
Alternatively, if the scatter is small enough that a deterministic S--N curve (e.g.\ the mean or a design percentile) is an acceptable approximation, the problem reduces to the deterministic Miner's rule and all constraints become linear.

\subsubsection{Conservative surrogates for the empirical damage-increment distribution}
\label{sec:cons_surrogates}

In the remaining-life setting the per-mission increment is the Palmgren--Miner contribution $\Delta D^{(j)} = \sum_i n_i/N_{f,i}$ of Equation~\eqref{eq:miner}, and in our intended use its distribution is not postulated but \emph{measured}: a rainflow count~\cite{MunizCalvente2022} of load histories collected for mission~$j$ yields an empirical probability density of the non-negative increment $\Delta D^{(j)}$.
We deliberately do \emph{not} fit a parametric law to this density---a fit ties the model to one data set and one functional form, undermining the generality the abstract formulation seeks---and instead summarise it by a few distributional functionals (moments, support, selected quantiles) and replace it by a \emph{conservative} surrogate.
Conservative here means that the surrogate never \emph{under}-states damage, so that the reliability constraint $\Pr(D_{i\ell k} > \tau) \le \varepsilon$ and the loop constraint~\eqref{eq:loop} remain valid by construction.
Two features of the problem delimit the admissible options.
First, both constraints act on the \emph{accumulated} damage $D = \sum_t \Delta D_t$ of Equation~\eqref{eq:discrete_damage}, not on a single increment, so the surrogate must be evaluated after summation.
Second, as observed in the contrast above, none of the standard S--N scatter families is closed under summation, so the law of $D$ is generally unavailable in closed form and must be either bounded or replaced by a convolution-closed surrogate.

Crucially, Route~I never requires the moments or support of the accumulated damage~$D$ in advance: it consumes only \emph{per-mission} information, which the schedule then assembles through the assignment counts.
Fix a component $(i,\ell)$ and an accumulation window, write $D = D_{i\ell}$ for its damage over that window, and let $n_{ij} = \sum_t x_{ijt}$ be the number of days mission~$j$ is served.
Modelling repeated executions of mission~$j$ on component~$\ell$ as independent draws from a mission-specific increment law---with mean $\mu_{\ell j}$, variance $v_{\ell j}$, support $[0,b_{\ell j}]$, and cumulant generating function $\psi_{\ell j}(s) = \ln \mathbb{E}\big[e^{s\Delta D^{(j)}}\big]$, where $\mathbb{E}[e^{sX}]$ is the moment-generating function---independence across days gives
\begin{equation}\label{eq:additive_functionals}
  \mathbb{E}[D] = \sum_j n_{ij}\,\mu_{\ell j}, \qquad
  \operatorname{Var}[D] = \sum_j n_{ij}\,v_{\ell j}, \qquad
  \ln M_D(s) = \sum_j n_{ij}\,\psi_{\ell j}(s).
\end{equation}
Each functional of~$D$ is therefore a count-weighted sum of per-mission quantities---linear in the schedule, with the support entering only through the per-mission $b_{\ell j}$, never through a shared worst-case value.
Repetition \emph{merges} into the tail description additively: serving mission~$j$ a total of $n_{ij}$ times multiplies its cumulant generating function by $n_{ij}$, and distinct missions combine by summation, so a gentle mission keeps its small $\psi_{\ell j}$ beside a punishing one rather than being averaged with it.
These per-mission functionals support two routes.
\emph{Route~I} bounds the upper tail of $D$ directly from the functionals in~\eqref{eq:additive_functionals}.
\emph{Route~II} replaces each increment by a stochastically larger surrogate drawn from a family closed under convolution and propagates the order to the sum, using the following standard closure property of the usual stochastic order $\le_{\mathrm{st}}$ (Definition~\ref{def:st}).

\begin{proposition}[Closure of $\le_{\mathrm{st}}$ under independent convolution]\label{prop:st_closure}
Let $\Delta D_t \le_{\mathrm{st}} \widetilde{\Delta D}_t$ for $t = 1,\dots,k$, with $\{\Delta D_t\}_t$ mutually independent and $\{\widetilde{\Delta D}_t\}_t$ mutually independent.
Then $\sum_t \Delta D_t \le_{\mathrm{st}} \sum_t \widetilde{\Delta D}_t$.
More generally, the conclusion holds with the sum replaced by any coordinatewise non-decreasing map $g$ (see Shaked and Shanthikumar~\cite[\S1.A]{ShakedShanthikumar2007}).
\end{proposition}
\begin{proof}
By the coupling characterisation of $\le_{\mathrm{st}}$ (Definition~\ref{def:st}), for each $t$ there exist $\hat X_t \stackrel{d}{=} \Delta D_t$ and $\hat Y_t \stackrel{d}{=} \widetilde{\Delta D}_t$ on a common probability space with $\hat X_t \le \hat Y_t$ almost surely.
Realise these couplings independently across $t$.
Then $\sum_t \hat X_t \le \sum_t \hat Y_t$ almost surely, while independence gives $\sum_t \hat X_t \stackrel{d}{=} \sum_t \Delta D_t$ and $\sum_t \hat Y_t \stackrel{d}{=} \sum_t \widetilde{\Delta D}_t$; hence $\sum_t \Delta D_t \le_{\mathrm{st}} \sum_t \widetilde{\Delta D}_t$.
Any non-decreasing $g$ preserves the almost-sure inequality, yielding the general statement.
\end{proof}

Proposition~\ref{prop:st_closure} is what makes Route~II usable beyond a single sum: the dynamics~\eqref{eq:dyn_full} is a coordinatewise non-decreasing map of the increments and the (independent) reset variables, so a per-increment stochastic upper bound propagates to a bound on the entire state $D_{i\ell k}$---repairs and replacements included---and hence to $\Pr(D_{i\ell k} > \tau)$ at every step.
Table~\ref{tab:cons_surrogates} collects the resulting options together with their guarantees and trade-offs.

\begin{table}[htbp]
\centering
\caption{Ways to summarise the empirical per-mission damage-increment distribution and propagate it to the accumulated damage $D = \sum_t \Delta D_t$ for the constraint $\Pr(D > \tau) \le \varepsilon$.
``Conservative'' means a guaranteed upper bound on the failure probability (Route~I) or $\le_{\mathrm{st}}$ dominance (Route~II); ``$\widetilde{D}$'' means the method yields an explicit surrogate \emph{law} for the sum, as opposed to a one-sided bound.
All inputs are \emph{per mission}; the schedule assembles them through the assignment counts $n_{ij}$ via Equation~\eqref{eq:additive_functionals}, so the support enters only per mission and is never shared.}
\label{tab:cons_surrogates}
\footnotesize
\setlength{\tabcolsep}{4.5pt}
\begin{tabular}{@{}p{1.7cm}p{1.9cm}p{0.9cm}p{0.7cm}p{2.9cm}p{2.9cm}p{1.5cm}@{}}
\toprule
\textbf{Method} & \textbf{Inputs} & \textbf{Cons.?} & \textbf{$\widetilde{D}$?} & \textbf{Pros} & \textbf{Cons} & \textbf{Ref.} \\
\midrule
\multicolumn{7}{@{}l}{\textit{Route I --- bound the tail of $D$ directly (no surrogate law)}} \\
Markov
  & mean
  & Yes & No
  & needs only the mean; always valid
  & extremely loose
  & \cite{BertsimasPopescu2005} \\[3pt]
Cantelli
  & mean, var
  & Yes & No
  & only two moments; tight \emph{as} a two-moment bound
  & ignores independence; loose in the tail
  & \cite{BertsimasPopescu2005} \\[3pt]
Bernstein
  & per-mission mean, var, support $b_j$
  & Yes & No
  & exploits independence; few parameters; moment inputs linear in the counts
  & needs per-mission support; one-sided only
  & \cite{BoucheronLugosiMassart2013} \\[3pt]
Chernoff
  & per-mission c.g.f.\ $\psi_j(s)$ (or mean $+$ support)
  & Yes & No
  & tightest in the tail; count-weighted; valid for any fixed $s>0$
  & needs the c.g.f.; empirical c.g.f.\ is an estimate, not a bound
  & \cite{Chernoff1952} \\[3pt]
\multicolumn{7}{@{}l}{\textit{Route II --- dominate each increment, then convolve (explicit $\widetilde{D}$)}} \\
Moment-matched Gamma / IG
  & mean, var
  & No & Yes
  & explicit law; closed under $+$; ordering tables apply
  & matches moments but does not dominate (survival curves cross)
  & \cite{vanNoortwijk2009,YeChen2014} \\[3pt]
Dominating shared-rate Gamma
  & moments $+$ empirical survival
  & Yes & Yes
  & conservative \emph{and} explicit; serves reliability and loop; shape linear in the schedule
  & shared rate $+$ dominance margin add slack; commits to Gamma
  & \cite{ShakedShanthikumar2007} \\
\bottomrule
\end{tabular}
\end{table}

The Route~I entries are elementary, which is why they require no distributional assumption.
With $D \ge 0$ and $\tau > 0$, Markov's inequality $\Pr(D > \tau) \le \mathbb{E}[D]/\tau$ follows from $\tau\,\mathbf{1}\{D > \tau\} \le D$.
Applying Markov to $(D - \mathbb{E}[D] + c)^2$ and optimising over $c > 0$ gives the one-sided Chebyshev (Cantelli) bound, for $\tau > \mathbb{E}[D]$,
\begin{equation}\label{eq:cantelli_surrogate}
  \Pr(D > \tau) \;\le\; \frac{\operatorname{Var}[D]}{\operatorname{Var}[D] + (\tau - \mathbb{E}[D])^2}.
\end{equation}
Both are the \emph{tight} bounds attainable from the first one and two moments, respectively~\cite{BertsimasPopescu2005}, and both are attained (for the closed event $D \ge \tau$) by two-point distributions---which is exactly why they ignore, and are loose under, the independent-sum structure.
The exponential (Cram\'er--Chernoff) bound,
\begin{equation}\label{eq:chernoff_surrogate}
  \Pr(D > \tau) \;\le\; \inf_{s > 0} e^{-s\tau} M_D(s)
  \;=\; \inf_{s > 0} \exp\!\Big( -s\tau + \textstyle\sum_t \ln M_{\Delta D_t}(s) \Big),
\end{equation}
follows from Markov applied to $e^{sD}$ and is valid for every fixed $s > 0$~\cite{Chernoff1952}; Bernstein's inequality is the specialisation for independent summands bounded in $[0,b]$, $\Pr(D - \mathbb{E}[D] > a) \le \exp\big( -\tfrac{a^2/2}{\operatorname{Var}[D] + b\,a/3} \big)$ with $a = \tau - \mathbb{E}[D]$~\cite{BoucheronLugosiMassart2013}.
Unlike Markov and Cantelli, these use the additive log-MGF of~\eqref{eq:additive_functionals} and therefore exploit independence.

\paragraph{Per-mission inputs do not flatten the missions.}
Collapsing the summands to a single worst-case support $b = \max_j b_{\ell j}$ would make Bernstein treat every mission as the most damaging one, erasing the differences the schedule exploits; this flattening is an artefact of the closed form, not of Route~I.
Kept per term, each bound retains every mission's own support and tail.
Hoeffding reads $\Pr(D - \mathbb{E}[D] > a) \le \exp\!\big( -2a^2 / \sum_j n_{ij} b_{\ell j}^2 \big)$ with $a = \tau - \mathbb{E}[D]$, where mission~$j$ enters only through its own $b_{\ell j}^2$, and the Chernoff bound~\eqref{eq:chernoff_surrogate} becomes the count-weighted reliability constraint
\begin{equation}\label{eq:chernoff_permission}
  \Pr(D > \tau) \;\le\; \inf_{s > 0} \exp\!\Big( -s\tau + \sum_j n_{ij}\,\psi_{\ell j}(s) \Big),
\end{equation}
in which each mission contributes its own cumulant generating function $\psi_{\ell j}$.
For a fixed~$s$ the exponent is linear in the assignment counts, so~\eqref{eq:chernoff_permission} is a linear reliability constraint; the inner infimum couples~$s$ to the counts---for a fixed schedule it is a one-dimensional convex problem whose minimiser $s^\star$ solves $\sum_j n_{ij}\,\psi_{\ell j}'(s) = \tau$ (the exponentially tilted mean of~$D$ equals~$\tau$), but since $s^\star$ depends on the counts one either fixes a conservative~$s$ (any $s > 0$ remains a valid bound) or sweeps a small grid.

Bernstein's inequality is the one Route~I bound that resists this per-mission treatment, and the reason is the same coupling seen from the other side: its closed form requires a \emph{single} support constant~$b_\ell$ in the correction term (only the variance stays per-mission).
Between replacements the damage is a weighted sum, with weights in $[0,1]$, of the terms of the initial or replacement damage and of the increments, and Bernstein's inequality needs a bound on every centred term.
So the constant is $b_\ell = \max\{b^{\mathrm{start}}_\ell,\, b^{\mathrm{new}}_\ell,\, \max_j b_{\ell j}\}$, where $b^{\mathrm{start}}_\ell$ and $b^{\mathrm{new}}_\ell$ bound the terms of the initial and of the replacement damage, and are $0$ for a known damage.
Explicitly, with $\kappa = \ln(1/\varepsilon)$, the reliability constraint $\Pr(D > \tau) \le \varepsilon$ becomes
\begin{equation}\label{eq:bernstein_constraint}
  \sum_j n_{ij}\,\mu_{\ell j}
  \;+\; \frac{\kappa\, b_\ell}{3}
  \;+\; \sqrt{ \Big( \tfrac{\kappa\, b_\ell}{3} \Big)^2 + 2\kappa \sum_j n_{ij}\, v_{\ell j} }
  \;\le\; \tau,
\end{equation}
a constraint whose left-hand side is \emph{concave} in the assignment counts (equivalently $(\kappa b_\ell/3)^2 + 2\kappa\sum_j n_{ij} v_{\ell j} \le (\tau - \kappa b_\ell/3 - \sum_j n_{ij}\mu_{\ell j})^2$ together with the side row $\sum_j n_{ij}\mu_{\ell j} \le \tau - \kappa b_\ell/3$, which keeps the bracket non-negative): by~\eqref{eq:additive_functionals} the mean and the variance are \emph{affine} in the counts, so the square root is concave and the feasible set is reverse-convex rather than second-order-cone-representable---in the classical SOCP chance constraint the variance is a convex quadratic of the decisions, whereas here independence across days makes it affine, which flips the curvature.
For the MILP this is a feature rather than an obstacle: every tangent line to the concave square root \emph{over}-estimates it, so replacing the square root in~\eqref{eq:bernstein_constraint} by a fixed tangent---or by the piecewise-linear pointwise minimum of several tangents, encoded exactly as a PWL equality (SOS2)---yields linear constraints in the counts that imply~\eqref{eq:bernstein_constraint} and therefore remain conservative.
Retaining a per-mission support $b_{\ell j}$ in the correction is possible only by keeping the per-term cumulant sum $\sum_j n_{ij}\,(s^2 v_{\ell j}/2)/(1 - s\,b_{\ell j}/3)$ and minimising over~$s$---which is exactly the Chernoff/Bennett bound~\eqref{eq:chernoff_permission}.
The $s$-free closed form~\eqref{eq:bernstein_constraint} and per-mission support resolution therefore cannot be obtained simultaneously.

\paragraph{Guaranteed bound versus estimate.}
For~\eqref{eq:chernoff_permission} to be a true bound, the per-mission ingredient must be an \emph{upper} bound on $\psi_{\ell j}$.
Two per-mission choices suffice and keep the support local: the mean-and-support bound $\psi_{\ell j}(s) \le \ln\!\big( 1 + (\mu_{\ell j}/b_{\ell j})(e^{s b_{\ell j}} - 1) \big)$, attained by the two-point law on $\{0, b_{\ell j}\}$; or the empirical cumulant generating function $\widehat\psi_{\ell j}(s) = \ln\!\big( \tfrac{1}{m} \sum_{r=1}^{m} e^{s \Delta D^{(j)}_r} \big)$ built from the $m$ recorded executions of mission~$j$, which is an \emph{estimate}, not a bound, and so must carry a one-sided confidence margin.
Repeating a mission tightens both ingredients---more executions sharpen $\widehat\psi_{\ell j}$, $\mu_{\ell j}$ and $b_{\ell j}$ and shrink the margin---which is the precise sense in which undergoing a mission many times refines its tail description.

The choice among the Route~I bounds depends on how deep into the tail the constraint sits.
In the regime relevant to a reliability constraint with small~$\varepsilon$---several standard deviations above $\mathbb{E}[D]$---the exponential bounds are markedly tighter: in a representative numerical experiment (non-negative increments bounded in $[0,1]$, $k \in \{30,100\}$ summands), at $\varepsilon \sim 10^{-4}$ the Chernoff bound stayed within roughly one order of magnitude of the empirical tail, whereas the Cantelli bound over-stated it by more than two orders; Bernstein fell between them and overtook Cantelli only as the tail deepened and the number of summands grew.\footnote{Recorded as a numerical illustration, not a proof; the test code and output are kept in the verification workspace. Asymptotically the gap is structural: Cantelli decays only polynomially, as $\operatorname{Var}[D]/(\tau-\mathbb{E}[D])^2$, whereas~\eqref{eq:chernoff_surrogate} decays exponentially in the deviation.}
Cantelli, though the simplest, is therefore best reserved for short accumulation windows or as a coarse always-valid fallback.

Route~II yields not a bound but an explicit surrogate law.
Matching only the first two moments of a convolution-closed family (Gamma, with shape $\alpha = \mathbb{E}[\Delta D]^2/\operatorname{Var}[\Delta D]$ and rate $\beta = \mathbb{E}[\Delta D]/\operatorname{Var}[\Delta D]$; or inverse Gaussian) is convenient---the sum stays in the family under the closure conditions of Table~\ref{tab:discrete_models}---but is \emph{not} conservative: a moment-matched law generally crosses the empirical survival function rather than dominating it.
Conservatism is recovered by using more of the empirical information than its moments.
Fix a common rate $\beta^\star$ and, per mission, inflate the shape $\alpha_j$ until the surrogate survival dominates the empirical one, $\overline{\mathrm{Ga}}(x;\alpha_j,\beta^\star) \ge \widehat{\bar F}_j(x)$ for all $x$---always achievable, since $\overline{\mathrm{Ga}}(x;\alpha,\beta^\star) \to 1$ as $\alpha \to \infty$.
By Proposition~\ref{prop:st_closure} the accumulated surrogate is then $\widetilde{D} = \mathrm{Ga}(A,\beta^\star)$ with cumulative shape $A = \sum_{t,j} x_{ijt}\,\alpha_j$ linear in the schedule, the reliability constraint reduces to the scalar condition $A \le A^\star(\varepsilon)$, and the loop constraint, in the absence of repairs, to the row $A_T \le A_{H_1}$ on the cumulative shape.
The price is the slack introduced by the shared rate and the dominance margin.

In summary, the formulation certifies the reliability constraint of a component with measured increments by one of five Route-I bounds, chosen per component: Markov, Cantelli, Hoeffding, Bernstein or Chernoff.
The loop constraint acts on the statistics that the chosen bound tracks (Section~\ref{sec:loop_routeI}), so it needs no surrogate law, and Route~II is not part of the formulation.
The irreducible trade-off is that the bounds assume nothing about the shape but yield only a tail bound, whereas the dominating Gamma yields a propagatable law at the cost of a family commitment and some conservatism.

\paragraph{The five bounds of the formulation.}
The rainflow model of the formulation is \emph{distribution-free}.
The law of a mission increment is measured from recorded load histories by rainflow counting (Palmgren--Miner damage), and it enters the model only through summary statistics: its mean and, as the bound requires, its variance, the width of its range, or its cumulant generating function at a fixed point~$s$.
No parametric family is assumed.
Each bound tracks the mean $\mu_k$ of the damage on day~$k$ and at most one more statistic.
The second statistic is the variance~$v_k$, the \emph{Hoeffding budget}~$R_k$ (the sum of the squared range widths of the terms of the damage), or the cumulant generating function $K_k = \ln\mathbb{E}\bigl[e^{sD_k}\bigr]$ at a fixed point $s > 0$.
Table~\ref{tab:five_bounds} lists the rows, with $\kappa = \ln(1/\varepsilon)$ and $c = \kappa\, b_\ell/3$.
Each row is the corresponding tail bound of this section, solved for the target~$\varepsilon$ and written in the tracked statistics.
If the tracked values satisfy a row and are at least the true statistics of the damage, the bound gives $\Pr(D_k > \tau) \le \varepsilon$.
The rows rest on four modelling assumptions.
\begin{itemize}[nosep]
  \item \textbf{A0 (non-negative damage).} Every increment, the initial damage and every replacement draw are non-negative.
  \item \textbf{A1 (independence).} The mission increments are mutually independent, and independent of the initial damage and of every replacement draw.
    Each replacement installs a new draw, independent of everything else.
    Every bound except Markov needs A1.
  \item \textbf{A2 (bounded ranges; Hoeffding and Bernstein).} Each mission increment lies in an interval of known width.
    The initial damage is a known part plus independent terms, each in an interval; its Hoeffding budget is the sum of their squared widths, and its support is the largest width.
    The replacement draw is described in the same way.
    A known damage has budget and support~$0$.
  \item \textbf{A3 (repetition).} The row of the bound holds on every day $k = 1, \ldots, T$, and after day~$T$ the operating phase repeats with period~$H_2$, schedule and increment laws included.
    This is the setting of Theorem~\ref{thm:perpetual}.
\end{itemize}

\begin{table}[htbp]
\centering
\caption{Reliability rows of the five Route-I bounds, on every day~$k$ of the horizon.
$\mu_k$: tracked mean; $v_k$: tracked variance; $R_k$: tracked Hoeffding budget; $K_k$: tracked cumulant generating function at the fixed point~$s$.}
\label{tab:five_bounds}
\small
\begin{tabular}{@{}lll@{}}
\toprule
\textbf{Bound} & \textbf{Tracked statistics} & \textbf{Reliability row} \\
\midrule
Markov    & $\mu$       & $\mu_k \le \varepsilon\,\tau$ \\
Cantelli  & $\mu$, $v$  & $(1-\varepsilon)\,v_k \le \varepsilon\,(\tau - \mu_k)^2$, \quad $\mu_k \le \tau$ \\
Hoeffding & $\mu$, $R$  & $(\tau - \mu_k)^2 \ge (\kappa/2)\,R_k$, \quad $\mu_k \le \tau$ \\
Bernstein & $\mu$, $v$  & $(\tau - c - \mu_k)^2 \ge c^2 + 2\kappa\, v_k$, \quad $\mu_k \le \tau - c$ \\
Chernoff  & $\mu$, $K$  & $K_k - s\,\tau \le \ln\varepsilon$ \\
\bottomrule
\end{tabular}
\end{table}

\paragraph{Width of a range.}
The support of an increment~$X$ is the width~$b$ of an interval $[l, l + b]$ that contains $X$ almost surely; the width keeps this one meaning in every bound, and a width of~$0$ means that the increment is known.
Hoeffding's inequality uses the widths.
Bernstein's inequality needs an upper bound of the centred term, and the width is one: $X - \mathbb{E}X \le l + b - l = b$, since $\mathbb{E}X \ge l$.
A width~$w$ is consistent with a variance~$v$ only if $v \le w^2/4$ (Popoviciu's inequality).
Proof: let $X$ take values in an interval of width~$w$ with midpoint~$m_c$; then $\operatorname{Var} X \le \mathbb{E}\,(X - m_c)^2 \le (w/2)^2$.
The width is the only information on the range, and the range may start at any point $l \ge 0$.
The inequality $\operatorname{Var} X \le w^2/4$ holds for every such position, so it is a necessary condition for the pair (variance, width).
For example, increments in an interval of width $0.2$ have variance at most $0.01$; a two-point law on $\{0, 0.2\}$ with probability $1/2$ each has mean $0.1$ and reaches $0.01$.
A condition of the form $v \le \mu\,(b - \mu)$ assumes that the range starts at~$0$, so it is not necessary for a width.

\paragraph{Day maps of the statistics.}
Between replacements the damage is $D_k = \sum_t w_t X_t$: the $X_t$ are the known part, the independent terms of the initial damage or of the replacement draw, and the mission increments.
Each weight $w_t \in [0, 1]$ is a product of factors $1 - \rho$, and a term of weight~$0$, possible when $\rho = 1$, is dropped.
Under $\text{ARD}_1$, write $D^- = G + \Delta$, with $G$ the damage just after the last intervention and $\Delta$ the weighted increments added since.
$G$ and $\Delta$ are independent by A1, and a repair gives $D^+ = G + (1-\rho)\,\Delta$.
Before the first intervention, $G$ is the part of the initial damage that the starting anchors describe (Section~\ref{sec:repair_connection}).
For a damage $D = \sum_t w_t X_t$, write $m(D) = \mathbb{E}D$, $V(D) = \operatorname{Var} D$, $R(D) = \sum_t w_t^2 r_t$ with $r_t$ the squared width of~$X_t$, and $K(D) = \ln\mathbb{E}\,e^{sD}$.

\begin{proposition}[Day maps of the statistics]\label{prop:rf_maps}
Let $\theta \in \{m, V, R\}$, with factor $a_\theta = 1-\rho$ for $m$ and $a_\theta = (1-\rho)^2$ for $V$ and~$R$.
Under A1:
\begin{enumerate}[nosep]
  \item a mission day adds the statistic of the increment: $\theta(D + X) = \theta(D) + \theta(X)$ and $K(D + X) = K(D) + K(X)$;
  \item under $\text{ARD}_\infty$: $\theta\bigl((1-\rho)D\bigr) = a_\theta\,\theta(D)$ and $K\bigl((1-\rho)D\bigr) \le (1-\rho)\,K(D)$;
  \item under $\text{ARD}_1$: $\theta(D^+) = a_\theta\,\theta(D^-) + (1 - a_\theta)\,\theta(G)$ and $K(D^+) \le (1-\rho)\,K(D^-) + \rho\,K(G)$.
\end{enumerate}
\end{proposition}
\begin{proof}
Item~1: the mean and $R$ are additive by definition, the variance and the cumulant generating function by independence.
Item~2: $m$ and $V$ scale by $1-\rho$ and $(1-\rho)^2$; every weight is multiplied by $1-\rho$, so $R$ scales by $(1-\rho)^2$.
A cumulant generating function is convex with value~$0$ at~$0$, so $K\bigl((1-\rho)D\bigr) = K_D\bigl((1-\rho)s\bigr) \le (1-\rho)\,K_D(s) + \rho\,K_D(0) = (1-\rho)\,K(D)$, with $K_D$ the cumulant generating function of~$D$ as a function of its argument.
Item~3: by independence, $\theta(D^-) = \theta(G) + \theta(\Delta)$; for $R$ the weights of the terms of~$G$ are unchanged and those of~$\Delta$ are multiplied by $1-\rho$.
Hence
\[
  \theta(D^+) = \theta(G) + a_\theta\,\theta(\Delta)
  = \theta(G) + a_\theta\bigl(\theta(D^-) - \theta(G)\bigr)
  = a_\theta\,\theta(D^-) + (1 - a_\theta)\,\theta(G),
\]
and, with item~2 applied to~$\Delta$,
\[
  K(D^+) = K(G) + K_\Delta\bigl((1-\rho)s\bigr)
  \le K(G) + (1-\rho)\bigl(K(D^-) - K(G)\bigr)
  = (1-\rho)\,K(D^-) + \rho\,K(G). \qedhere
\]
\end{proof}

So each tracked statistic follows the repair map $a\,\theta + (1-a)\,g$ of Section~\ref{sec:repair_connection}, with $a = 1-\rho$ for $\mu$ and $K$, and $a = (1-\rho)^2$ for $v$ and~$R$; its anchor is its value just after the last intervention.
The maps of $\mu$, $v$ and $R$ are exact; the map of $K$ is an upper bound, which is what the domination argument of Section~\ref{sec:encoding} needs.
Hence each true statistic of~$D_k$, written $m$, $V$, $R^{\mathrm{t}}$ and $K^{\mathrm{t}}$, is at most the tracked statistic of the same name (Proposition~\ref{prop:domination}).

\paragraph{Derivation and validity of the rows.}
Write $d_k = \tau - \mu_k$.
\begin{itemize}
  \item \emph{Markov.} By A0 and Markov's inequality, $\Pr(D_k > \tau) \le m/\tau \le \mu_k/\tau$, so the row $\mu_k \le \varepsilon\,\tau$ suffices.
    It uses the mean only and does not need A1 or A2.
  \item \emph{Cantelli.} Cantelli's inequality gives $\Pr(D - m \ge a) \le V/(V + a^2)$ for $a > 0$; the row is equivalent to $\mu_k + \kappa_C\sqrt{v_k} \le \tau$ with $\kappa_C = \sqrt{(1-\varepsilon)/\varepsilon}$.
    If $d_k > 0$,
    \[
      \Pr(D_k > \tau) \le \Pr(D_k - m \ge \tau - m) \le \frac{V}{V + (\tau - m)^2} \le \frac{v_k}{v_k + d_k^2} \le \varepsilon,
    \]
    because $v/(v + d^2)$ increases in~$v$ and decreases in~$d$, and $\tau - m \ge d_k$.
    If $d_k = 0$, the row gives $v_k = 0$, so $V = 0$ and $D_k = m \le \tau$.
  \item \emph{Hoeffding.} $D_k - m$ is a sum of independent centred terms $w_t(X_t - \mathbb{E}X_t)$ with ranges of widths $w_t\sqrt{r_t}$, and the sum of their squares is~$R^{\mathrm{t}}$.
    Hoeffding's inequality gives $\Pr(D_k - m \ge a) \le \exp(-2a^2/R^{\mathrm{t}})$, and $\exp(-2d^2/R) \le \varepsilon$ exactly when $d^2 \ge (\kappa/2)\,R$.
    If $R^{\mathrm{t}} = 0$, every term is constant and $D_k = m \le \tau$; this covers $d_k = 0$, where the row gives $R_k = 0$.
    Otherwise $d_k > 0$, and $\Pr(D_k > \tau) \le \exp\bigl(-2(\tau - m)^2/R^{\mathrm{t}}\bigr) \le \exp(-2d_k^2/R_k) \le \varepsilon$, because $\exp(-2d^2/R)$ increases in~$R$ and decreases in~$d$.
  \item \emph{Bernstein.} The support constant $b_\ell$ of~\eqref{eq:bernstein_constraint} is a constant, not a decision variable; the replacement support enters it only when replacements are allowed.
    Every weighted term satisfies $w_t(X_t - \mathbb{E}X_t) \le w_t\, b_t \le b_\ell$, with $b_t$ the width of the range of~$X_t$, whatever the schedule, because $w_t \le 1$.
    Bernstein's inequality for independent terms bounded above by~$b_\ell$ gives $\Pr(D - m \ge a) \le \exp\bigl(-a^2/(2(V + b_\ell\,a/3))\bigr)$ for $a > 0$; the right side is at most~$\varepsilon$ exactly when $a^2 - 2c\,a - 2\kappa V \ge 0$.
    With $v_k \ge 0$, the row and its side row give $d_k \ge 2c$.
    The function $a^2/(2(V + b_\ell\,a/3))$ increases in $a > 0$ and decreases in~$V$, and $\tau - m \ge d_k \ge 2c$; so $\Pr(D_k > \tau) \le \varepsilon$ when $b_\ell > 0$.
    When $b_\ell = 0$, every term is constant and $D_k = m \le \tau$.
  \item \emph{Chernoff.} For $s > 0$, Markov's inequality applied to $e^{sD}$ gives $\Pr(D_k \ge \tau) \le \exp(K^{\mathrm{t}} - s\tau)$.
    Since $K^{\mathrm{t}} \le K_k$, the row gives $\Pr(D_k > \tau) \le \exp(K_k - s\tau) \le \varepsilon$.
    The row bounds $K$ by $s\tau - \kappa$, its natural range, and $K \ge 0$, because $K(X) \ge s\,\mathbb{E}X \ge 0$ by Jensen's inequality and A0.
    So a known zero damage satisfies the row only if $s\tau \ge \kappa$, and a replacement day satisfies it only if the replacement value of~$K$ is at most $s\tau - \kappa$.
\end{itemize}

\paragraph{The Bernstein side row.}
The quadratic Bernstein row alone also admits the point $\mu_k = \tau$, $v_k = 0$: a known damage equal to~$\tau$.
That point is safe, but it is cut off from the rest of the feasible set.
The side row $\mu_k \le \tau - c$ removes it and keeps the branch $\tau - \mu_k \ge c + \sqrt{c^2 + 2\kappa v_k}$ of~\eqref{eq:bernstein_constraint}.
At $\mu_k = v_k = 0$ the row needs $\tau \ge 2c$, and a positive variance is feasible only if $\tau > 2c = 2\kappa b_\ell/3$.

\paragraph{Starting and replacement values of the Chernoff statistic.}
The statistic $K_k$ is the cumulant generating function of the whole damage, initial damage included.
By Jensen's inequality, $\ln\mathbb{E}[e^{sD}] \ge s\,\mathbb{E}[D]$, with equality exactly when $D$ is almost surely constant.
So a known initial damage $\mu_0$ gives the starting value $K_0 = s\,\mu_0$, and a known replacement damage $\mu^{\mathrm{new}}$ gives the reset value $s\,\mu^{\mathrm{new}}$.
A random initial or replacement damage needs its own cumulant generating function at~$s$, which is larger.
For the same reason the cumulant generating function of every increment at~$s$ is at least $s$ times its mean, and the starting value and the replacement value of~$K$ are at least $s$ times the corresponding mean; an input check enforces these inequalities (Section~\ref{sec:check_reasons}).
A start at $K_0 = 0$ ignores the initial damage and can accept an unsafe schedule.
Example: $\tau = 1$, $\varepsilon = 0.1$, $s = 5$, a known initial damage $\mu_0 = 0.6$, and one mission that adds $0$ or $0.5$ with probability $1/2$ each, whose cumulant generating function at~$s$ is $\ln(0.5 + 0.5\,e^{2.5}) = 1.886$.
The true failure probability is $0.5$.
With $K_0 = 0$ the row reads $1.886 - 5 = -3.11 \le \ln 0.1 = -2.30$ and accepts the schedule.
With $K_0 = s\,\mu_0 = 3$ it reads $4.886 - 5 = -0.11 > -2.30$ and rejects it.
In the same way, the Hoeffding budget of a known damage is $0$, because its range has width~$0$, and a random initial or replacement damage needs its own budget.
So the defaults of the formulation describe a known initial damage and a known replacement damage: zero variance, zero support, zero budget, and $K = s$ times the mean.

\subsubsection{The loop constraint under the Route-I bounds}
\label{sec:loop_routeI}

The probability comparison~\eqref{eq:loop_prob} asks for $\Pr(D_T > \tau) \le \Pr(D_{H_1} > \tau)$ so that concatenating the schedule keeps the failure probability below~$\varepsilon$ indefinitely.
Under Route~I we never observe these probabilities; we have only an upper bound $U_k \ge \Pr(D_k > \tau)$ from one of the bounds of Table~\ref{tab:five_bounds}, for instance the count-weighted Chernoff bound~\eqref{eq:chernoff_permission} or the Bernstein constraint~\eqref{eq:bernstein_constraint}.
We write $U^{(n)}$ for the bound on day $H_1 + nH_2$, the end of the $n$-th repetition of the operating phase, so $U^{(0)} = U_{H_1}$ and $U^{(1)} = U_T$.
The tempting translation---replace both sides by their bounds and impose $U^{(1)} \le U^{(0)}$---is \emph{not} sufficient for repeatability, and it is worth seeing precisely why.

\paragraph{Why $U^{(1)} \le U^{(0)}$ does not chain.}
The bound $U$ is a many-to-one function of the state's parameters, and its level sets are not invariant under the dynamics.
Bernstein's bound, for instance, depends on $\big(\mathbb{E}[D],\operatorname{Var}[D]\big)$ through a single scalar, so $U^{(1)} = U^{(0)}$ can hold with $\big(\mathbb{E}[D_T],\operatorname{Var}[D_T]\big) \ne \big(\mathbb{E}[D_{H_1}],\operatorname{Var}[D_{H_1}]\big)$---a larger mean offset by a smaller variance, say.
One further period of the dynamics maps these two distinct parameter points to \emph{different} bound values, so $U^{(2)} > U^{(1)} = U^{(0)}$ is possible: the certified bound can creep upward over repetitions, and the single-period inequality does not induct.\footnote{A concrete instance: $\tau = 10$, $b_\ell = 3$, $\rho = 0.5$, repeated block $=$ one repair day followed by mission days with $\sum_j n_j\mu_{\ell j} = 3$ and $\sum_j n_j v_{\ell j} = 3.99$; from $\bigl(\mathbb{E}[D_{H_1}],\operatorname{Var}[D_{H_1}]\bigr) = (4,9)$ the Cantelli bound gives $U^{(0)} = 0.2000$ and $U^{(1)} = 0.1997 \le U^{(0)}$, yet $U^{(2)} = 0.2151 > U^{(0)}$, with $U^{(n)} \to 0.2495$. An analogous Bernstein instance (block variance $3.15$) violates at the second repetition as well. The loop constraint~\eqref{eq:loop} rejects both, since $\mathbb{E}[D_T] = 5 > 4$. Test scripts are kept in the verification workspace.}
Adding a fixed margin $\delta > 0$, $U^{(1)} \le U^{(0)} - \delta$, does not repair this, because the defect is the non-invariance of the level sets, not a constant gap.

\paragraph{The certificate stays valid at every period.}
This is not an undetectable worsening.
The statistics that Route~I tracks---the mean, the variance, the Hoeffding budget, or the cumulant generating function at a fixed point~$s$---follow monotone recursions under the dynamics~\eqref{eq:dyn_full}: increments add them, imperfect repair maps a statistic~$\theta$ to $a\,\theta + (1-a)\,g$ with its anchor~$g$ and a factor $a \in [0,1]$ (Section~\ref{sec:repair_connection}), and replacement resets them.
The recursions are exact for the mean, the variance and the Hoeffding budget, and an upper bound for the cumulant generating function (see below).
Consequently $U_k$ is a valid bound on \emph{every} day of the endless schedule, $\Pr(D_k > \tau) \le U_k$, and the true probability can exceed~$\varepsilon$ only if $U_k$ itself does.
Repeatability therefore needs a loop condition that forces $U_k \le \varepsilon$ on every such day, which means imposing it on a statistic that both determines $U$ and propagates monotonically---not on the lossy scalar~$U$.

\paragraph{Looping on the moments (the native condition).}
For the Cantelli and Bernstein bounds, the loop constraint~\eqref{eq:loop} acts on the tracked moments,
\begin{equation}\label{eq:loop_moments}
  \mathbb{E}[D_T] \;\le\; \mathbb{E}[D_{H_1}]
  \qquad\text{and}\qquad
  \operatorname{Var}[D_T] \;\le\; \operatorname{Var}[D_{H_1}],
\end{equation}
and, under $\text{ARD}_1$, on their anchors $g^{\mu}$ and $g^{v}$ in the same way.
Under the dynamics, the mean and variance evolve by an increment ($\mathbb{E} \mathrel{+}= \sum_j n_{ij}\mu_{\ell j}$, $\operatorname{Var} \mathrel{+}= \sum_j n_{ij} v_{\ell j}$), by imperfect repair, or by replacement (reset to the constants of $D^{\mathrm{new}}$).
A repair maps $\mathbb{E} \mapsto (1-\rho)\,\mathbb{E}$ and $\operatorname{Var} \mapsto (1-\rho)^2\operatorname{Var}$ under $\text{ARD}_\infty$, and $\mathbb{E} \mapsto (1-\rho)\,\mathbb{E} + \rho\, g^{\mu}$ and $\operatorname{Var} \mapsto (1-\rho)^2\operatorname{Var} + \bigl(1-(1-\rho)^2\bigr)\, g^{v}$ under $\text{ARD}_1$.
Each operation is non-decreasing in every moment and anchor separately, so the moments evolve as monotone recursions.
Componentwise domination therefore chains: writing $\Pi$ for the map of one repetition of the operating phase, \eqref{eq:loop_moments} states that $\Pi$ does not increase either moment at the state of day~$H_1$, whence the moments at the end of each repetition are at most those at the end of the previous one (Theorem~\ref{thm:perpetual}).
Because the Bernstein bound~\eqref{eq:bernstein_constraint}---and likewise the Cantelli bound~\eqref{eq:cantelli_surrogate}---is increasing in both $\mathbb{E}[D]$ and $\operatorname{Var}[D]$, the certified risk is then non-increasing across periods, $U^{(n)} \le U^{(0)} \le \varepsilon$, and the schedule is repeatable.
Both inequalities in~\eqref{eq:loop_moments} are linear in the tracked moment states $\bigl(\mathbb{E}[D_k],\operatorname{Var}[D_k]\bigr)$, which the model carries as auxiliary variables (big-$M$ rows or indicator constraints encode their recursion linearly), so the loop adds one linear constraint per tracked moment and per anchor; they are strictly stronger than the single relation $U^{(1)} \le U^{(0)}$, and that strengthening is exactly what makes the induction valid.

\paragraph{Looping on the cumulant under imperfect repair.}
If reliability is instead certified by the fixed-$s$ Chernoff bound, the loop acts on the tracked cumulant generating function,
\begin{equation}\label{eq:loop_cumulant}
  K_T \;\le\; K_{H_1}.
\end{equation}
Without imperfect repair, $K_k$ is its value at the last replacement (or on day~$0$) plus the count-weighted cumulant $\sum_j n^{(k)}_{ij}\,\psi_{\ell j}(s)$, where $n^{(k)}_{ij}$ is the post-replacement cumulative count of mission~$j$ up to day~$k$.
There $M_{D}(s)$ propagates multiplicatively, $M_{D + \Delta D}(s) = M_D(s)\,M_{\Delta D}(s)$, so $K$ is a self-propagating scalar and~\eqref{eq:loop_cumulant} inducts.
Imperfect (ARD) repair does not keep this scalar exact. With $D^+ = (1-\rho)D^-$,
\begin{equation}\label{eq:mgf_rescale}
  M_{D^+}(s) \;=\; \mathbb{E}\big[e^{s(1-\rho)D^-}\big] \;=\; M_{D^-}\big((1-\rho)s\big):
\end{equation}
the repaired state's transform at $s$ is the old transform at the \emph{shifted} argument $(1-\rho)s$, which is not determined by $M_{D^-}(s)$---the repair slides the state along the transform curve.
An upper bound, however, is determined by it: the cumulant generating function is convex with value~$0$ at~$0$, so $K_{D^+}(s) \le (1-\rho)\,K_{D^-}(s)$ under $\text{ARD}_\infty$, and $K_{D^+}(s) \le (1-\rho)\,K_{D^-}(s) + \rho\,K_G(s)$ under $\text{ARD}_1$ (Proposition~\ref{prop:rf_maps}).
So the tracked statistic follows $K \mapsto (1-\rho)\,K$ under $\text{ARD}_\infty$ and $K \mapsto (1-\rho)\,K + \rho\,g^K$ under $\text{ARD}_1$, with the anchor~$g^K$.
These maps are non-decreasing in $(K, g^K)$, and the tracked $K$ stays an upper bound on the true cumulant generating function, which is all that the Chernoff certificate needs.
Theorem~\ref{thm:perpetual} then applies: the loop row on $K$, and on its anchor under $\text{ARD}_1$, chains under both repair laws.
The price is slack: after a repair, the tracked $K$ exceeds the true value.
The moments need no such bound, because a repair acts on each of them through a fixed power of $1-\rho$, which keeps their recursion exact.

\paragraph{Summary.}
Under Route~I the loop constraint must be imposed on sufficient, self-propagating statistics, never on the scalar bound~$U$.
The formulation imposes the loop constraint~\eqref{eq:loop} on the statistics of the chosen bound: the mean for Markov; the mean and the variance~\eqref{eq:loop_moments} for Cantelli and Bernstein; the mean and the Hoeffding budget for Hoeffding; the mean and the cumulant generating function~\eqref{eq:loop_cumulant} for Chernoff.
Under $\text{ARD}_1$ the anchors of these statistics enter the loop too.
The loop is valid for every bound and for both repair laws.
In every case the guarantee is non-worsening of the \emph{certified} risk, hence $\Pr(D_k > \tau) \le \varepsilon$ on every day of the endless schedule---the operative meaning of repeatability for a bound-based model.

\subsubsection{Two refinements of the Route-I constraints: schedule-dependent support and piecewise-linear tightening}
\label{sec:routeI_refinements}

Two identifiable sources of slack remain in the Route-I constraints, and both can be removed with standard mixed-integer constructions at a quantifiable price in binaries and rows.
The formulation includes the second one, the piecewise-linear tightening, as the form \texttt{pwl} of its reliability rows (Section~\ref{sec:linear_forms}).
The first one, a schedule-dependent support, is not part of the formulation, which uses the constant support of Section~\ref{sec:cons_surrogates}.
Neither touches the loop argument of Section~\ref{sec:loop_routeI}: the moment condition~\eqref{eq:loop_moments} acts on the exactly tracked moments and is agnostic to how tightly the certificate is enforced on them.\footnote{Every monotonicity and equivalence claim in this subsection, both worked examples, and the union-versus-intersection geometry are checked numerically; the test scripts and outputs are kept in the verification workspace (\texttt{refine\_schedule\_support}, \texttt{refine\_pwl\_tightening}).}

\paragraph{Schedule-dependent Bernstein support.}
The closed form~\eqref{eq:bernstein_constraint} carries the single constant $b_\ell = \max\{b^{\mathrm{start}}_\ell,\, b^{\mathrm{new}}_\ell,\, \max_j b_{\ell j}\}$, which includes the supports $b^{\mathrm{start}}_\ell$, $b^{\mathrm{new}}_\ell$ of the initial and replacement draws, so a component pays the most damaging mission's support even if the schedule never assigns it that mission.
The effect is not marginal, because the correction $\kappa b_\ell/3$ enters twice at zero variance: even the empty schedule satisfies~\eqref{eq:bernstein_constraint} only if $2\kappa b_\ell/3 \le \tau$.
Concretely, take $\varepsilon = 0.01$ ($\kappa = \ln 100 \approx 4.605$), $\tau = 1$, two missions with supports $b_{\ell 1} = 0.02$ and $b_{\ell 2} = 0.30$, reset supports $0.02$, and mission-1 per-day moments $\mu_{\ell 1} = 0.01$, $v_{\ell 1} = 4\cdot 10^{-6}$.
For a component that never serves mission~2, the constant-$b$ constraint ($\kappa b_\ell/3 = 0.4605$) leaves a mean budget of $0.079$ and admits at most $n_{i1} = 7$ mission-1 days, whereas the same constraint evaluated with the supports actually present ($\kappa b_\ell/3 = 0.0307$) leaves $0.939$ and admits $n_{i1} = 90$---a $12.9\times$ gap attributable entirely to the unused $b_{\ell 2}$; and raising $b_{\ell 2}$ to $0.35$ makes $2\kappa b_\ell/3 = 1.075 > \tau$, so the constant-$b$ model becomes infeasible outright for a component that never touches mission~2.
The tightened constant is $b_\ell(\zeta) = \max\bigl\{b^{\mathrm{start}}_\ell,\, b^{\mathrm{new}}_\ell,\, \max_{j:\,\zeta_{ij}=1} b_{\ell j}\bigr\}$, where $\zeta_{ij} \in \{0,1\}$ flags the missions member~$i$ actually serves ($\zeta_{ij} \ge x_{ijt}$ for all~$t$; only this forcing direction is needed): between resets the state unrolls as a $[0,1]$-weighted sum of one reset draw and of increments of \emph{served} missions only, so $b_\ell(\zeta)$ bounds every summand and Bernstein applies verbatim, at every step and under every concatenation of the schedule.
What makes $b_\ell(\zeta)$ implementable without nonlinear max-tracking is a monotonicity observation: the certified value $\mu + c + \sqrt{c^2 + 2\kappa v}$ has $c$-derivative $1 + c/\sqrt{c^2 + 2\kappa v} \in [1,2]$, so it is increasing in the support constant, and the Bernstein row written for a larger constant implies every row written for a smaller one.
Consequently, imposing the row with the per-mission constant $c_j = \kappa\max\{b_{\ell j}, b^{\mathrm{start}}_\ell, b^{\mathrm{new}}_\ell\}/3$ \emph{conditionally on} $\zeta_{ij} = 1$, together with one unconditional row at $c_0 = \kappa\max\{b^{\mathrm{start}}_\ell, b^{\mathrm{new}}_\ell\}/3$, is exactly equivalent to the single row with the schedule-dependent maximum: the active row with the largest constant dominates the rest, and that largest constant is precisely $\kappa b_\ell(\zeta)/3$.
The price is one binary per member--mission pair ($F \cdot M$ in total), their linking rows, and per-mission conditional certificate rows, all linear in the new material: the tangent-surrogate rows attach to $\zeta_{ij}$ by indicator constraints or by big-$M$ (a uniform valid constant is $2\tau(\tau - c_0)$), whereas the exact quadratic rows must use big-$M$, because indicator constraints admit only linear bodies.
Two side effects deserve record.
First, big-$M$ constants derived from the largest admissible variance must be recomputed from $c_0$: the feasible region has grown, and a bound inherited from the old constant $b_\ell$ would cut it.
Second, the feasibility precondition $2\kappa b_\ell/3 < \tau$ localises: it is required only of $c_0$, while a mission with $2c_j \ge \tau$ simply forces $\zeta_{ij} = 0$---that mission is automatically never assigned to that component, instead of rendering the whole model inconsistent.
Deliberately \emph{not} captured: the repair weights further shrink each summand's support to $w_t\,b_{\ell j}$, but exploiting this requires a running maximum over schedule-dependent products of repair factors and mission indicators---a genuinely nonlinear refinement; the window-level $\zeta$ captures the first-order gap, namely missions never served at all.

\paragraph{Piecewise-linear tightening of the tangent surrogates.}
The linear surrogate obtained by fixing one tangent (the device noted after~\eqref{eq:bernstein_constraint}) is conservative in a way that grows quadratically with the mean.
Cantelli's bound~\eqref{eq:cantelli_surrogate} is at most $\varepsilon$ iff $\mathbb{E}[D] + \kappa_C\sqrt{\operatorname{Var}[D]} \le \tau$ with $\kappa_C = \sqrt{(1-\varepsilon)/\varepsilon}$; writing $(\mu, v)$ for the tracked moments, the exact feasible set is $\kappa_C^2\, v \le f(\mu) := (\tau - \mu)^2$, and the tangent at $\mu_0 = 0$ gives the linear row $\kappa_C^2\, v \le \tau^2 - 2\tau\mu$.
Since $f(\mu) - (\tau^2 - 2\tau\mu) = \mu^2$, the surrogate loses exactly $\mu^2/\kappa_C^2$ of admissible variance, reaches zero at $\mu = \tau/2$, and excludes \emph{every} point with $\mu > \tau/2$---a deterministic component included---while the exact set still admits $v = (\tau - \mu)^2/\kappa_C^2 > 0$ up to $\mu = \tau$.
The refinement takes tangents at several points $\mu_r$: each row $\kappa_C^2\, v \le T_r(\mu) := (\tau - \mu_r)(\tau + \mu_r - 2\mu)$ is individually conservative, because tangents lie below the convex parabola.
The directionality is the critical point: imposing all rows \emph{simultaneously} realises their intersection $\kappa_C^2 v \le \min_r T_r(\mu)$, which is \emph{more} conservative than any single tangent; the tightening is the \emph{union} of the single-tangent regions, $\kappa_C^2 v \le \max_r T_r(\mu)$---a disjunction.
Two encodings realise it: (i)~one selector binary per tangent, $\sum_r \lambda_r = 1$, with indicator (or big-$M$) rows $\lambda_r = 1 \Rightarrow \kappa_C^2 v \le T_r(\mu)$; or (ii)~an auxiliary $y$ tied by a piecewise-linear \emph{equality} (SOS2 $\lambda$-formulation) to the pointwise minimum of the tangents to $\sqrt{v}$ at $v_r = (\tau - \mu_r)^2/\kappa_C^2$, followed by the single linear row $\mu + \kappa_C\, y \le \tau$.
The two describe the same set: the $\mu$-side tangent at $\mu_r$ and the $v$-side tangent at $v_r$ bound the identical half-plane $\kappa_C^2 v \le (\tau - \mu_r)(\tau + \mu_r - 2\mu)$.
The tempting shortcuts both fail: requiring $y \ge$ every tangent realises their pointwise \emph{maximum} (a region strictly inside the single-tangent one, worse than not refining), and requiring $y \le$ every tangent lets $y$ under-run $\sqrt{v}$ and destroys the certificate; likewise the equality must interpolate the vertices of the tangent envelope---consecutive tangents meet at $v = \sqrt{v_r v_{r+1}}$ with value $(\sqrt{v_r} + \sqrt{v_{r+1}})/2$, above the curve---not points on the curve, whose chords under-estimate the concave square root.
The recovery is rapid, and exactly quantifiable because $f$ is quadratic: tangent~$r$'s gap is $f(\mu) - T_r(\mu) = (\mu - \mu_r)^2$, so the union's residual gap is $\min_r (\mu - \mu_r)^2$---quadratic in the distance to the nearest tangency, with worst case $(\Delta/2)^2$ for spacing $\Delta$, so doubling the tangents quarters it.
With $\tau = 1$, $\varepsilon = 0.01$ ($\kappa_C^2 = 99$) and tangents at $\mu_r \in \{0, 0.25, 0.5, 0.75\}$, the admissible variance recovers as follows ($v \times 10^{3}$):
\begin{center}
\footnotesize
\begin{tabular}{@{}lccccc@{}}
\toprule
$\mu$ & $0$ & $0.2$ & $0.4$ & $0.5$ & $0.6$ \\
\midrule
exact $(\tau-\mu)^2/\kappa_C^2$ & $10.10$ & $6.46$ & $3.64$ & $2.53$ & $1.62$ \\
single tangent ($\mu_0 = 0$)    & $10.10$ & $6.06$ & $2.02$ & $0$    & infeasible \\
four tangents (union)           & $10.10$ & $6.44$ & $3.54$ & $2.53$ & $1.52$ \\
\bottomrule
\end{tabular}
\end{center}
and the union recovers $94\%$ of the feasible-region area the single tangent forfeits.
The construction transfers verbatim to Bernstein---the parabola $(\tau - \kappa b_\ell/3 - \mu)^2$ obeys the same gap law, and the tangents to the concave $\sqrt{(\kappa b_\ell/3)^2 + 2\kappa v}$ over-estimate it globally---and composes soundly with the schedule-dependent support above: each conditional row carries its own mission constant and remains conservative for it.

\subsubsection{Linear forms of the reliability rows}
\label{sec:linear_forms}

The formulation offers three forms for each quadratic reliability row.
These rows are the Gaussian and Wiener row $\kappa_G^2\, v \le (\tau - \mu)^2$ with $\kappa_G = \Phi^{-1}(1-\varepsilon)$ (Section~\ref{sec:loop_per_model}), and the Cantelli, Hoeffding and Bernstein rows of Table~\ref{tab:five_bounds}.
Each has the form
\begin{equation}\label{eq:quad_row}
  \chi \;\le\; f(\mu), \qquad f(\mu) = c_2\,(\tau - \mu)^2 + c_1\,(\tau - \mu), \qquad \mu \le \bar\mu_{\mathrm{side}},
\end{equation}
where $\chi$ is the variance or the Hoeffding budget, with the constants of Table~\ref{tab:quad_constants}.
The third relation is the side row; it selects the branch of the parabola on which the row is a valid certificate.
Since $c_2 > 0$, $f$ is convex in~$\mu$, with $f'(\mu) = -2c_2(\tau - \mu) - c_1$.

\begin{table}[htbp]
\centering
\caption{Constants of the quadratic reliability rows~\eqref{eq:quad_row}.}
\label{tab:quad_constants}
\small
\begin{tabular}{@{}lllll@{}}
\toprule
\textbf{Row} & $\chi$ & $c_2$ & $c_1$ & $\bar\mu_{\mathrm{side}}$ \\
\midrule
Gaussian, Wiener   & $v$ & $1/\kappa_G^2$ & $0$ & $\tau$ \\
Rainflow, Cantelli & $v$ & $1/\kappa_C^2 = \varepsilon/(1-\varepsilon)$ & $0$ & $\tau$ \\
Rainflow, Hoeffding & $R$ & $2/\kappa$ & $0$ & $\tau$ \\
Rainflow, Bernstein & $v$ & $1/(2\kappa)$ & $-b_\ell/3$ & $\tau - c$ \\
\bottomrule
\end{tabular}
\end{table}

Every form keeps the side row.
\begin{itemize}
  \item \texttt{exact}: the row~\eqref{eq:quad_row} itself.
    The feasible set $\{(\mu, \chi) : \chi \le f(\mu)\}$ lies below a convex function, so it is not convex: the row is a nonconvex quadratic constraint.
    Example (Gaussian row): $(\mu, v) = (\tau, 0)$ and $(0, \tau^2/\kappa_G^2)$ satisfy it, and their midpoint $(\tau/2, \tau^2/(2\kappa_G^2))$ does not, since $\tau^2/2 > \tau^2/4$.
    The solver Gurobi handles such rows through its nonconvex quadratic support.
  \item \texttt{tangent}: one linear row per day, the tangent of~$f$ at a point $\mu_r \in [0, \tau)$ chosen per component; the default is $\mu_r = \tau/2$.
    At $\mu_r = 0$ the Gaussian row becomes $\kappa_G^2\, v + 2\tau\mu \le \tau^2$.
    The point $\mu_r = \tau$ is excluded: for $c_1 = 0$ the tangent at~$\tau$ is $\chi \le 0$.
  \item \texttt{pwl}: the interval $[0, \tau]$ is split into $P$ equal segments, by default $P = 8$.
    One binary per segment and day selects the segment that contains the mean, and the tangent at the midpoint of that segment caps~$\chi$.
    This is encoding~(i) of the piecewise-linear tightening (Section~\ref{sec:routeI_refinements}), with each selector tied to its segment.
    With $P = 1$ the form is the single tangent at $\tau/2$, with no binaries.
\end{itemize}

\begin{proposition}[Safety of the linear forms]\label{prop:tangent_safe}
For every $\mu_r$ and every $\mu$, $f(\mu_r) + f'(\mu_r)\,(\mu - \mu_r) = f(\mu) - c_2\,(\mu - \mu_r)^2 \le f(\mu)$.
Therefore every point that satisfies the \texttt{tangent} or the \texttt{pwl} form satisfies the \texttt{exact} form, and both linear forms lie inside the certificate set $\mathcal{C}_{i\ell}$ of Section~\ref{sec:encoding}.
\end{proposition}
\begin{proof}
$f$ is a polynomial of degree~$2$ in~$\mu$ with leading coefficient~$c_2$, so its Taylor expansion at~$\mu_r$ is exact: $f(\mu) = f(\mu_r) + f'(\mu_r)(\mu - \mu_r) + c_2(\mu - \mu_r)^2$.
Since $c_2 > 0$, the tangent lies below~$f$.
Each linear form caps $\chi$ by one such tangent and keeps the side row.
Therefore it implies~\eqref{eq:quad_row}.
\end{proof}

\paragraph{Accuracy.}
The \texttt{tangent} form loses $c_2(\mu - \mu_r)^2$ in~$\chi$: it is exact at $\mu = \mu_r$ and loosest far from it.
A tangent point near the operating mean of the component is the best choice.
In the \texttt{pwl} form the loss is $c_2(\mu - \hat\mu_p)^2 \le c_2\,\tau^2/(4P^2)$, with $\hat\mu_p$ the midpoint of the segment that contains~$\mu$.
That midpoint is the nearest tangent point to~$\mu$.
So the \texttt{pwl} form accepts exactly the points with $\chi \le \max_p\, [f(\hat\mu_p) + f'(\hat\mu_p)(\mu - \hat\mu_p)]$, the union of the $P$ single-tangent regions and the largest set that the $P$ tangents allow.
The price is $P$ binaries per component and day.

\begin{example}[Cantelli, $\varepsilon = 0.1$, $\kappa_C^2 = 9$, $\tau = 1$, $v = 0.004$]
The largest admissible mean is $1 - \sqrt{9 \cdot 0.004} \approx 0.810$ with \texttt{exact}; $0.482$ with \texttt{tangent} at~$0$ ($9v + 2\mu \le 1$); $0.714$ with \texttt{tangent} at $0.5\,\tau$ ($9v \le 0.75 - \mu$); and $0.810$ with \texttt{pwl} and $P = 8$.
The tangent at~$0$ never admits a mean above $\tau/2$, even with zero variance.
\end{example}

\paragraph{The tangent point of a Bernstein row.}
For Bernstein with $b_\ell > 0$, the tangent point must satisfy $\mu_r < \sqrt{\tau(\tau - 2c)}$, with $c = \kappa b_\ell/3$, and the \texttt{pwl} form needs $\tau/(2P) < \sqrt{\tau(\tau - 2c)}$ for its first segment.
The reason: a tangent is linear in~$\mu$, so its largest value on $[0, \tau - c]$ is at an end point.
At $\tau - c$ it is at most $f(\tau - c) = -c^2/(2\kappa) < 0$.
At~$0$ it equals $f(0) - \mu_r^2/(2\kappa) = \bigl(\tau(\tau - 2c) - \mu_r^2\bigr)/(2\kappa)$, by Proposition~\ref{prop:tangent_safe} with $c_2 = 1/(2\kappa)$.
So the tangent row admits a point with $v \ge 0$ exactly when $\mu_r^2 \le \tau(\tau - 2c)$, and at equality only $\mu = v = 0$.
Since $\sqrt{\tau(\tau - 2c)} < \tau - c$, this covers the case $\mu_r \ge \tau - c$, in which the tangent has a non-negative slope.
The \texttt{pwl} form accepts the points below the largest of its tangents, a convex function; at $\mu = 0$ the largest is the tangent at the midpoint $\tau/(2P)$ of the first segment, so the same argument applies with $\mu_r = \tau/(2P)$.
Otherwise no segment admits a point with $v > 0$: the first segment fails as above, and on every later segment $\mu \ge \tau/P > \tau - 2c$, where $f(\mu) < 0$.

\subsubsection{Random fatigue life in the Palmgren--Miner rule}
\label{sec:miner_randomN}

The rainflow pipeline of Section~\ref{sec:cons_surrogates} measures the law of the per-mission increment $\Delta D^{(j)}$ directly and nonparametrically.
A complementary pipeline keeps the rainflow \emph{counts} deterministic and makes the fatigue life $N_f$ in Miner's rule~\eqref{eq:miner} the random object: one fits a scatter family of Table~\ref{tab:sn_scatter} to constant-amplitude test data or in-service failure records, tying the location to the stress amplitude through Basquin's law~\eqref{eq:basquin}---for the log-normal, $\ln N_f \sim \mathcal{N}(\mu(\sigma_a),\sigma^2)$ with $\mu(\sigma_a) = \ln C - m \ln \sigma_a$ and a stress-independent $\sigma$; for the Weibull, scale $\lambda(\sigma_a) = C\sigma_a^{-m}$ and common shape $k$; run-outs enter the likelihood as censored observations---and then \emph{derives} the increment from the mission's rainflow histogram $\{n^{(j)}_c\}$ as $\Delta D^{(j)} = \sum_c n^{(j)}_c / N_f(\sigma_c)$.
The retrieval is the same in spirit---collect data, approximate a distribution---but dual in object and commitment: the empirical route is distribution-free and captures load and material variability jointly, whereas the random-$N_f$ route is parametric, pools information across stress levels, extrapolates to unobserved amplitudes through~\eqref{eq:basquin}, and separates the load spectrum (the counts) from the material scatter (the law of $N_f$).

\paragraph{Inverse moments of the scatter families.}
Everything the framework consumes downstream is a moment of $1/N_f$.
For the three families of Table~\ref{tab:sn_scatter} and $r > 0$,
\begin{equation}\label{eq:miner_invmom}
  \mathbb{E}\big[N_f^{-r}\big] \;=\;
  \begin{cases}
    \exp\big({-r\mu + r^2\sigma^2/2}\big), & \ln N_f \sim \mathcal{N}(\mu,\sigma^2) \quad \text{(finite for all $r$)},\\[3pt]
    \lambda^{-r}\,\Gamma(1 - r/k), & N_f \sim \mathrm{Wei}(k,\lambda) \quad \text{(finite iff $r < k$)},\\[3pt]
    e^{-r\mu}\,\Gamma(1 - r\beta), & \ln N_f \sim \mathrm{Gumbel}_{\min}(\mu,\beta) \quad \text{(finite iff $r\beta < 1$)}.
  \end{cases}
\end{equation}
The first line is the normal moment generating function at argument $-r$; the second follows from the substitution $u = (x/\lambda)^k$, which turns $\mathbb{E}[N_f^{-r}]$ into $\lambda^{-r}\int_0^\infty u^{-r/k}e^{-u}\,\mathrm{d}u$, divergent \emph{at the origin} when $r \ge k$; the third writes $\ln N_f = \mu + \beta G$ with $\Pr(G \le g) = 1 - e^{-e^{g}}$, whose transform $\mathbb{E}[e^{tG}] = \int_0^\infty u^{t}e^{-u}\,\mathrm{d}u = \Gamma(1+t)$ (substitute $u = e^g$) is finite iff $t > -1$.
In particular $\operatorname{Var}[1/N_f]$ equals $e^{-2\mu+\sigma^2}(e^{\sigma^2}-1)$, $\lambda^{-2}\big[\Gamma(1-2/k)-\Gamma(1-1/k)^2\big]$, and $e^{-2\mu}\big[\Gamma(1-2\beta)-\Gamma(1-\beta)^2\big]$ respectively, matching the increment moments already recorded for the log-normal below Table~\ref{tab:sn_scatter}.

\begin{remark}[Two tail behaviours, one real restriction]\label{rem:miner_tails}
The Gumbel-minimum log-life is \emph{exactly} a Weibull life: $\Pr(N_f > x) = \exp\big(-e^{(\ln x - \mu)/\beta}\big) = \exp\big(-(x/e^{\mu})^{1/\beta}\big)$, i.e.\ $\mathrm{Wei}(k = 1/\beta,\ \lambda = e^{\mu})$, and the existence conditions in~\eqref{eq:miner_invmom} map onto each other.
The three families of Table~\ref{tab:sn_scatter} therefore realise two tail behaviours for $1/N_f$: log-normal (all polynomial moments finite) and Fr\'echet-tailed with index $k$ (moments capped at order $k$).
The latter is a genuine restriction, not a formality: the mean of the increment requires $k > 1$ and the variance---the input of the Cantelli certificate~\eqref{eq:cantelli_surrogate}---requires $k > 2$, so the fitted shape must be checked before the machinery is invoked; for $k \in (1,2]$ only the Markov bound survives, and for $k \le 1$ the untruncated model has an infinite mean increment and cannot enter the framework at all.
Moreover, for \emph{all three} untruncated families $1/N_f$ has unbounded support and an infinite moment generating function, so the exponential certificates---Chernoff~\eqref{eq:chernoff_permission}, Bernstein~\eqref{eq:bernstein_constraint}, Hoeffding---are inadmissible; they are restored only by a physically defended lower life bound $N_{\min} > 0$ (truncated family), which yields the per-mission support $b_{\ell j} = n_j / N_{\min}$ and requires the moments to be recomputed under the truncated law (Gaussian-tail and incomplete-gamma closed forms).
\end{remark}

\paragraph{What ``random $N_f$'' means: two interpretations.}
The marginal law of $N_f$ does not determine the model; the joint law of the scatter \emph{across days} must be chosen, and the two extremes lead to different mathematics.

\emph{(I) Block-independent scatter.}
Each execution of mission~$j$ draws a fresh, independent life---appropriate when cycle-to-cycle and day-to-day variability dominates and the component does not carry one fixed $N_f$.
With a single dominant stress level per mission (for a binned spectrum the means add over bins, while the variance depends on the within-mission dependence assumption and is computed from~\eqref{eq:miner_invmom} in the same way), the per-mission inputs are
\begin{equation}\label{eq:miner_routeI}
  \mu_{\ell j} \;=\; n_j\,\mathbb{E}\big[N_f^{-1}\big],
  \qquad
  v_{\ell j} \;=\; n_j^2\,\operatorname{Var}\big[N_f^{-1}\big],
\end{equation}
and independence across days is precisely the hypothesis of Equation~\eqref{eq:additive_functionals}: the Cantelli certificate~\eqref{eq:cantelli_surrogate} and the moment loop~\eqref{eq:loop_moments} apply \emph{verbatim}, with parametric instead of empirical per-mission moments.
This is the interpretation the rainflow route of Section~\ref{sec:cons_surrogates} already embodies; the random-$N_f$ model is then nothing but an alternative input pipeline for the same formulation, subject to the admissibility conditions of Remark~\ref{rem:miner_tails}.

\emph{(II) Unit-specific scatter.}
The component carries \emph{one} realisation of the scatter for its entire life.
Under the stress-independent scale/shape of the fitted families this is a single multiplicative factor,
\begin{equation}\label{eq:miner_unit}
  N_{f,j} \;=\; A\,\bar N_j, \qquad A > 0 \ \text{drawn once at installation},
\end{equation}
with $\bar N_j$ the median (log-normal, Gumbel: $e^{\mu_j}$) or characteristic (Weibull: $\lambda_j$) life and $A \sim F_A$ the unit factor ($\ln A \sim \mathcal{N}(0,\sigma^2)$, $A \sim \mathrm{Wei}(k,1)$, or $A = e^{\beta G}$ respectively).
This is exactly the ``random critical damage sum'' of source~(ii) above: dividing the Miner sum by $A$ is the same as comparing a deterministic sum to a random capacity.
The damage state becomes $D_k = M_k / A$, where $M_k$ is the \emph{deterministic} Miner state with dynamics mirroring~\eqref{eq:dyn_full}: a mission day adds $\sum_j x_{ijk}\, n_j/\bar N_j$, ARD repair maps $M \mapsto (1-\rho)M$, and replacement resets $M$ to zero \emph{and redraws} $A$ (a new physical component carries a fresh factor with the same law).
For continuous $F_A$ the reliability constraint is then \emph{exact}, deterministic, and linear in the assignment counts:
\begin{equation}\label{eq:miner_quantile}
  \Pr(D_k \ge 1) \;=\; \Pr(A \le M_k) \;=\; F_A(M_k) \;\le\; \varepsilon
  \quad\Longleftrightarrow\quad
  M_k \;\le\; q_\varepsilon \;:=\; F_A^{-1}(\varepsilon),
\end{equation}
with closed-form quantiles $q_\varepsilon = e^{\sigma\Phi^{-1}(\varepsilon)}$ (log-normal), $q_\varepsilon = (-\ln(1-\varepsilon))^{1/k}$ (Weibull), $q_\varepsilon = (-\ln(1-\varepsilon))^{\beta}$ (Gumbel-minimum); indeed no family is needed at all beyond one scalar---an empirical $\varepsilon$-quantile of $A$ with a confidence margin serves equally.
Rewriting $M_k \le q_\varepsilon$ as $\sum_j n_{ij}/(q_\varepsilon \bar N_j) \le 1$ recovers the classical design practice of running deterministic Miner on the $\varepsilon$-percentile lives $N^{(\varepsilon)}_j = q_\varepsilon \bar N_j$ (these are exactly the $\varepsilon$-quantiles of $N_{f,j}$)---here as an \emph{exact} chance constraint rather than a heuristic, in the per-step sense of~\eqref{eq:abs_rel}.
The loop constraint collapses likewise: $F_A$ is strictly increasing on $(0,\infty)$, so the probability comparison~\eqref{eq:loop_prob} is equivalent to the single linear condition
\begin{equation}\label{eq:miner_loop_det}
  M_T \;\le\; M_{H_1},
\end{equation}
and---unlike the lossy scalar bound of Section~\ref{sec:loop_routeI}---it \emph{chains}: $M$ is a sufficient statistic (the failure probability is the fixed increasing function $F_A$ of the scalar state), and the map $\Pi$ of the operating phase is a composition of $m \mapsto m + c$ ($c \ge 0$), $m \mapsto (1-\rho)m$ and $m \mapsto 0$, each non-decreasing, so from $M_T = \Pi(M_{H_1}) \le M_{H_1}$ induction gives $\Pi^{n+1}(M_{H_1}) \le \Pi^{n}(M_{H_1})$ for all $n$; the same monotonicity applied to the partial flows dominates every within-period trajectory by the trajectory of the operating phase on the horizon, so imposing~\eqref{eq:miner_quantile} on all $k = 1,\dots,T$ extends the per-step reliability to every concatenation, replacements (which redraw $A$) included.

\begin{remark}[The dichotomy is not a conservatism knob]\label{rem:miner_dichotomy}
Interpretations~(I) and~(II) can share the \emph{same} per-day marginal law and differ only in dependence---zero versus perfect correlation of the scatter across days---yet they answer differently.
Means coincide, while by Cauchy--Schwarz $\operatorname{Var}_{\mathrm{II}}/\operatorname{Var}_{\mathrm{I}} = \big(\sum_t m_t\big)^2 / \sum_t m_t^2$ (the number of days, for equal daily increments $m_t$): under~(II) the dispersion never averages out.
In the small-$\varepsilon$ regime~(II) is by far the heavier-tailed---in a representative experiment (log-normal factor, $\sigma = 0.5$, fifty days, median Miner sum $0.6$) the failure probability was $0.15$ under~(II) and $\sim 10^{-6}$ under~(I)---while at high accumulated damage the ordering flips, so neither interpretation is uniformly conservative.\footnote{Test code and output (\texttt{miner\_a}--\texttt{miner\_c}) are kept in the verification workspace, covering the closed forms~\eqref{eq:miner_invmom}, the divergent sample variance at Weibull $k \le 2$, the exactness of~\eqref{eq:miner_quantile} along a schedule with repair and replacement, the chaining of~\eqref{eq:miner_loop_det}, and the size of the (I)/(II) gap.}
The choice is an identification question: scatter traced to specimen-to-specimen variability (material batch, manufacturing) is unit-specific~(II); scatter from load-history randomness within a mission is block-independent~(I); rainflow-counted in-service data from a fleet mixes both, and a between-unit/within-unit variance decomposition of the recorded increments estimates the split.
When both sources matter one writes $N_f = A\,\widetilde N_f$ with a unit factor and block-independent residuals; conservative composites exist---e.g.\ $\varepsilon$-splitting, $\Pr(\text{fail}) \le \delta + \Pr(\text{fail} \mid A = a_\delta)$ with $a_\delta$ the $\delta$-quantile of $A$, after which the increments are again independent across days with moments scaled by $a_\delta^{-1}$ and $a_\delta^{-2}$ and the Route-I machinery runs on the residual budget $\varepsilon - \delta$---but their calibration is left as an outlook.
\end{remark}

\paragraph{Implications for the abstract formulation.}
Interpretation~(I) requires no change to Problem~\eqref{eq:abstract}: it supplies the per-mission moments~\eqref{eq:miner_routeI} to the same Route-I pipeline that consumes the rainflow-measured ones, and inherits its certificates and loop condition unchanged.
Interpretation~(II) defines a strictly simpler model family: the stochastic state collapses to a deterministic scalar $M_{i\ell k}$ with linear dynamics, the chance constraint becomes the exact linear inequality~\eqref{eq:miner_quantile}, and the loop becomes the linear condition~\eqref{eq:miner_loop_det}---no moment tracking, no tail-bound slack, and no big-$M$ linearisation of a variance recursion---with the entire stochastic description compressed into one quantile $q_\varepsilon$ per component.
It is therefore a natural candidate for a dedicated (deterministic-equivalent) model type alongside the process-based ones, with the caveat of Remark~\ref{rem:miner_dichotomy}: which interpretation matches a given component is an empirical question, not a formulation preference.

\subsection{Stochastic orderings}
\label{sec:stoch_orders}

Comparing the degradation state of a component at different points in the schedule requires a notion of ``one random variable being larger than another.''
Several stochastic orders formalise this idea with increasing strength~\cite{ShakedShanthikumar2007,WikiStochDom}.
Let $X$ and $Y$ be two non-negative random variables with CDFs $F_X$, $F_Y$, survival functions $\bar F_X = 1-F_X$, $\bar F_Y = 1-F_Y$, and densities $f_X$, $f_Y$ (when they exist).

\begin{definition}[Usual stochastic order]\label{def:st}
  $X \le_{\mathrm{st}} Y$ if $F_X(t) \ge F_Y(t)$ for all $t \in \mathbb{R}$,
  or equivalently $\Pr(X > t) \le \Pr(Y > t)$ for all~$t$.
\end{definition}

\begin{definition}[Hazard rate order]\label{def:hr}
  $X \le_{\mathrm{hr}} Y$ if $\bar F_Y(t)/\bar F_X(t)$ is non-decreasing in $t$ over $\{t : \bar F_X(t) > 0\}$.
  Equivalently, the hazard rates satisfy $r_X(t) \ge r_Y(t)$ for all~$t$, where $r(t) = f(t)/\bar F(t)$.
\end{definition}

\begin{definition}[Likelihood ratio order]\label{def:lr}
  $X \le_{\mathrm{lr}} Y$ if $f_Y(t)/f_X(t)$ is non-decreasing in $t$ over the union of their supports.
\end{definition}

These three orders are progressively stronger~\cite[Ch.~1]{ShakedShanthikumar2007}.
The implications are strict---none of the converses holds in general:
\begin{equation}\label{eq:order_implications}
  X \le_{\mathrm{lr}} Y
  \;\Longrightarrow\;
  X \le_{\mathrm{hr}} Y
  \;\Longrightarrow\;
  X \le_{\mathrm{st}} Y
  \;\Longrightarrow\;
  \mathbb{E}[X] \le \mathbb{E}[Y].
\end{equation}

\begin{remark}[Connection to the loop constraint]
The probability comparison~\eqref{eq:loop_prob} requires $\Pr(D_{i\ell T} > \tau) \le \Pr(D_{i\ell H_1} > \tau)$, which is a \emph{pointwise} comparison of the survival functions at the single threshold~$\tau$.
This is strictly weaker than $D_{i\ell T} \le_{\mathrm{st}} D_{i\ell H_1}$, which demands the inequality at every point.
Establishing the usual stochastic order is therefore a convenient sufficient condition for~\eqref{eq:loop_prob}, though not a necessary one.
The loop constraint~\eqref{eq:loop} of the formulation compares tracked statistics instead (Section~\ref{sec:loop_constraint}).
\end{remark}

\subsubsection{Stochastic ordering within each degradation family}

For the fleet management problem, the relevant comparison is between two damage states of the same component---for instance the state at the end of the transitory phase versus the end of the horizon---both belonging to the same parametric family.
Tables~\ref{tab:ordering_lr}--\ref{tab:ordering_st} give, for each family from Table~\ref{tab:deg_models}, the necessary and sufficient conditions for each of the three orderings.
An entry marked ``---'' indicates that no simple closed-form characterisation is known to the author; for the inverse Gaussian rows the $\le_{\mathrm{lr}}$ conditions remain a valid sufficient criterion via~\eqref{eq:order_implications}, while for the compound Poisson rows see the discussion below.

\begin{table}[htbp]
\centering
\caption{Conditions for $X \le_{\mathrm{lr}} Y$ (likelihood ratio order) within each family.  Conditions are necessary and sufficient except where noted.}
\label{tab:ordering_lr}
\small
\begin{tabular}{@{}p{2.1cm}p{3.8cm}p{4.0cm}p{4.0cm}@{}}
\toprule
\textbf{Family} & $X$ & $Y$ & \textbf{Condition} \\
\midrule
Gamma
  & $\mathrm{Ga}(\alpha_1,\,\beta_1)$
  & $\mathrm{Ga}(\alpha_2,\,\beta_2)$
  & $\alpha_1 \le \alpha_2$ and $\beta_1 \ge \beta_2$ \\[4pt]
Gaussian
  & $\mathcal{N}(\mu_1,\,\sigma_1^2)$
  & $\mathcal{N}(\mu_2,\,\sigma_2^2)$
  & $\mu_1 \le \mu_2$ and $\sigma_1 = \sigma_2$ \\[4pt]
Inv.\ Gaussian
  & $\mathrm{IG}(\mu_1,\,\lambda_1)$
  & $\mathrm{IG}(\mu_2,\,\lambda_2)$
  & $\lambda_2 \ge \lambda_1$ and $\lambda_1/\mu_1^2 \ge \lambda_2/\mu_2^2$ \\[4pt]
Comp.\ Poisson
  & $\mathrm{CP}(\nu_1,\,F_{Y_1})$
  & $\mathrm{CP}(\nu_2,\,F_{Y_2})$
  & --- (can fail even for $F_{Y_1}{=}F_{Y_2}$, $\nu_1 \le \nu_2$; holds, e.g., for deterministic jumps; see text) \\
\bottomrule
\end{tabular}
\end{table}

\begin{table}[htbp]
\centering
\caption{Conditions for $X \le_{\mathrm{hr}} Y$ (hazard rate order) within each family.}
\label{tab:ordering_hr}
\small
\begin{tabular}{@{}p{2.1cm}p{3.8cm}p{4.0cm}p{4.0cm}@{}}
\toprule
\textbf{Family} & $X$ & $Y$ & \textbf{Condition} \\
\midrule
Gamma
  & $\mathrm{Ga}(\alpha_1,\,\beta_1)$
  & $\mathrm{Ga}(\alpha_2,\,\beta_2)$
  & $\alpha_1 \le \alpha_2$ and $\beta_1 \ge \beta_2$ \\[4pt]
Gaussian
  & $\mathcal{N}(\mu_1,\,\sigma_1^2)$
  & $\mathcal{N}(\mu_2,\,\sigma_2^2)$
  & $\mu_1 \le \mu_2$ and $\sigma_1 = \sigma_2$ \\[4pt]
Inv.\ Gaussian
  & $\mathrm{IG}(\mu_1,\,\lambda_1)$
  & $\mathrm{IG}(\mu_2,\,\lambda_2)$
  & --- (suf.: $\le_{\mathrm{lr}}$ conditions) \\[4pt]
Comp.\ Poisson
  & $\mathrm{CP}(\nu_1,\,F_{Y_1})$
  & $\mathrm{CP}(\nu_2,\,F_{Y_2})$
  & --- (no general criterion; the $\le_{\mathrm{lr}}$ route is unavailable, see text) \\
\bottomrule
\end{tabular}
\end{table}

\begin{table}[htbp]
\centering
\caption{Conditions for $X \le_{\mathrm{st}} Y$ (usual stochastic order) within each family.}
\label{tab:ordering_st}
\small
\begin{tabular}{@{}p{2.1cm}p{3.8cm}p{4.0cm}p{4.0cm}@{}}
\toprule
\textbf{Family} & $X$ & $Y$ & \textbf{Condition} \\
\midrule
Gamma
  & $\mathrm{Ga}(\alpha_1,\,\beta_1)$
  & $\mathrm{Ga}(\alpha_2,\,\beta_2)$
  & $\alpha_1 \le \alpha_2$ and $\beta_1 \ge \beta_2$ \\[4pt]
Gaussian
  & $\mathcal{N}(\mu_1,\,\sigma_1^2)$
  & $\mathcal{N}(\mu_2,\,\sigma_2^2)$
  & $\mu_1 \le \mu_2$ and $\sigma_1 = \sigma_2$ \\[4pt]
Inv.\ Gaussian
  & $\mathrm{IG}(\mu_1,\,\lambda_1)$
  & $\mathrm{IG}(\mu_2,\,\lambda_2)$
  & --- (suf.: $\le_{\mathrm{lr}}$ conditions) \\[4pt]
Comp.\ Poisson
  & $\mathrm{CP}(\nu_1,\,F_{Y_1})$
  & $\mathrm{CP}(\nu_2,\,F_{Y_2})$
  & $\nu_1 \le \nu_2$, $F_{Y_1} = F_{Y_2}$, $Y \ge 0$ \\
\bottomrule
\end{tabular}
\end{table}

\paragraph{Gamma.}
The log-likelihood ratio of $\mathrm{Ga}(\alpha_2,\beta_2)$ over $\mathrm{Ga}(\alpha_1,\beta_1)$ has derivative $(\alpha_2 - \alpha_1)/x - (\beta_2 - \beta_1)$, which is non-negative for all $x > 0$ if and only if $\alpha_2 \ge \alpha_1$ and $\beta_2 \le \beta_1$.
Conversely, if $\alpha_1 > \alpha_2$ then $F_X(t) \sim (\beta_1 t)^{\alpha_1}/(\alpha_1\,\Gamma(\alpha_1)) \ll (\beta_2 t)^{\alpha_2}/(\alpha_2\,\Gamma(\alpha_2)) \sim F_Y(t)$ as $t \to 0^+$, violating $\le_{\mathrm{st}}$;
if $\beta_1 < \beta_2$ then $\bar F_X(t)/\bar F_Y(t) \to \infty$ as $t \to \infty$, again violating $\le_{\mathrm{st}}$.
Hence all three orders coincide for the Gamma family~\cite[Sec.~1.C]{ShakedShanthikumar2007}.
Under the closure condition (same~$\beta$), they all reduce to $\alpha_1 \le \alpha_2$.

\paragraph{Gaussian.}
The CDF of $\mathcal{N}(\mu,\sigma^2)$ is $\Phi\bigl((t-\mu)/\sigma\bigr)$.
The inequality $\Phi\bigl((t-\mu_1)/\sigma_1\bigr) \ge \Phi\bigl((t-\mu_2)/\sigma_2\bigr)$ for all~$t$ requires $(t-\mu_1)/\sigma_1 \ge (t-\mu_2)/\sigma_2$ for all~$t$, which is linear in~$t$ and can hold identically only if $\sigma_1 = \sigma_2$, in which case it reduces to $\mu_1 \le \mu_2$.
The likelihood ratio is then $\exp\bigl((\mu_2-\mu_1)\,x/\sigma^2 + \text{const}\bigr)$, non-decreasing in~$x$.
Thus all three orders coincide: $\sigma_1 = \sigma_2$ and $\mu_1 \le \mu_2$.

\paragraph{Inverse Gaussian.}
The log-likelihood ratio of $\mathrm{IG}(\mu_2,\lambda_2)$ over $\mathrm{IG}(\mu_1,\lambda_1)$ has derivative
\[
  \frac{\lambda_1}{2\mu_1^2} - \frac{\lambda_2}{2\mu_2^2} + \frac{\lambda_2 - \lambda_1}{2x^2}.
\]
This is non-negative for all $x > 0$ if and only if $\lambda_2 \ge \lambda_1$ (from $x \to 0^+$) and $\lambda_1/\mu_1^2 \ge \lambda_2/\mu_2^2$ (from $x \to \infty$).
Under the closure condition (same $\mu^2/\lambda$), these simplify to $\mu_1 \le \mu_2$.
For the general two-parameter case, closed-form necessary and sufficient conditions for $\le_{\mathrm{hr}}$ and $\le_{\mathrm{st}}$ do not appear in the standard references; the $\le_{\mathrm{lr}}$ conditions serve as a sufficient criterion.

\paragraph{Compound Poisson.}
With the same jump-size distribution~$F_Y$ ($Y \ge 0$), $\mathrm{CP}(\nu_1, F_Y) \le_{\mathrm{st}} \mathrm{CP}(\nu_2, F_Y)$ if and only if $\nu_1 \le \nu_2$ (provided $\Pr(Y>0) > 0$): sufficiency follows from the superposition coupling $\mathrm{CP}(\nu_2, F_Y) \stackrel{d}{=} \mathrm{CP}(\nu_1, F_Y) + \mathrm{CP}(\nu_2 - \nu_1, F_Y)$ with independent non-negative summands, and necessity from the atom at zero, $\Pr(X = 0) = e^{-\nu \Pr(Y > 0)}$.
The stronger orders do \emph{not} follow in general: with $Y$ uniform on $\{1, 10\}$ and $\nu_1 = 0.1$, $\nu_2 = 0.2$, the likelihood ratio drops from ${\approx}463$ at $x = 9$ (reachable only by nine unit jumps) to ${\approx}1.8$ at $x = 10$ (reachable by a single large jump), so $\le_{\mathrm{lr}}$ fails, and the survival ratio is likewise non-monotone, so $\le_{\mathrm{hr}}$ fails as well; both hold in special cases such as deterministic jump sizes, where the process is a scaled Poisson.
When the jump-size distributions differ, the ordering conditions depend on the specific form of~$F_Y$ and no universal closed-form characterisation is available.

\paragraph{Implications for the fleet problem.}
Under the closure conditions of Table~\ref{tab:discrete_models}, the cumulative damage $D_{i\ell k}$ depends on the schedule only through cumulative parameters.
For the one-parameter families the comparison of two damage states is scalar:
for Gamma (shared rate~$\beta$), the total shape $\sum_t \alpha_t$;
for IG (shared dispersion $\mu^2/\lambda$), the total mean $\sum_t \mu_t$.
For the Gaussian family the state carries \emph{two} cumulative parameters, $\bigl(\sum_t \mu_t,\ \sum_t \sigma_t^2\bigr)$, and the variance is not schedule-independent: it accumulates with the mission days served and is reduced only by maintenance.
Since the usual stochastic order between Gaussians requires equal variances (Table~\ref{tab:ordering_st}), the Gaussian comparison proceeds through the single-threshold condition~\eqref{eq:loop_gauss} or through componentwise domination of mean and variance, not through the mean alone.

The formulation does not rely on closure for the Gamma and inverse Gaussian models, for two reasons.
First, the terms of the damage need not share the closure parameter.
The missions of one component may have different Gamma rates, and an inverse Gaussian term of mean~$a$ in the formulation has the dispersion $\mu^2/\lambda = \eta\, a$, which grows with~$a$ (the model is stated below).
Second, a repair acts on the damage path, which leaves both families: $(1-\rho)\,\mathrm{Ga}(A,\beta) = \mathrm{Ga}\bigl(A,\,\beta/(1-\rho)\bigr)$, and $c\,\mathrm{IG}(\mu,\lambda) = \mathrm{IG}(c\mu,\,c\lambda)$ multiplies the dispersion $\mu^2/\lambda$ by~$c$.
For these two models the formulation therefore tracks the mean only.
It bounds the failure probability over a whole class of laws that contains every possible law of the damage, and it imposes a cap on the tracked mean.
The inverse Gaussian cap uses a further order.

\begin{definition}[Increasing convex order]\label{def:icx}
For integrable non-negative $X$ and $Y$, $X \le_{\mathrm{icx}} Y$ if $\mathbb{E}(X-d)_+ \le \mathbb{E}(Y-d)_+$ for every $d \ge 0$.
\end{definition}

The quantity $\mathbb{E}(X-d)_+$ is the stop-loss transform of~$X$ at~$d$.
This order does not compare tails directly: $X \le_{\mathrm{icx}} Y$ allows $\Pr(X > \tau) > \Pr(Y > \tau)$.
For example, the constant $X = 1$ and the variable $Y$ equal to $0$ or $2$ with probability $1/2$ each satisfy $X \le_{\mathrm{icx}} Y$, while $\Pr(X > 0.9) = 1 > \Pr(Y > 0.9) = 1/2$.

\paragraph{Inverse Gaussian: a cap on the mean.}
The inverse Gaussian law $\mathrm{IG}(\mu, \lambda)$ with mean $\mu > 0$ and shape $\lambda > 0$ has the density
\[
  f(x; \mu, \lambda) = \sqrt{\frac{\lambda}{2\pi x^3}}\,\exp\!\Bigl(-\frac{\lambda\,(x - \mu)^2}{2\mu^2 x}\Bigr), \qquad x > 0,
\]
and the variance $\mu^3/\lambda$; for $c > 0$, $c\,\mathrm{IG}(\mu, \lambda) = \mathrm{IG}(c\mu, c\lambda)$.
The formulation gives each inverse Gaussian component one dispersion $\eta > 0$.
A term of mean $a$ has the law $\mathrm{IG}(a, a/\eta) = a\,Z$, with $Z \sim \mathrm{IG}(1, 1/\eta)$.
Its density is $\sqrt{a/(2\pi\eta x^3)}\,\exp\bigl(-(x - a)^2/(2\eta a x)\bigr)$, and $\eta$ is its squared coefficient of variation.
Every term has this law with its own mean: each mission increment, the initial damage and each replacement draw.
A known damage~$a$ is the constant term $a \cdot 1$.
The terms may depend on each other in any way.
Between replacements, under either repair law, the damage is $D_k = \sum_t w_t\, a_t\, Z_t$, where each $Z_t$ has the law of~$Z$ or equals~$1$, and each weight $w_t \in [0,1]$ is a product of factors $1-\rho$.
A term of weight~$0$, possible when $\rho = 1$, is dropped.
$\text{ARD}_\infty$ multiplies every weight by $1-\rho$, and $\text{ARD}_1$ multiplies only the weights of the terms added since the last intervention.
When the anchor before the first intervention is below the initial mean~$\mu_0$, the initial damage $\mu_0 Z_0$ is read as the two terms $g\,Z_0$ and $(\mu_0 - g)\,Z_0$, with $g$ the anchor.
The mean $\sum_t w_t a_t$ is at most the tracked mean $\mu_k$ (Proposition~\ref{prop:domination}).

\begin{proposition}[Stop-loss order]\label{prop:ig_order}
Let $D = \sum_t w_t\, a_t\, Z_t$ with weights $w_t \in [0,1]$, means $a_t \ge 0$ and $\sum_t w_t a_t \le \mu$, where each $Z_t$ has the law of~$Z$ or equals~$1$, under any joint law.
Then $D \le_{\mathrm{icx}} \mu Z$.
\end{proposition}
\begin{proof}
Let $c_t = w_t a_t$ and $A = \sum_t c_t$.
If $A = 0$, then $D = 0$.
Otherwise, the map $x \mapsto (x-d)_+$ is convex and non-decreasing, and $(A - d)_+ \le \mathbb{E}(AZ - d)_+$ by Jensen's inequality, since $\mathbb{E}Z = 1$.
Hence, for every $d \ge 0$,
\[
  \mathbb{E}(D - d)_+
  = \mathbb{E}\Bigl(\sum_t \tfrac{c_t}{A}\,(A Z_t) - d\Bigr)_+
  \le \sum_t \tfrac{c_t}{A}\,\mathbb{E}(A Z_t - d)_+
  \le \mathbb{E}(A Z - d)_+
  \le \mathbb{E}(\mu Z - d)_+ . \qedhere
\]
\end{proof}

Let $q$ be the $(1-\varepsilon)$-quantile of~$Z$, so $\Pr(Z > q) = \varepsilon$, and let $\bar q = \mathbb{E}[Z \mid Z > q]$, the conditional value at risk (CVaR) of~$Z$ at the level $1-\varepsilon$.

\begin{theorem}[CVaR cap]\label{thm:ig_cap}
Let $D \ge 0$ satisfy $D \le_{\mathrm{icx}} \mu Z$ with $0 \le \mu \le \tau/\bar q$.
Then $\Pr(D > \tau) \le \varepsilon$.
\end{theorem}
\begin{proof}
If $\mu = 0$, then $\mathbb{E}D \le 0$ and $D = 0$.
Let $\mu > 0$ and $d = \mu q$; then $d < \mu\bar q \le \tau$.
For every $x$, $\mathbf{1}\{x > \tau\} \le (x-d)_+/(\tau - d)$, and $\mathbb{E}(Z - q)_+ = \varepsilon\,(\bar q - q)$.
Therefore
\[
  \Pr(D > \tau)
  \le \frac{\mathbb{E}(D-d)_+}{\tau - d}
  \le \frac{\mu\,\mathbb{E}(Z-q)_+}{\tau - \mu q}
  = \frac{\varepsilon\,\mu\,(\bar q - q)}{\tau - \mu q}
  \le \varepsilon,
\]
where the last step is $\mu\bar q \le \tau$.
\end{proof}

So the formulation tracks the mean of an inverse Gaussian component and imposes the linear row $\mu_k \le \bar\mu = \tau/\bar q$ on every day, under both repair laws.
By Proposition~\ref{prop:ig_order} and Theorem~\ref{thm:ig_cap}, the row gives $\Pr(D_k > \tau) \le \varepsilon$, whatever the dependence between the terms.
The quantile~$q$ follows from the distribution function~\eqref{eq:ig_cdf} of $Z$, with $\mu = 1$ and $\lambda = 1/\eta$.
The tail mean has the closed form $\mathbb{E}[Z;\, Z > q] = F_Z(1/q)$, so $\bar\mu = \tau\,\varepsilon/F_Z(1/q)$.
Derivation: substituting into the density $f_Z$ of~$Z$ gives $z\,f_Z(z) = z^{-2} f_Z(1/z)$, and the change of variable $u = 1/z$ gives $\int_q^\infty z\, f_Z(z)\, \mathrm{d}z = \int_0^{1/q} f_Z(u)\, \mathrm{d}u = F_Z(1/q)$.
Table~\ref{tab:ig_cap} gives values of the cap.

\begin{table}[htbp]
\centering
\caption{The CVaR cap $\bar\mu/\tau$ of the inverse Gaussian model, for three reliability levels and five dispersions.}
\label{tab:ig_cap}
\small
\begin{tabular}{@{}lccccc@{}}
\toprule
 & $\eta = 0.05$ & $0.25$ & $1$ & $4$ & $10$ \\
\midrule
$\varepsilon = 0.001$ & 0.491 & 0.242 & 0.101 & 0.0370 & 0.0190 \\
$\varepsilon = 0.01$ & 0.570 & 0.322 & 0.156 & 0.0679 & 0.0401 \\
$\varepsilon = 0.1$ & 0.693 & 0.477 & 0.299 & 0.183 & 0.141 \\
\bottomrule
\end{tabular}
\end{table}

No larger cap is safe for every law below $\mu Z$ in the order of Definition~\ref{def:icx}: for $\mu > \bar\mu$ and a level $\varepsilon' > \varepsilon$ close to~$\varepsilon$, with quantile~$q'$, still $\mu\,\mathbb{E}[Z \mid Z > q'] > \tau$.
The two-point law $\mu\,\mathbb{E}\bigl[Z \mid \mathbf{1}\{Z > q'\}\bigr]$ lies below $\mu Z$ in this order, by Jensen's inequality, and exceeds~$\tau$ with probability~$\varepsilon'$.
The schedules produce only some of the laws below $\mu Z$, so the row can reject some safe schedules.
The quantile cap $\tau/q$, at which the single law $\mu Z$ meets the target exactly, is not safe for sums.
Example: $\eta = 10$ and $\varepsilon = 0.1$, so $\tau/q = 0.473\,\tau$; three independent missions of equal mean, with total mean $\tau/q$, give $\Pr(D > \tau) = 0.118$.
The CVaR cap is $0.141\,\tau$ there.
At $\varepsilon = 0.01$ and $\eta$ from $0.05$ to $10$, the CVaR cap $\bar\mu$ is $0.93$ to $0.59$ times $\tau/q$.

\paragraph{Gamma: a sharp cap on the mean.}
The formulation lets each mission of a Gamma component have its own rate: an increment of mean $m > 0$ and rate~$\beta$ has the law $\mathrm{Ga}(\beta m, \beta)$.
The initial damage and each replacement draw are Gamma terms with their own rates, or known constants.
All terms are independent.
A repair acts on the damage path by the law of the component: $\text{ARD}_\infty$ gives $D^+ = (1-\rho)\,D^-$, and $\text{ARD}_1$ gives $D^+ = D^{\mathrm{last}} + (1-\rho)\,(D^- - D^{\mathrm{last}})$.
After a repair the damage is, in general, not a Gamma variable: it is a weighted sum of independent Gamma terms and known constants.
The cap below holds for every such sum.
The day map of the tracked mean is exact for the true mean $\mathbb{E}[D_k]$, because expectation is linear.
Let $\beta^*$ be the smallest rate among the terms of the component with a positive mean, over the missions and the days of both phases, and $t = \beta^*\tau$.
The formulation computes $\beta^*$ from the data; no input gives it.
A mission with mean~$0$ adds no term, and a known damage is a constant, which does not enter~$\beta^*$; the rate of the replacement draw enters only when replacements are allowed.
If no term has a rate, the damage is a known constant, and the cap below reads $\mu_k \le \tau$.
Units of $1/\beta^*$ turn every law below into a law of rate at least~$1$.
Write $Q(a, x) = \Pr(\mathrm{Ga}(a,1) > x)$ for the regularised upper incomplete gamma function.

\begin{definition}[Class of weighted Gamma sums]\label{def:gamma_class}
For $m \ge 0$, the class $\mathcal{G}(m)$ holds the laws of the random variables $S = \delta + \sum_{i=1}^{n} c_i\, Y_i$ with $Y_i \sim \mathrm{Ga}(a_i, 1)$ independent, $c_i \in (0,1]$, $a_i > 0$, $\delta \ge 0$, $n \ge 0$, and mean $\delta + \sum_i c_i a_i = m$.
The term $c\,Y$ with $Y \sim \mathrm{Ga}(m'/c, 1)$ has mean~$m'$; as $c \to 0$ it tends to the constant~$m'$.
\end{definition}

\begin{lemma}[True states belong to the class]\label{lem:gamma_class}
On every day~$k$ of a schedule, under either repair law, the law of $\beta^* D_k$ belongs to $\mathcal{G}(\beta^*\,\mathbb{E}[D_k])$.
\end{lemma}
\begin{proof}
Since the last replacement (or day~$0$), $D_k = \delta + \sum_t w_t X_t$: the $X_t$ are the independent Gamma terms added since then, $\delta \ge 0$ collects the known constants, and each weight $w_t \in [0,1]$ is a product of factors $1-\rho$; a term of weight~$0$, possible when $\rho = 1$, is dropped.
A mission adds a term with weight~$1$.
$\text{ARD}_\infty$ multiplies every weight, constants included, by $1-\rho$.
$\text{ARD}_1$ multiplies by $1-\rho$ the weights of the terms and constants added since the last intervention, and keeps the others.
A replacement starts a new sum with the replacement draw.
When the anchor before the first intervention is below the initial mean, the initial damage is read as two independent terms of its rate, with means equal to the anchor and to the rest; their sum has the law of the initial damage.
A known initial damage is read as two constants in the same way.
Write $X_t = Y_t/\beta_t$ with $Y_t \sim \mathrm{Ga}(a_t, 1)$ and $\beta_t \ge \beta^*$.
Then $\beta^* D_k = \beta^*\delta + \sum_t (w_t\beta^*/\beta_t)\, Y_t$, and every scale $w_t\beta^*/\beta_t$ lies in $(0,1]$.
The $Y_t$ are independent, so the law of $\beta^* D_k$ is in $\mathcal{G}(\beta^*\,\mathbb{E}[D_k])$.
\end{proof}

For $0 < m < t$, let $W(m, t)$ be the supremum of $\Pr(S > t)$ over the random variables $S$ whose law is in $\mathcal{G}(m)$.
The one-term member $c\,\mathrm{Ga}(m/c, 1)$ has the tail $Q(m/c,\, t/c)$ at~$t$.
A member of $\mathcal{G}(m)$ is described by its \emph{atoms}: the pairs $(m_i, c_i)$ of a term mean $m_i = c_i a_i$ and a scale~$c_i$, and the constant~$\delta$, read as an atom at scale $c = 0$.
Terms with equal scales merge into one term, because a sum of independent Gamma variables of one rate is a Gamma variable.

\begin{lemma}[First-order rates]\label{lem:gamma_rates}
Let $S$ be a member with at least one term, with density~$f_S$.
Since $\delta \le m < t$, $f_S$ is continuous at~$t$.
For $c \in (0, 1]$ let
\[
  \omega_S(c) = \frac{1}{c}\int_0^\infty \Pr(t - y < S \le t)\,\frac{e^{-y/c}}{y}\, \mathrm{d}y,
  \qquad \omega_S(0) = f_S(t).
\]
(a) Adding mean $\Delta m$ at scale~$c$ (a term $c\,\mathrm{Ga}(\Delta m/c, 1)$, or the constant $\Delta m$ when $c = 0$) raises $\Pr(S > t)$ at the rate $\omega_S(c)$ as $\Delta m \to 0$; removing mean from an atom at scale~$c$ lowers it at the same rate.
(b) Moving the scale of an atom $(m_i, c_i)$, $c_i \in (0,1)$, at fixed mean~$m_i$ changes $\Pr(S > t)$ at the rate $m_i\, \partial\omega_S/\partial c\,(c_i)$.
\end{lemma}
\begin{proof}
$S$ is infinitely divisible, with drift~$\delta$ and L\'evy density $\Lambda(y) = \sum_i (m_i/c_i)\, e^{-y/c_i}/y$ on $y > 0$.
A change of the parameters at the rates $(\dot\delta, \dot\Lambda)$ changes the Laplace transform $\mathbb{E}e^{-\vartheta S} = \exp\bigl(-\vartheta\delta - \int_0^\infty (1 - e^{-\vartheta y})\,\Lambda(y)\, \mathrm{d}y\bigr)$ at the rate $\int_0^\infty \mathbb{E}\bigl[e^{-\vartheta(S+y)} - e^{-\vartheta S}\bigr]\,\dot\Lambda(y)\, \mathrm{d}y - \vartheta\dot\delta\,\mathbb{E}e^{-\vartheta S}$.
This is the transform of the rate $\int_0^\infty \Pr(t - y < S \le t)\,\dot\Lambda(y)\, \mathrm{d}y + \dot\delta\, f_S(t)$ of $\Pr(S > t)$.
Case~(a) is $\dot\Lambda(y) = c^{-1}e^{-y/c}/y$, or $\dot\delta = 1$ when $c = 0$.
Case~(b) is $\dot\Lambda(y) = m_i\,\partial_c\bigl(c^{-1}e^{-y/c}/y\bigr)$ at $c = c_i$.
\end{proof}

\begin{theorem}[One term is the worst case]\label{thm:gamma_oneterm}
For $0 < m < t$,
\[
  W(m, t) \;=\; \max_{c \in (0,1]} Q\!\left(\frac{m}{c},\, \frac{t}{c}\right).
\]
\end{theorem}
\begin{proof}
Let $\mathcal{G}_N(m)$ be the members with at most $N$ atoms, the constant included, and let $N \ge 3$.
\begin{enumerate}
  \item \emph{A maximiser exists.} $\mathcal{G}_N(m)$ is the image of the compact set of atoms $(m_i, c_i)$ with $m_i \ge 0$, $\sum_i m_i = m$, $c_i \in [0, 1]$.
    The law of~$S$ depends continuously on the atoms, and only the constant member $S = m$ has an atom of probability, at $m \ne t$.
    So $\Pr(S > t)$ is continuous and attains its maximum at some~$S^*$.
    Since $\Pr(m > t) = 0 < Q(m, t)$, $S^*$ has at least one term.
    Every member has finitely many atoms, so $W(m, t) = \sup_N \max_{\mathcal{G}_N(m)} \Pr(S > t)$.
  \item \emph{First-order conditions.} Write $\omega^\circ = \omega_{S^*}$.
    By Lemma~\ref{lem:gamma_rates}: (i) moving mean between two atoms of positive mean must not raise the tail, so $\omega^\circ$ takes one value $\omega^*$ at every atom of positive mean; (ii) at an atom with $c \in (0, 1)$, $\partial\omega^\circ/\partial c = 0$; (iii) if $S^*$ has fewer than $N$ atoms of positive mean, a new atom can be opened at any $c \in [0, 1]$, so $\omega^\circ(c) \le \omega^*$ on $[0, 1]$.
  \item \emph{A single-peak function.} Put $\gamma = 1/c$ and $h(u) = \Pr(t - u < S^* \le t)/u$.
    Then $\omega^\circ(\gamma) = \gamma\int_0^\infty h(u)\, e^{-\gamma u}\, \mathrm{d}u$, and $\omega^\circ(\infty) = h(0^+) = f_{S^*}(t)$ is the value at the constant atom.
    $S^*$ is a constant plus a sum of independent Gamma variables, so its density is unimodal~\cite[Theorem~4, p.~11]{RoostaKhorasaniSzekely2015}, proved in~\cite[proof of Theorem~4, pp.~194--195]{SzekelyBakirov2003}; the constant shifts the mode only.
    Then $p(v) = f_{S^*}(t - v)$ for $v < t$ and $p(v) = 0$ for $v \ge t$ is unimodal on $[0, \infty)$, and so is its running average $h(u) = u^{-1}\int_0^u p$: before the mode of~$p$, $p \ge h$; after it, $u\,(p(u) - h(u)) = u\,p(u) - \int_0^u p$ is non-increasing.
    Since $h' = (p - h)/u$, $h$ rises and then falls.
  \item \emph{Zero count.} For every constant $\omega'$, $h - \omega'$ changes sign at most twice.
    The Laplace transform of a function with at most $n$ sign changes has at most $n$ zeros in $\gamma > 0$, counted with multiplicity: at a sign change~$u_1$, multiply the transform by $e^{\gamma u_1}$; its derivative in~$\gamma$ is the transform of a function with one sign change fewer; Rolle's theorem and induction conclude.
    So $\omega^\circ - \omega' = \gamma\int_0^\infty (h - \omega')\, e^{-\gamma u}\, \mathrm{d}u$ has at most two zeros in $\gamma > 0$.
    Also $\omega^\circ(\gamma) \to 0$ as $\gamma \to 0^+$.
    Hence $\omega^\circ$ cannot rise above a level $\omega' > 0$, fall below it and rise above it again: with $\omega^\circ(0^+) = 0$ that would give three zeros.
    So $\omega^\circ$ is quasi-concave.
    It is analytic in $\gamma > 0$ and not constant, so its maximum over $\gamma \in [1, \infty]$ is attained at one point only: two maximisers would make $\omega^\circ$ constant between them.
  \item \emph{One term.} By (i) and (ii), an atom with $c \in (0, 1)$ is a zero of $\omega^\circ - \omega^*$ of multiplicity at least~$2$, and an atom at $c = 1$ a zero of multiplicity at least~$1$; the constant atom is at $\gamma = \infty$.
    By step~4, $S^*$ has at most two atoms of positive mean.
    As $N \ge 3$, (iii) applies, so $\omega^*$ is the maximum of $\omega^\circ$ over $[1, \infty]$, attained at every atom of positive mean.
    By step~4 there is only one such atom, and by step~1 it is not the constant.
    So $S^* = c^*\,\mathrm{Ga}(m/c^*, 1)$ with $c^* \in (0, 1]$.
\end{enumerate}
Therefore $\max_{\mathcal{G}_N(m)} \Pr(S > t) = \max_{c \in (0,1]} Q(m/c, t/c)$ for every $N \ge 3$, and the claim follows from step~1.
\end{proof}

\begin{lemma}[Monotone in the mean]\label{lem:gamma_monotone}
$W(m, t)$ is non-decreasing in $m \in (0, t)$.
\end{lemma}
\begin{proof}
Let $m < m' < t$ and let the law of $S$ be in $\mathcal{G}(m)$.
Then the law of $S + (m' - m)$ is in $\mathcal{G}(m')$, and $\Pr(S + m' - m > t) \ge \Pr(S > t)$.
Therefore $W(m, t) \le W(m', t)$.
\end{proof}

For $A \in [0, t)$ let $w(A) = \max_{c \in (0,1]} Q(A/c, t/c)$; it is continuous and non-decreasing in~$A$, because $Q(a, x)$ increases in~$a$.
Also $w(A) < 1/2$ for every $A < t$: with $r = A/t < 1$ and $s = t/c$, $Q(rs, s) \le Q(rs, rs) < 1/2$, because the median of a Gamma law lies below its mean.
And $w(A) \to 1/2$ as $A \to t$, because $Q(s, s) \to 1/2$ as $s \to \infty$.
So, for $\varepsilon < 1/2$, the \emph{sharp cap} $A_1(t,\varepsilon)$ is the largest $A \in [0, t)$ with $w(A) \le \varepsilon$, and it lies below~$t$.
For $\varepsilon \ge 1/2$ no largest value exists, because every $A < t$ qualifies; the cap is then a value below~$t$ fixed by the computation.
The formulation tracks the mean of a Gamma component and imposes, on every day and under both repair laws, the linear row
\begin{equation}\label{eq:gamma_cap}
  \mu_k \;\le\; \frac{A_1(\beta^*\tau,\, \varepsilon)}{\beta^*}.
\end{equation}

\begin{proposition}[The sharp cap is safe]\label{prop:gamma_safe}
If the tracked mean $\mu_k$ satisfies~\eqref{eq:gamma_cap} and $\mu_k \ge \mathbb{E}[D_k]$, then $\Pr(D_k > \tau) \le \varepsilon$, for any mix of rates.
\end{proposition}
\begin{proof}
Let $m = \beta^*\,\mathbb{E}[D_k]$ and $A = A_1(\beta^*\tau, \varepsilon)$, so $m \le \beta^*\mu_k \le A < t$.
If $m = 0$, then $D_k = 0$.
Otherwise
\[
  \Pr(D_k > \tau) = \Pr(\beta^* D_k > t) \le W(m, t) \le W(A, t) = w(A) \le \varepsilon,
\]
by Lemma~\ref{lem:gamma_class}, Lemma~\ref{lem:gamma_monotone}, Theorem~\ref{thm:gamma_oneterm} and the choice of~$A$: in both cases of the definition, $A < t$ and $w(A) \le \varepsilon$.
\end{proof}

\paragraph{Computation of the cap.}
The cap $A_1$ is computed once per component, before the model is built, by bisection on~$A$ with an inner maximum over~$c$.
Three facts justify the computation.
First, the exact one-rate cap $\hat A$, the root of $Q(\hat A, t) = \varepsilon$, is found by bisection on its first argument, because $Q$ increases in its first argument; the inverse routine \texttt{gammaincinv} solves for the second argument and does not apply.
Second, the function $c \mapsto Q(A/c, t/c)$ tends to~$0$ as $c \to 0$ and has a single peak on $(0, 1]$; this was checked numerically for $t$ from $10^{-3}$ to $100$ and $A/t$ from $0.05$ to $0.999$.
So a logarithmic grid of values of~$c$, refined by a golden-section search between the neighbours of the best grid point, finds the maximum.
Third, the bisection on~$A$ starts from the lower end $\varepsilon t$, which satisfies $w(\varepsilon t) \le \varepsilon$ by Markov's inequality $Q(a, x) \le a/x$; it returns the lower end of its last interval, so $w(A_1) \le \varepsilon$ holds at the returned value.

\paragraph{The sharp cap and the exact one-rate cap.}
$\hat A$ is the exact cap of a single Gamma law of rate~$\beta^*$, and $A_1 \le \hat A$, because $c = 1$ is in the maximum.
Let $t_1(\varepsilon)$ be the value of~$t$ at which $\partial_c Q(\hat A/c, t/c)$ at $c = 1$ changes sign: $t_1 = 0.436$, $0.448$ and $0.614$ for $\varepsilon = 0.001$, $0.01$ and $0.1$.
For $t \ge t_1$, $A_1 = \hat A$: on a grid of~$t$ in $[t_1, 1000]$, no $c < 1$ gives a larger tail than $c = 1$.
Below $t_1$, $A_1 < \hat A$; for $\varepsilon = 0.01$:
\begin{center}
\small
\begin{tabular}{@{}lccccc@{}}
\toprule
$t = \beta^*\tau$ & 0.01 & 0.1 & 0.3 & 1 & 10 \\
\midrule
$A_1/\hat A$ & 0.142 & 0.642 & 0.961 & 1 & 1 \\
\bottomrule
\end{tabular}
\end{center}
When every term has the same rate and no repair occurs, the true state is a single Gamma law and $\hat A$ is exact for it.
With several rates, or with repairs, the true state is a weighted sum, and below~$t_1$ the exact one-rate cap is not safe.
Example: $\varepsilon = 0.01$ and $\beta^*\tau = 0.1$, so $\hat A = 0.0055012$.
A new part serves one mission of mean $\hat A/\beta^*$ and rate $4.56\,\beta^*$.
Its mean meets the exact one-rate cap, but its tail is $1.565\,\varepsilon$.
The sharp cap there is $A_1 = 0.642\,\hat A$.

\paragraph{Below the threshold.}
With $r = A/t$ and $s = t/c$, $Q(A/c, t/c) = Q(rs, s)$, and $c \in (0, 1]$ means $s \ge t$.
For $t \le t_1(\varepsilon)$ the peak of $s \mapsto Q(rs, s)$ at the cap lies at $s \ge t$, so the maximum does not depend on~$t$, and $A_1(t, \varepsilon) = r^*(\varepsilon)\, t$.
The mean cap is then $A_1/\beta^* = r^*(\varepsilon)\,\tau$, whatever the rates: $r^* = 0.00355$, $0.0353$ and $0.331$ for $\varepsilon = 0.001$, $0.01$ and $0.1$.

\paragraph{The choice of $\beta^*$.}
A larger~$t$ shrinks the range $s \ge t$, so $A_1(t, \varepsilon)/t$ is non-decreasing in~$t$, and the mean cap $A_1(\beta^*\tau, \varepsilon)/\beta^* = \tau\, A_1(t, \varepsilon)/t$ is non-decreasing in~$\beta^*$.
Lemma~\ref{lem:gamma_class} holds for every $\beta^*$ at most the smallest rate.
So the smallest rate gives the largest safe cap of this form.
The class $\mathcal{G}(m)$ also contains laws that no schedule produces, so below~$t_1$ the row can reject some safe schedules.
The row keeps each mission's own mean; $\beta^*$ only bounds how dispersed each term may be.

\subsubsection{Loop constraint for each degradation model}
\label{sec:loop_per_model}

The loop constraint~\eqref{eq:loop} is one rule for every model: each tracked statistic, and its anchor under $\text{ARD}_1$, must be no larger at the end of the horizon than at the end of the transitory phase.
Table~\ref{tab:loop_per_model} lists the tracked statistics and the reliability row of each model of the formulation.
Every model accepts both repair laws, and Theorem~\ref{thm:perpetual} applies to each of them.
We drop the $(i,\ell)$ subscripts for clarity.

\begin{table}[htbp]
\centering
\caption{Tracked statistics and reliability rows of the five degradation models of the formulation, on every day~$k$ of the horizon.
Under $\text{ARD}_1$ each tracked statistic also has an anchor, except the Wiener variance.
The loop constraint~\eqref{eq:loop} covers every tracked statistic and every anchor.
$\kappa_G = \Phi^{-1}(1-\varepsilon)$, with $\varepsilon < 0.5$.}
\label{tab:loop_per_model}
\small
\begin{tabular}{@{}lll@{}}
\toprule
\textbf{Model} & \textbf{Tracked statistics} & \textbf{Reliability row} \\
\midrule
Gaussian         & mean $\mu$, variance $v$ & $\kappa_G^2\, v_k \le (\tau - \mu_k)^2$, \quad $\mu_k \le \tau$ \\
Wiener           & mean $\mu$, variance $v$ & as for the Gaussian model \\
Inverse Gaussian & mean $\mu$ & $\mu_k \le \tau/\bar q$ (Theorem~\ref{thm:ig_cap}) \\
Gamma            & mean $\mu$ & $\mu_k \le A_1(\beta^*\tau, \varepsilon)/\beta^*$, Eq.~\eqref{eq:gamma_cap} \\
Rainflow         & mean $\mu$ and the statistic of the bound & Table~\ref{tab:five_bounds} \\
\bottomrule
\end{tabular}
\end{table}

\paragraph{Gamma.}
With one shared rate~$\beta$ and no repair, the damage is $\mathrm{Ga}(A_k,\,\beta)$, where $A_k$ is the cumulative shape at step~$k$.
Its survival function is $\Pr(D > \tau) = Q(A,\,\beta\tau)$, with $Q$ the regularised upper incomplete gamma function of Section~\ref{sec:stoch_orders}.
Since $Q(a,x)$ is strictly increasing in~$a$ for fixed $x > 0$, a loop on the shape would then be a scalar inequality.
A caveat on maintenance: the \emph{pathwise} ARD repair $D \mapsto (1-\rho)D$ leaves the shared-rate sub-family, since $(1-\rho)\,\mathrm{Ga}(A,\beta) = \mathrm{Ga}\bigl(A,\,\beta/(1-\rho)\bigr)$.
A repair modelled in parameter space, as the shape reduction $A \mapsto (1-\rho)A$ at fixed rate~$\beta$, matches ARD in expectation, $\mathbb{E}[D^+] = (1-\rho)\,\mathbb{E}[D^-]$, but is a distinct operator pathwise, and it can accept an unsafe schedule.
Example: $\tau = 1$, $\varepsilon = 0.01$, $\rho = 0.5$ and one rate $\beta = 0.01$, so the exact one-rate cap is $\hat A = 0.0024877$.
A new part serves a mission of shape~$\hat A$ on day~1, is repaired on day~2, and serves a mission of shape $0.5\,\hat A$ on day~3.
The shapes in parameter space are $\hat A$, $0.5\,\hat A$ and $\hat A$, so every cap row holds, but the true tail on day~3 is $0.0133$, that is $1.33\,\varepsilon$.
The formulation therefore repairs the damage path, tracks the mean, and certifies reliability by the sharp cap~\eqref{eq:gamma_cap}, which is safe for every weighted sum of Gamma terms (Proposition~\ref{prop:gamma_safe}).
Its loop constraint acts on the mean and, under $\text{ARD}_1$, on the anchor of the mean.
Facts (M1)--(M3) of Section~\ref{sec:encoding} hold: (M1) the repair factor $1-\rho$ lies in $[0, 1]$, and the mean increments are non-negative; (M2) the true mean equals the day map applied to the true statistics of the day before, and the starting values satisfy $\theta_0 \le \xi_0$ as for the Gaussian model below; (M3) if a value $\mu \le A_1/\beta^*$ is at least the true mean of day~$k$, then $\Pr(D_k > \tau) \le \varepsilon$, because the proof of Proposition~\ref{prop:gamma_safe} uses only these two facts.
The Gaussian family is closed under both readings of the repair, $(1-\rho)\,\mathcal{N}(\mu,\sigma^2) = \mathcal{N}\bigl((1-\rho)\mu,\,(1-\rho)^2\sigma^2\bigr)$.

\paragraph{Gaussian.}
Let $D_{H_1} \sim \mathcal{N}(\mu_{H_1},\,\sigma_{H_1}^2)$ and $D_T \sim \mathcal{N}(\mu_T,\,\sigma_T^2)$.
The survival function is $\Pr(D > \tau) = \Phi\bigl((\mu - \tau)/\sigma\bigr)$, where $\Phi$ is the standard normal CDF.
Since $\Phi$ is strictly increasing, the probability comparison~\eqref{eq:loop_prob} is equivalent to
\begin{equation}\label{eq:loop_gauss}
  \frac{\mu_T - \tau}{\sigma_T} \;\le\; \frac{\mu_{H_1} - \tau}{\sigma_{H_1}},
\end{equation}
or equivalently $(\tau - \mu_T)/\sigma_T \ge (\tau - \mu_{H_1})/\sigma_{H_1}$.
Note that $\sigma_T = \sigma_{H_1}$ does \emph{not} follow from serving the same missions in both phases: absent maintenance the variance accumulates, $\sigma^2_T = \sigma^2_{H_1} + \sum_{k = H_1+1}^{T} \sigma^2_k > \sigma^2_{H_1}$; equality requires the operating phase's repairs or replacements (which scale the variance by $(1-\rho)^2$ or reset it) to remove exactly the variance the mission days add, in which case~\eqref{eq:loop_gauss} reduces to $\mu_T \le \mu_{H_1}$.
In general, condition~\eqref{eq:loop_gauss} couples the mean and standard deviation: a higher variance at~$T$ can be compensated by a sufficiently lower mean.
Because~\eqref{eq:loop_gauss} constrains only the single scalar $(\mu - \tau)/\sigma$ of a two-parameter state, it does not by itself make the schedule repeatable (cf.\ Section~\ref{sec:loop_routeI}); repeatability is restored by the componentwise condition $\mu_T \le \mu_{H_1}$ and $\sigma^2_T \le \sigma^2_{H_1}$---the exact-law analogue of~\eqref{eq:loop_moments}---which implies~\eqref{eq:loop_gauss} whenever $\mu_{H_1} < \tau$ and chains because both moments propagate monotonically under the dynamics.
This componentwise condition is the loop constraint~\eqref{eq:loop} of the Gaussian model, together with the rows on the two anchors under $\text{ARD}_1$.

In the formulation, a mission increment of a Gaussian component is $\mathcal{N}(\mu^{\mathrm{inc}}, v^{\mathrm{inc}})$, the initial damage is $\mathcal{N}(\mu_0, v_0)$, and each replacement draw is $\mathcal{N}(\mu^{\mathrm{new}}, v^{\mathrm{new}})$; a known damage has variance~$0$.
These terms are mutually independent (assumption N1).
A repair acts on the damage path by the law of the component.
The Gaussian law gives negative damage a small probability; the input checks $\mu \ge 3\sqrt{v}$ for the initial damage and for each increment keep it at most $\Phi(-3) = 0.135\,\%$ for each of these terms.

\begin{proposition}[The state stays Gaussian]\label{prop:gauss_exact}
Under N1, on every day of a schedule $D_k \sim \mathcal{N}(m_k, V_k)$, where $(m_k, V_k)$ follow the day maps of the mean and the variance, with factors $1-\rho$ and $(1-\rho)^2$, and with anchors equal to the mean and the variance just after the last intervention.
\end{proposition}
\begin{proof}
A mission adds an independent Gaussian term, and $\text{ARD}_\infty$ scales the damage by $1-\rho$; both keep the law Gaussian and give the first two cases.
Under $\text{ARD}_1$, $D^- = G + \Delta$, with $G$ the damage just after the last intervention and $\Delta$ the increments added since; $G$ and $\Delta$ are independent Gaussian variables by N1.
Then $D^+ = G + (1-\rho)\Delta$ is Gaussian, with mean $m_G + (1-\rho)(m^- - m_G)$ and variance
\[
  V_G + (1-\rho)^2\, V_\Delta = V_G + (1-\rho)^2\,(V^- - V_G) = (1-\rho)^2\, V^- + \bigl(1 - (1-\rho)^2\bigr)\, V_G .
\]
A replacement starts from the replacement draw.
Therefore the law is Gaussian on every day, with the moments of the maps.
\end{proof}

So $\text{ARD}_1$ is exact for this model; its variance map uses the factor $(1-\rho)^2$ and the variance anchor.
For $D \sim \mathcal{N}(m, V)$, $\Pr(D > \tau) = 1 - \Phi\bigl((\tau - m)/\sqrt{V}\bigr) \le \varepsilon$ exactly when $m + \kappa_G\sqrt{V} \le \tau$, with $\kappa_G = \Phi^{-1}(1-\varepsilon)$.
For $\varepsilon < 0.5$, so that $\kappa_G > 0$, and with $\tau - \mu \ge 0$, this is equivalent to the squared row of Table~\ref{tab:loop_per_model}.
\emph{Validity.} Let $(m_k, V_k) \le (\mu_k, v_k)$ be the true moments (Proposition~\ref{prop:gauss_exact} and Proposition~\ref{prop:domination}).
If $V_k = 0$, then $D_k = m_k \le \tau$.
Otherwise $\mu_k < \tau$, because $\mu_k = \tau$ forces $v_k = 0$, and $\Pr(D_k > \tau) = 1 - \Phi\bigl((\tau - m_k)/\sqrt{V_k}\bigr) \le 1 - \Phi\bigl((\tau - \mu_k)/\sqrt{v_k}\bigr) \le \varepsilon$, because $(\tau - \mu)/\sqrt{v}$ decreases in~$\mu$ and in~$v$.

\begin{remark}[The level $\varepsilon$ and the range $V_{\max}$]\label{rem:vmax}
$V_{\max} = \tau^2/\kappa_G^2$ is the largest variance that the Gaussian row admits, at $\mu = 0$.
So it is the natural range of the variance, and the big-M constant of its rows (Section~\ref{sec:impl_encoding}).
$V_{\max}/\tau^2$ is $0.18$ at $\varepsilon = 0.01$, $0.61$ at $0.1$, $3.6$ at $0.3$ and about $1600$ at $0.49$.
It grows without bound as $\varepsilon \to 0.5$, and at $\varepsilon \ge 0.5$ the squared row does not describe the reliability requirement, because $\kappa_G \le 0$.
Hence the condition $\varepsilon < 0.5$ for Gaussian and Wiener components; the range is large only close to~$0.5$.
The other models and bounds have no limit below $\varepsilon = 1$: for example, the variance cap $\tau^2\varepsilon/(1-\varepsilon)$ of the Cantelli row is finite for every $\varepsilon < 1$.
The variance ranges of the rainflow rows behave like $V_{\max}$ at the other end: the Cantelli and Bernstein variance ranges and the Hoeffding range $2\tau^2/\kappa$ grow without bound as $\varepsilon \to 1$.
\end{remark}

Facts (M1)--(M3) of Section~\ref{sec:encoding} hold.
(M1) The repair factors $1-\rho$ and $(1-\rho)^2$ lie in $[0, 1]$, and the mean and variance increments are non-negative.
(M2) $\theta_0 \le \xi_0$: the initial values are the moments of the initial damage, and each starting anchor is at least the moment of the part of the damage present just after the last intervention before day~$1$ (Section~\ref{sec:repair_connection}); on every later day the true moments equal the maps applied to the true statistics of the day before (Proposition~\ref{prop:gauss_exact}).
(M3) If values $(\mu, v)$ that satisfy the row are at least the true moments of day~$k$, then $\Pr(D_k > \tau) \le \varepsilon$ (the validity paragraph).

\paragraph{Wiener.}
The damage is a Wiener process with drift, observed once a day.
A mission adds a Gaussian increment with a mission-dependent mean and the variance~$\sigma^2$, the same for every mission, because $\sigma$ describes the material of the component, not the route.
The initial damage is $\mathcal{N}(\mu_0, v_0)$ with $v_0 \ge 0$, and each replacement draw is $\mathcal{N}(\mu^{\mathrm{new}}, v^{\mathrm{new}})$; these terms are mutually independent (assumption W1).
In the formulation a repair acts on the mean only, a modelling choice: it shifts the damage by a known amount, $D^+ = D^- - \rho\, m^-$ under $\text{ARD}_\infty$ and $D^+ = D^- - \rho\,(m^- - g^{\mu})$ under $\text{ARD}_1$, with $m^- = \mathbb{E}D^-$.
The law stays Gaussian, the variance is unchanged, and only a replacement lowers it.
Under W1 the damage is $\mathcal{N}(m_k, V_k)$ on every day, with $(m_k, V_k)$ given by these maps: a mission adds an independent Gaussian term, and a repair or a replacement acts as stated.
The reliability row is the Gaussian one, with the same validity argument, forms and remark on~$V_{\max}$.
A repair of the damage path, as in the Gaussian model, would give the same mean and the variance $(1-\rho)^2 V$ under $\text{ARD}_\infty$, or $(1-\rho)^2 V + (1 - (1-\rho)^2)\,V_G$ under $\text{ARD}_1$, with $V_G \le V$ the variance just after the last intervention.
Both are at most~$V$, and the reliability row increases in~$v$, so the rows of the Wiener model are safe for that reading as well.
Facts (M1)--(M3) hold as for the Gaussian model, with the factor $1$ for the variance, which has no anchor.
The loop row on the variance forces a replacement in the operating phase whenever the member serves a mission there: without one, $v_T \ge v_{H_1} + n\,\sigma^2 > v_{H_1}$ after $n \ge 1$ mission days.
No input check enforces this; when replacements are not allowed, the member of a Wiener component cannot serve a mission in the operating phase.

\paragraph{Inverse Gaussian.}
The CDF of $\mathrm{IG}(\mu,\lambda)$ admits the closed form
\begin{equation}\label{eq:ig_cdf}
  F_{\mathrm{IG}}(x;\,\mu,\lambda) \;=\;
    \Phi\!\Biggl(\sqrt{\frac{\lambda}{x}}\Bigl(\frac{x}{\mu}-1\Bigr)\Biggr)
    \;+\; e^{2\lambda/\mu}\;
    \Phi\!\Biggl(-\sqrt{\frac{\lambda}{x}}\Bigl(\frac{x}{\mu}+1\Bigr)\Biggr).
\end{equation}
For $D_{H_1} \sim \mathrm{IG}(\mu_{H_1},\,\lambda_{H_1})$ and $D_T \sim \mathrm{IG}(\mu_T,\,\lambda_T)$, the probability comparison~\eqref{eq:loop_prob} reads
\[
  F_{\mathrm{IG}}(\tau;\,\mu_T,\lambda_T) \;\ge\; F_{\mathrm{IG}}(\tau;\,\mu_{H_1},\lambda_{H_1}).
\]
This is an analytical but nonlinear inequality in the four parameters $(\mu_{H_1},\lambda_{H_1},\mu_T,\lambda_T)$.
Under the closure condition (same $\mu^2/\lambda$), the $\le_{\mathrm{lr}}$ results of Table~\ref{tab:ordering_lr} guarantee that $\mu_T \le \mu_{H_1}$ is a sufficient condition.
The formulation does not use this comparison, because its damage is not an inverse Gaussian variable in general (Section~\ref{sec:stoch_orders}).
It tracks the mean, imposes the CVaR cap of Theorem~\ref{thm:ig_cap}, and loops on the mean and, under $\text{ARD}_1$, on the anchor of the mean.
Facts (M1)--(M3) of Section~\ref{sec:encoding} hold: (M1) the repair factor $1-\rho$ lies in $[0, 1]$, and the mean increments are non-negative; (M2) the starting values satisfy $\theta_0 \le \xi_0$ as for the Gaussian model, and the true mean equals the day map applied to the true statistics of the day before; (M3) if a value $\mu \le \bar\mu$ is at least the true mean of day~$k$, then $\Pr(D_k > \tau) \le \varepsilon$, by Proposition~\ref{prop:ig_order} with this~$\mu$ and Theorem~\ref{thm:ig_cap}.
Each day of each repetition has its own set of terms; the argument holds for each of them, because it needs neither independence nor a fixed set of terms.

\paragraph{Rainflow.}
The tracked statistics and the rows are those of Section~\ref{sec:cons_surrogates}.
Facts (M1)--(M3) of Section~\ref{sec:encoding} hold.
(M1) The repair factors $1-\rho$ (for the mean and the cumulant generating function) and $(1-\rho)^2$ (for the variance and the Hoeffding budget) lie in $[0, 1]$, and the increments of the four statistics are non-negative; the cumulant generating function of an increment is at least $s$ times its mean, which is non-negative, by Jensen's inequality and A0.
(M2) The starting values satisfy $\theta_0 \le \xi_0$ as for the Gaussian model; on every later day the true statistics are at most the day maps applied to the true statistics of the day before, with equality for the mean, the variance and the Hoeffding budget, and an inequality for the cumulant generating function (Proposition~\ref{prop:rf_maps}).
(M3) If values that satisfy the row of the bound are at least the true statistics of day~$k$, then $\Pr(D_k > \tau) \le \varepsilon$ (the derivations of Section~\ref{sec:cons_surrogates}).
The loop constraint covers the mean, the second statistic of the bound and, under $\text{ARD}_1$, their anchors.

\paragraph{Compound Poisson.}
When $D \sim \mathrm{CP}(\nu,F_Y)$, the survival function does not admit a universal closed form.
Conditioning on the number of shocks gives the series
\begin{equation}\label{eq:surv_cp}
  \Pr(D > \tau) \;=\; \sum_{n=0}^{\infty} \frac{e^{-\nu}\,\nu^n}{n!}\;\Pr\!\Bigl(\sum_{k=1}^{n} Y_k > \tau\Bigr),
\end{equation}
which can be evaluated numerically (it converges rapidly) or bounded.

\paragraph{Upper bounds when no closed form is available.}
When the survival function $\Pr(D > \tau)$ cannot be computed in closed form---as for the compound Poisson process, or when the closure condition fails and the marginal of the sum is not within a standard family---the \emph{reliability} constraint can still be enforced through conservative upper bounds on the failure probability: Markov ($\Pr(D > \tau) \le \mathbb{E}[D]/\tau$), Cantelli~\eqref{eq:cantelli_surrogate}, and Chernoff~\eqref{eq:chernoff_surrogate}.
Section~\ref{sec:cons_surrogates} develops these bounds in per-mission form, together with the Bernstein and Hoeffding refinements, their relative tightness in the small-$\varepsilon$ regime, and their MILP encodings; for the compound Poisson process the Chernoff bound is particularly convenient, since $M_D(s) = \exp\bigl(\nu\,(M_Y(s)-1)\bigr)$ is available in closed form whenever $M_Y$ is.
The \emph{loop} constraint, by contrast, must not be translated into the same bounds: imposing $U^{(1)} \le U^{(0)}$ on a scalar bound~$U$ does not make the schedule repeatable, and the loop is instead imposed on the tracked statistics of the bound, for example the componentwise moment condition~\eqref{eq:loop_moments}, whose rows are linear in the tracked statistics (Section~\ref{sec:loop_routeI}).

\subsection{Maintenance strategies}
\label{sec:maint_strategies}

Maintenance actions on a repairable component range between two extremes: \emph{minimal repair}, which restores the component to the state it was in just before failure (``as bad as old'', ABAO), and \emph{perfect repair} or \emph{replacement}, which resets it to a like-new condition (``as good as new'', AGAN).
Real-world interventions almost always fall between these two extremes; they are collectively referred to as \emph{imperfect maintenance}~\cite{PhamWang1996}.
This section surveys the principal modelling families.

\subsubsection{Taxonomy of repair actions}

Table~\ref{tab:repair_actions} classifies five canonical maintenance actions by their effect on the component state.

\begin{table}[htbp]
\centering
\caption{Taxonomy of maintenance actions, from least to most restorative.}
\label{tab:repair_actions}
\small
\begin{tabular}{@{}p{3.0cm}p{4.5cm}p{5.8cm}@{}}
\toprule
\textbf{Action} & \textbf{Effect on damage $D$} & \textbf{Model} \\
\midrule
No intervention
  & $D^+ = D^-$
  & Damage unchanged \\[3pt]
Minimal repair (ABAO)
  & $D^+ = D^-$
  & Failure intensity restored to pre-failure level; damage history preserved~\cite{BrownProschan1983} \\[3pt]
Imperfect repair
  & $D^- > D^+ > D^{\mathrm{new}}$
  & Partial reduction of damage or virtual age (see below) \\[3pt]
Perfect repair (AGAN)
  & $D^+ = D^{\mathrm{new}}$
  & Component returned to like-new state \\[3pt]
Replacement
  & $D^+ = D^{\mathrm{new}}$
  & Physical component substituted; equivalent to perfect repair \\
\bottomrule
\end{tabular}
\end{table}

\noindent Here $D^-$ and $D^+$ denote the damage immediately before and after maintenance, and $D^{\mathrm{new}}$ is the damage of a new component.

\subsubsection{Virtual age models}

Kijima~\cite{Kijima1989} formalised imperfect repair through the concept of \emph{virtual age}.
After the $n$-th repair, the component behaves as if it had age $V_n$ (rather than its calendar age), and the failure intensity at calendar time~$t$ is $\lambda(t - T_n + V_n)$, where $T_n$ is the time of the $n$-th repair and $\lambda(\cdot)$ is the baseline intensity.
Two variants are standard:
\begin{itemize}
  \item \textbf{Kijima~I.} The repair acts only on the damage accumulated since the previous repair:
    \[
      V_n \;=\; V_{n-1} + (1-\rho)\,(T_n - T_{n-1}),
    \]
    where $\rho \in [0,1]$ is the repair efficiency.
    Only the most recent inter-repair period is partially rejuvenated (``memory-one'').
  \item \textbf{Kijima~II.} The repair acts on the entire accumulated virtual age:
    \[
      V_n \;=\; (1-\rho)\,(V_{n-1} + T_n - T_{n-1}).
    \]
    The full history is rejuvenated (``infinite memory'').
\end{itemize}
In both cases $\rho = 0$ gives minimal repair (ABAO) and $\rho = 1$ gives perfect repair (AGAN).

\subsubsection{Reduction-based models}

Doyen and Gaudoin~\cite{DoyenGaudoin2004} classify imperfect repair models by the \emph{quantity} that the repair reduces: the virtual age (classes ARA, arithmetic reduction of age) or the failure intensity (classes ARI, arithmetic reduction of intensity).
Each class has a memory: the largest number of previous failure times that enter the failure intensity~\cite[p.~8]{DoyenGaudoin2004}.
The memory-one variant carries the subscript~1, and the infinite-memory variant the subscript~$\infty$.
No class of Doyen and Gaudoin acts on a degradation level.
Mercier and Castro~\cite{MercierCastro2019} define a reduction that acts on the degradation level itself.

\paragraph{Arithmetic Reduction of Age (ARA).}
The repair removes a fraction~$\rho$ of the virtual age~$V$:
\begin{align}
  \text{ARA}_1: &\quad V^+ = V^- - \rho\,(V^- - V^{\mathrm{last}}), \label{eq:ara1} \\
  \text{ARA}_\infty: &\quad V^+ = (1-\rho)\,V^-, \label{eq:arainf}
\end{align}
where $V^{\mathrm{last}}$ is the virtual age just after the previous repair (\cite[eq.~(22), p.~12]{DoyenGaudoin2004} for $\text{ARA}_1$ and \cite[eq.~(20), p.~11]{DoyenGaudoin2004} for $\text{ARA}_\infty$).
$\text{ARA}_1$ reduces only the age accumulated since the previous repair, and $\text{ARA}_\infty$ reduces the whole virtual age.
With $V^- = V^{\mathrm{last}} + (T_n - T_{n-1})$, $\text{ARA}_1$ coincides with Kijima~I and $\text{ARA}_\infty$ with Kijima~II.

\paragraph{Arithmetic Reduction of Deterioration (ARD).}
The repair removes a fraction~$\rho$ of the degradation level~$D$:
\begin{align}
  \text{ARD}_1: &\quad D^+ = D^- - \rho\,(D^- - D^{\mathrm{last}}), \label{eq:ard1} \\
  \text{ARD}_\infty: &\quad D^+ = (1-\rho)\,D^-. \label{eq:ardinf}
\end{align}
$\text{ARD}_1$ is the ``arithmetic reduction of deterioration of order~1'' of Mercier and Castro~\cite[\S3.1, eq.~(4)]{MercierCastro2019}.
It removes the fraction~$\rho$ of the degradation accumulated since the last maintenance.
It acts on the degradation path, and the age is unchanged.
The reference level $D^{\mathrm{last}}$ is the level just after the last intervention, repair or replacement; before the first intervention, it is the initial level.
Neither source names the law $\text{ARD}_\infty$, and its name is an analogy.
$\text{ARD}_\infty$ applies the rule of $\text{ARA}_\infty$ to the whole degradation level, as $\text{ARD}_1$ applies the rule of $\text{ARA}_1$ to the degradation accumulated since the last intervention.

\paragraph{Arithmetic Reduction of Intensity (ARI).}
Instead of reducing the virtual age, the repair directly reduces the failure intensity:
\begin{align}
  \text{ARI}_1: &\quad \lambda^+ = \lambda^- - \rho\,(\lambda^- - \lambda^{\mathrm{last}}), \label{eq:ari1} \\
  \text{ARI}_\infty: &\quad \lambda^+ = (1-\rho)\,\lambda^-. \label{eq:ariinf}
\end{align}
ARI models are natural when maintenance targets the failure rate (e.g.\ condition-based interventions) rather than the underlying degradation state.
They are in general not equivalent to ARA/ARD models: a reduction in intensity does not correspond to a unique reduction in virtual age unless the baseline intensity is strictly monotone~\cite{DoyenGaudoin2004}.

\subsubsection{Connection to the fleet management formulation}
\label{sec:repair_connection}

The abstract repair operator $\mathcal{R}(D^-,\,\rho)$ introduced in Section~\ref{sec:overview} encompasses ARD and replacement models.
In particular:
\begin{itemize}
  \item $\text{ARD}_\infty$: $\mathcal{R}(D^-,\rho) = (1-\rho)\,D^-$, the default repair law of the formulation;
  \item $\text{ARD}_1$: $\mathcal{R}(D^-,\rho) = D^- - \rho\,(D^- - D^{\mathrm{last}})$, requiring the additional state variable $D^{\mathrm{last}}$;
  \item Replacement / AGAN: $\mathcal{R}(D^-,\rho) = D^{\mathrm{new}}$ (captured by the separate replacement variable $x^r_{i\ell k}$ in the formulation).
\end{itemize}
The formulation offers both repair laws for every degradation model, chosen per component.
The choice of repair model affects the degradation dynamics~\eqref{eq:dyn_full}, the definition of the repair cost variable~$z_{i\ell k}$, and the tractability of the resulting optimisation problem.
ARI models are not directly compatible with the damage-tracking formulation of Section~\ref{sec:overview} and would require reformulation in terms of the failure intensity rather than the damage state.
No law of the formulation acts on a virtual age.

\paragraph{Anchors.}
Under $\text{ARD}_1$ the formulation tracks, for each statistic of the damage, an \emph{anchor}: the value of the statistic just after the last intervention, repair or replacement.
On an intervention day the anchor takes the new value of its statistic; on any other day it keeps its value.
Example: mean $0.6$ and anchor $0.2$ before a repair with $\rho = 0.5$.
The repair gives the mean $0.6 - 0.5\,(0.6 - 0.2) = 0.4$ and the anchor $0.4$.
On the next day the gap between the mean and the anchor is $0$, and it grows again with each mission.
Before the first intervention of the horizon, the anchor is the value just after the last intervention before the horizon.
When that value is unknown, the formulation takes the initial value of the statistic.
Damage only grows between interventions, so the true value lies between the value of a new part and the initial value.
A larger anchor gives a weaker repair, so the initial value never overstates what a repair achieves.
Example: initial mean $0.5$, $\rho = 0.5$, and a repair on day~1 before any mission.
With the anchor $0$ the mean drops to $0.25$; with the anchor $0.5$ it stays at $0.5$, and the repair removes only damage accumulated within the horizon.

\paragraph{Repair maps of the tracked statistics.}
A model tracks moments or other summaries of the damage~$D$, not $D$ itself.
A repair maps a tracked statistic~$\theta$ with anchor~$g$ to $a\,\theta + (1-a)\,g$, with a factor $a \in [0,1]$ and $g = 0$ under $\text{ARD}_\infty$.
For the mean, $a = 1-\rho$, because expectation is linear: Eq.~\eqref{eq:ard1} gives $\mathbb{E}[D^+] = (1-\rho)\,\mathbb{E}[D^-] + \rho\,\mathbb{E}[D^{\mathrm{last}}]$.
For a variance and for the Hoeffding budget, $a = (1-\rho)^2$.
Under $\text{ARD}_1$ the reason is that $D^+ = G + (1-\rho)\Delta$, where $G$ is the damage just after the last intervention and $\Delta$ the damage added since, independent of~$G$; so $\operatorname{Var}[D^+] = \operatorname{Var}[G] + (1-\rho)^2\operatorname{Var}[\Delta]$.
For the cumulant generating function of the Chernoff bound, $a = 1-\rho$ gives an upper bound (Proposition~\ref{prop:rf_maps}).
The Wiener variance is the exception: a repair leaves it unchanged ($a = 1$), and it has no anchor (Section~\ref{sec:loop_per_model}).
These maps are non-decreasing in the statistic and in its anchor, as fact (M1) and Lemma~\ref{lem:monotone} require.

\paragraph{Laws per model.}
Every model of the formulation accepts both laws.
Table~\ref{tab:laws_per_model} gives the reason for each model.
Under $\text{ARD}_1$ with $\rho < 1$, the loop constraint requires a replacement in the operating phase of a component that gains damage there (Proposition~\ref{prop:ard1_replacement}).

\begin{table}[htbp]
\centering
\caption{Repair laws accepted by each degradation model, with the reason.}
\label{tab:laws_per_model}
\small
\begin{tabular}{@{}lllp{7.4cm}@{}}
\toprule
\textbf{Model} & $\text{ARD}_\infty$ & $\text{ARD}_1$ & \textbf{Reason or note} \\
\midrule
Gaussian & yes & yes & both laws keep the state Gaussian, with exact moments (Proposition~\ref{prop:gauss_exact}) \\
Inverse Gaussian & yes & yes & the cap holds for any repair weights and any dependence between the terms (Proposition~\ref{prop:ig_order}) \\
Wiener & yes & yes & on the mean only; the variance changes only by replacement \\
Gamma & yes & yes & the sharp cap holds for every weighted sum of Gamma terms (Lemma~\ref{lem:gamma_class}) \\
Rainflow, every bound & yes & yes & Chernoff under $\text{ARD}_1$ uses a conservative recursion for the cumulant generating function (Proposition~\ref{prop:rf_maps}) \\
\bottomrule
\end{tabular}
\end{table}

\paragraph{Maintenance under the Palmgren--Miner damage model.}
The remaining-life framework of Section~\ref{sec:discrete_deg} represents damage as the normalised fraction $D_{\mathrm{PM}} = \sum_i n_i/N_{f,i} \in [0,1]$, with failure at $D_{\mathrm{PM}} \ge 1$.
Because $D_{\mathrm{PM}}$ is additive over loading blocks, it fits the same damage-tracking structure as the stochastic processes of Table~\ref{tab:deg_models}: each mission day contributes an increment $\Delta D_{\mathrm{PM}} = n_j / N_{f,j}$ (deterministic if the S--N parameters are fixed, or random if scatter is modelled~\cite{MunizCalvente2022}).
The repair operator~$\mathcal{R}$ and the replacement action therefore apply directly:
\begin{itemize}
  \item \emph{Replacement} resets the fatigue fraction to zero: $D_{\mathrm{PM}}^+ = 0$ (equivalently $D^{\mathrm{new}} = 0$).
    This corresponds to substituting the physical component with a new one whose cycle count is zero.
  \item $\text{ARD}_\infty$: $D_{\mathrm{PM}}^+ = (1-\rho)\,D_{\mathrm{PM}}^-$.
    Physically, this models an intervention that partially restores fatigue life (e.g.\ surface treatment, shot peening, localised crack repair) without resetting the full loading history.
  \item $\text{ARD}_1$: $D_{\mathrm{PM}}^+ = D_{\mathrm{PM}}^- - \rho\,(D_{\mathrm{PM}}^- - D_{\mathrm{PM}}^{\mathrm{last}})$.
    Only the damage accumulated since the last intervention (repair or replacement) is partially reversed.
\end{itemize}
A notable difference from the stochastic process models is that $D_{\mathrm{PM}}$ is bounded above by construction (failure at~1), so the threshold is fixed at $\tau = 1$ and the reliability constraint becomes $\Pr(D_{\mathrm{PM}} \ge 1) \le \varepsilon$.
In the deterministic variant of Miner's rule this is a hard feasibility check; in the stochastic extension it becomes a chance constraint analogous to~\eqref{eq:abs_rel}.

Closure under summation is trivially satisfied: $D_{\mathrm{PM}}$ is a sum of real-valued increments, so the marginal at any step is the convolution of the individual increments regardless of their distributional family.
In the deterministic variant of Miner's rule, comparing two schedules reduces to comparing two real numbers and no stochastic ordering is needed.
In the stochastic extension (scatter in $N_{f,i}$), the parametric ordering conditions of Section~\ref{sec:stoch_orders} apply only if the distributional assumption on the scatter places the increments $n_i/N_{f,i}$ in one of the families of Table~\ref{tab:deg_models}; otherwise the general ordering definitions still hold but must be verified case by case.

\paragraph{Remaining items for this section.}
\begin{itemize}
    \item[$\checkmark$] Fatigue models (either damage accumulation or remaining life)
    \item[$\square$] Degradation-aware fleet management
    \item[$\checkmark$] Maintenance strategies
    \item[$\checkmark$] Stochastic orderings
    \item[$\square$] Similar applications
    \item[$\square$] Novelties of the project and its placement in the literature
\end{itemize}

\section{Focus on electric vehicles}

A battery-electric road vehicle---a transit bus or a passenger car---is a complex repairable system whose parts degrade at very different rates and through very different physical mechanisms.
Not every part deserves an explicit place in the fleet-management model of Section~\ref{sec:overview}.
We retain a component if it satisfies two criteria:
\begin{enumerate}[label=(C\arabic*)]
  \item \textbf{Usage-driven degradation.}
    Its condition deteriorates measurably as a function of \emph{how} the vehicle is used, so that the mission assigned on a given day produces a quantifiable damage increment $\Delta D^{(j)}_{i\ell k}$ (Section~\ref{sec:overview}).
    A part that ages only with calendar time, independently of the mission served, carries no scheduling information and need not enter the per-mission degradation state.
  \item \textbf{Operational significance.}
    Its failure or preventive renewal is a material driver of at least one of \emph{safety}, \emph{availability}, or \emph{life-cycle maintenance cost}---the three quantities the schedule trades off.
\end{enumerate}
We deliberately do \emph{not} require that the component admit a degradation model expressible within the families of Section~\ref{sec:stoch_orders}: a component can be worth considering even when the abstract model of Section~\ref{sec:overview} does not yet capture it.
Such cases are listed here together with the modelling extension they would demand, rather than being silently dropped.

The components meeting these criteria fall into two broad groups.
The \emph{electrochemical and electrical powertrain}---traction battery, electric motor, power electronics, charging interface---degrades through electrochemical ageing and thermal and electrical cycling, driven by energy throughput, charge and discharge rates, and temperature.
The \emph{mechanical running gear}---tires, friction brakes, wheel bearings, suspension, reduction gear---degrades through wear and fatigue, driven by distance travelled, vehicle mass, road quality, and the acceleration/braking profile.
Table~\ref{tab:ev_components} lists the principal candidates in each group, the dominant degradation mechanism, the operational quantity that drives it, and the reason the component is operationally significant.

Two domain-specific facts shape the relative importance of these components.
First, the traction battery dominates: a recent predictive-maintenance study of electric buses centres entirely on the battery pack as the primary maintenance-critical component, models its end-of-life through the accumulated charge throughput, and shows that optimising the driving profile alone can cut the modelled capacity loss by 25\%~\cite[Secs.~2,~4.2]{Ibrahim2025}---direct evidence that battery damage is steerable by exactly the kind of mission-assignment decisions our framework optimises.
Second, regenerative braking absorbs most of the service braking, with estimated reductions of roughly 60--90\% in friction-brake wear emissions~\cite[Sec.~6.1]{Fussell2022}; the relative wear burden consequently shifts to the tires, whose wear is aggravated by the higher mass of electric vehicles---to the point that, on motorways, the added tire-wear particle emissions are not offset by the regenerative-braking savings~\cite[Sec.~6.1]{Fussell2022}.

\begin{table}[htbp]
\centering
\caption{Candidate electric-vehicle components for degradation-aware modelling.
Relevance flags: \textbf{S} safety-critical, \textbf{A} availability driver (its failure or renewal takes the vehicle out of service), \textbf{C} maintenance-cost driver.
The ``mission driver'' column names the operational quantity that determines the per-mission damage increment $\Delta D^{(j)}_{i\ell k}$.}
\label{tab:ev_components}
\small
\begin{tabular}{@{}p{3.0cm}p{4.6cm}p{3.5cm}c@{}}
\toprule
\textbf{Component} & \textbf{Dominant degradation mechanism} & \textbf{Mission driver} & \textbf{Rel.} \\
\midrule
\multicolumn{4}{@{}l}{\textit{Electrochemical / electrical powertrain}} \\
Traction battery pack
  & Capacity fade and impedance growth: SEI growth, lithium plating, loss of active material (cycle $+$ calendar ageing)
  & Energy throughput, DOD, C-rate, temperature
  & A, C \\[3pt]
Traction motor
  & Winding-insulation thermal ageing; bearing wear
  & Thermal/load cycles, current
  & A, C \\[3pt]
Power semiconductors (IGBT / SiC modules in the inverter and converters)
  & Thermal-cycling fatigue (bond-wire lift-off, solder cracking)
  & Power throughput, thermal cycles
  & A, C \\[3pt]
Converter capacitors (DC-link, filtering)
  & Electrolyte vaporisation / dielectric wear-out (ripple-current self-heating)
  & Ripple current, temperature, operating hours
  & A, C \\[3pt]
Control / power PCBs (gate drivers, BMS, converter boards)
  & Solder-joint and via thermo-mechanical fatigue; random infant/wear-out failures
  & Thermal cycles, power-on hours, vibration
  & A \\[3pt]
Charging interface (inlet, connector)
  & Contact wear and contact-resistance growth over mating cycles
  & Number of charging events
  & A \\
\midrule
\multicolumn{4}{@{}l}{\textit{Mechanical: running gear and chassis}} \\
Tires
  & Tread abrasion (friction work at the contact patch)
  & Distance, vehicle mass, driving style, road surface
  & S, C \\[3pt]
Brake friction elements (pads, discs)
  & Abrasive and thermal wear (mitigated by regenerative braking)
  & Number and energy of friction-braking events
  & S, C \\[3pt]
Wheel bearings
  & Rolling-contact fatigue spalling
  & Distance, load
  & S, A \\[3pt]
Suspension (springs, dampers, bushings)
  & Fatigue; loss of damping
  & Load cycles, road quality
  & A \\[3pt]
Reduction gear / drive shafts
  & Wear, contact fatigue, lubricant ageing
  & Transmitted torque, distance
  & A \\[3pt]
Auxiliaries (HVAC, battery thermal management, doors)
  & Wear, cycle-driven failures
  & Operating hours, actuation cycles
  & A \\
\bottomrule
\end{tabular}
\end{table}

\paragraph{Power-electronic and electronic components.}
The traction chain contains four electrical sub-families beyond the battery, and they fit the framework to differing degrees.
The \emph{power semiconductors} of the inverter and the DC--DC and on-board-charger \emph{converters} are governed by junction-temperature-swing fatigue and map cleanly onto the remaining-life family (treated in detail below).
The \emph{converter capacitors}---DC-link and filtering---wear out by electrolyte vaporisation and dielectric ageing driven by ripple-current self-heating; this is a monotone, usage-driven accumulation (Arrhenius in temperature, power law in ripple current) and therefore also fits the framework, with the mission acting through the delivered power and the resulting ripple and thermal load.\footnote{The capacitor and PCB degradation-model forms in this paragraph are stated from standard power-electronics reliability theory; located primary sources have not yet been secured for them (see \texttt{inaccessible\_papers.md}).}
The \emph{control and power printed circuit boards}---gate-driver, battery-management, and converter boards---are the borderline case: their solder-joint and plated-through-via fatigue is genuinely usage-driven (thermal cycling, power-on hours, vibration) and so satisfies~(C1)--(C2), but a substantial part of their failures are sudden and effectively random rather than a smooth accumulation crossing a threshold.
The damage-accumulation dynamics of Section~\ref{sec:overview} do not fully capture this mixed wear-out-plus-random behaviour; representing it faithfully would require either a compound-Poisson shock term (Section~\ref{sec:stoch_orders}) or a failure-intensity formulation rather than the damage-state formulation of Equation~\eqref{eq:dyn_full}.
We list PCBs here, with this caveat, precisely so the limitation is recorded rather than hidden by the choice of model.

\paragraph{Modelling shortlist.}
For a first concrete instantiation of the abstract framework, the most informative components are those that are simultaneously safety- or cost-critical \emph{and} equipped with mature, monotone degradation models that fit the closure and ordering analysis of Section~\ref{sec:stoch_orders}: the \emph{traction battery pack}, the \emph{tires}, the \emph{brake friction elements}, the \emph{wheel bearings}, the \emph{traction-motor insulation}, and the \emph{inverter and converter power modules}.
Each maps onto a single per-component degradation variable $D_{i\ell k}$ whose mission-dependent increment $\Delta D^{(j)}_{i\ell k}$ is governed by the operational driver in the third column of Table~\ref{tab:ev_components}.
A long highway mission, for instance, produces large tire-wear and deep-discharge battery increments, whereas a stop-intensive urban mission multiplies braking events and charge--discharge micro-cycles; this mission-dependence is exactly the coupling that the scheduling problem of Section~\ref{sec:overview} exploits.

\paragraph{Worked example: traction-battery capacity fade.}
The battery pack is the costliest component of an electric vehicle and the focus of its predictive maintenance~\cite[Secs.~2,~4.2]{Ibrahim2025}, so it serves as the reference instantiation of the abstract increment $\Delta D^{(j)}_{i\ell k}$.
Its degradation state is the relative capacity fade (one minus state-of-health), and the failure threshold~$\tau$ is the maximum admissible fade (the end-of-life criterion).
Mechanistically, capacity is lost through three side-reaction families---SEI film growth, lithium plating, and gas generation---whose contributions split into loss of lithium inventory and loss of active material~\cite[Secs.~2.1,~4.1, Eqs.~28--29]{Lv2025}.
At the fleet level this is abstracted by semi-empirical models that decompose the fade into a \emph{calendar} term and a \emph{cycling} term: the calendar term grows with an Arrhenius factor $a\exp(-E_a/RT)$ in temperature and a power law in time and state-of-charge, while the cycling term grows with the number of cycles modulated by temperature, C-rate, and depth-of-discharge stress factors~\cite[Eqs.~1--3]{Hosen2021}.
An even simpler engineering abstraction links end-of-life directly to the accumulated charge throughput~\cite[Sec.~4.2]{Ibrahim2025}.
For our framework the decisive structure is: the cycling fade is monotone and additive over days, and the daily increment is a function of the energy served on the mission---including the recharge that restores it, so that fast charging (high C-rate) makes the same mission costlier in damage---together with its depth-of-discharge and temperature profile.
This is precisely a mission-dependent accumulation process of the Gamma or inverse-Gaussian type of Section~\ref{sec:stoch_orders}; the non-stationary (power-law) time dependence is absorbed by the time-varying increment parameters that the formulation already allows ($\Delta D^{(j)}_{i\ell k}$ may depend on~$k$).
The calendar term, by contrast, accrues even on idle and maintenance days and is therefore \emph{not} captured by the dynamics of Section~\ref{sec:overview}, which leave the state unchanged on such days; over short horizons it can be neglected, while over long horizons the dynamics~\eqref{eq:dyn_full} must be extended with a small activity-independent drift added to every case, including idle.
Maintenance also takes a specific form: replacing a subset of battery modules removes a corresponding fraction of the pack-level damage---a physical realisation of the ARD repair operator $\mathcal{R}$ with $\rho$ the fraction of modules renewed---while full pack replacement is the AGAN action~$x^r$.

\subsection{Degradation models for the shortlisted components}
\label{sec:ev_deg_models}

The remaining shortlisted components degrade through the same two structural mechanisms identified for the battery---monotone usage-driven \emph{wear/accumulation} and \emph{cyclic-load fatigue life consumption}---and therefore map onto the same two model families of Section~\ref{sec:stoch_orders}.
Table~\ref{tab:ev_deg_models} summarises the governing model for each, the operational quantity that drives the per-mission increment, and the abstract family it instantiates.

\begin{table}[htbp]
\centering
\caption{Degradation models for the shortlisted electric-vehicle components and their mapping onto the abstract model families of Section~\ref{sec:stoch_orders}.
``Accumulation'' denotes a monotone damage-accumulation process (Gamma / inverse-Gaussian / compound-Poisson); ``remaining-life'' denotes a cycles-to-failure model accumulated by the Palmgren--Miner rule.}
\label{tab:ev_deg_models}
\small
\begin{tabular}{@{}p{2.4cm}p{4.9cm}p{3.0cm}p{2.6cm}@{}}
\toprule
\textbf{Component} & \textbf{Degradation model / law} & \textbf{Mission driver} & \textbf{Abstract family} \\
\midrule
Battery pack
  & Calendar $+$ cycle capacity fade: Arrhenius in $T$, power law in $t$/SOC, cycle-count term \cite[Eqs.~1--3]{Hosen2021}; mechanistic SEI / plating / LAM \cite[Sec.~2.1]{Lv2025}
  & Energy throughput, DOD, C-rate, temperature
  & Accumulation ($k$-dependent) \\[3pt]
Tires
  & Friction-work abrasion \cite[Eq.~7]{SchallamachTurner1960}; factors: tire build, road surface, vehicle operation \cite[Sec.~2.1.2]{Fussell2022}
  & Distance, mass, accel./braking, road surface
  & Accumulation \\[3pt]
Brake friction pad
  & Archard-type wear per braking event \cite[Eqs.~14,~22]{Sha2022}; regenerative-braking mitigation \cite[Sec.~6.1]{Fussell2022}
  & Number and energy of friction-braking events
  & Accumulation \\[3pt]
Wheel bearing
  & Subsurface-fatigue spall propagation, continuum damage \cite[Sec.~2.1.3]{Ohana2023}
  & Distance, speed, load
  & Accumulation / remaining-life \\[3pt]
Traction-motor insulation
  & Arrhenius--Dakin thermal-endurance life \cite{Dakin1948}; Montsinger rule \cite{Montsinger1930}; cumulative-damage extension \cite[Eqs.~1--2]{Zhang2025insul}
  & Thermal / load cycles, current
  & Remaining-life (Miner) \\[3pt]
Inverter power module
  & Coffin--Manson / Norris--Landzberg cycles-to-failure in $\Delta T_j$ \cite[Eqs.~3--5]{Alavi2024}; Bayerer/CIPS08 \cite[Eq.~1]{Bayerer2008}
  & Power throughput, thermal cycles
  & Remaining-life (Miner) \\
\bottomrule
\end{tabular}
\end{table}

\paragraph{Tires.}
Tread loss is governed by the frictional work expended at the tire--road contact patch: the classical theory of Schallamach and Turner assumes abrasion proportional to the frictionally dissipated energy~\cite[Eq.~7]{SchallamachTurner1960} and derives that, at small slip, wear grows approximately with the \emph{square} of the slip~\cite[Eq.~19]{SchallamachTurner1960}---so the wear rate rises steeply with the tangential forces generated by traction, braking, and cornering.
The principal factors are the tire characteristics (compound, construction), the road-surface characteristics, and the vehicle operation~\cite[Sec.~2.1.2]{Fussell2022}.
The higher kerb mass of electric vehicles raises these contact forces---enough that, on motorways, the added tire-wear particle emissions are not offset by the brake-wear savings of regenerative braking~\cite[Sec.~6.1]{Fussell2022}.
Tread depth decreases monotonically from its new value to the legal limit, which plays the role of the threshold~$\tau$; the per-mission increment depends on distance, load, road surface, and driving aggressiveness.
For bus fleets, retreading acts as a near-AGAN replacement of the wear state, while tire rotation merely redistributes wear across positions and is left outside the per-component model.

\paragraph{Brake friction elements.}
Friction-pad material loss is governed by the energy dissipated at the pad--disc interface during braking; a physics-based route applies Archard's law on a fractal, elastoplastic contact to obtain a per-cycle wear depth and hence a brake-cycle life $N = H_0/(S_0\,h_{\mathrm{wear}})$, with $H_0$ the admissible wear depth and $h_{\mathrm{wear}}$ the wear per braking~\cite[Eqs.~14,~22]{Sha2022}.
In an electric vehicle the regenerative brake absorbs most service braking, cutting friction-brake wear by an estimated 60--90\%~\cite[Sec.~6.1]{Fussell2022}; the residual increments concentrate in the events where regeneration is unavailable or insufficient---battery near full charge, very low temperature, emergency stops.
The cumulative pad wear remains a sum of non-negative, mission-dependent increments---one per friction-braking event---i.e.\ a monotone accumulation process whose rate now also depends on the energy-management strategy.

\paragraph{Wheel bearings.}
Bearing degradation proceeds in two stages: a long rolling-contact-fatigue \emph{initiation} phase followed by rapid \emph{spall propagation} once a surface defect forms.
Ohana et al.~\cite[Sec.~2.1.3]{Ohana2023} model the propagation phase with a continuum-damage-mechanics formulation in which a scalar damage variable accumulates per over-rolling cycle, $\mathrm{d}D/\mathrm{d}N = \bigl(\sigma_{\mathrm{eff}}/[\sigma_r(1-D)]\bigr)^m$, with growth driven by the tensile residual stresses generated at the spall edge by repeated rolling-element impacts; the model is intended as the basis for a physics-based prognostic tool that estimates remaining useful life once a critical spall size is reached.
Because the underlying driver is the number of load cycles---and hence running distance and speed---the bearing damage is a monotone accumulation process, well represented by a Gamma or inverse-Gaussian process (initiation) feeding a remaining-life threshold (propagation).

\paragraph{Traction-motor insulation.}
The winding insulation is the principal life-limiting element of the traction motor, and its dominant ageing mechanism is thermal: elevated winding temperature accelerates the chemical degradation of the dielectric.
Dakin recognised thermal ageing as a chemical reaction whose rate constant follows the Arrhenius exponential law $K = A\,e^{-B/T}$ in absolute temperature~\cite[Eq.~5]{Dakin1948}, giving the thermal-endurance life $L(T) = a\,\exp(b/T)$~\cite[Eq.~1]{Zhang2025insul}; the earlier empirical rule of Montsinger states that the ageing rate roughly doubles for each 8\,$^\circ$C temperature increase~\cite[Synopsis, p.~776]{Montsinger1930}---a halving of life per fixed rise that Dakin showed to be the constant-coefficient approximation of the Arrhenius law, accurate at one temperature only~\cite[Fig.~1 and discussion of Eq.~5]{Dakin1948}.
Under a mission-varying temperature profile, the fraction of life consumed in an interval $\mathrm{d}t$ at temperature $T_i(t)$ is $\mathrm{d}t/L[T_i(t)]$, and the total consumed fraction is the integral of this loss rate over the profile---a cumulative-damage (Palmgren--Miner) extension of the Arrhenius model~\cite[Eq.~2 and Fig.~10]{Zhang2025insul}.
The abstract description therefore coincides with the remaining-life model of Section~\ref{sec:discrete_deg}, with the mission driving the thermal load (current, cooling, ambient); notably, the same review also surveys Wiener- and Gamma-process formulations of insulation degradation~\cite[Sec.~5]{Zhang2025insul}, directly matching the stochastic families of Section~\ref{sec:stoch_orders}.

\paragraph{Inverter power modules.}
The IGBT or SiC modules of the traction inverter fail predominantly by thermo-mechanical fatigue---aluminium bond-wire lift-off and solder-layer cracking---driven by the coefficient-of-thermal-expansion mismatch under repeated junction-temperature swings $\Delta T_j$~\cite[Sec.~1]{Bayerer2008}.
The number of cycles to failure is classically captured by a Coffin--Manson power law $N_f = \alpha\,(\Delta T_j)^{-n}$~\cite[Eq.~3]{Alavi2024}, refined by adding an Arrhenius factor in the mean junction temperature $T_{jm}$, $N_f = A\,(\Delta T_j)^{\alpha}\exp\!\bigl(E_a/(k_b T_{jm})\bigr)$~\cite[Eq.~4]{Alavi2024}, and generalised by the multi-factor Bayerer (CIPS08) model
$N_f = K\,\Delta T_j^{\beta_1} e^{\beta_2/(T_{j}+273)}\, t_{\mathrm{on}}^{\beta_3}\, I^{\beta_4}\, V^{\beta_5}\, D^{\beta_6}$,
fitted by multiple regression over power-cycling tests spanning temperature swing, power-on time~$t_{\mathrm{on}}$, current per bond wire~$I$, blocking voltage~$V$, and bond-wire diameter~$D$~\cite[Eq.~1]{Bayerer2008}.
Being purely empirical, the model is valid only within the fitted test range and is not to be extrapolated~\cite[Sec.~3.1]{Bayerer2008}; its coefficients also require recalibration per device family---a recent re-fit on dedicated power-cycling tests shifted $(\beta_1,\beta_2,\beta_3)$ substantially and cut the predicted B10 inverter lifetime from 43 to 13~years~\cite[Abstract and Sec.~5]{Alavi2024}.
Under a realistic mission mix the thermal cycles have heterogeneous amplitudes, and the accumulated damage is summed by the Palmgren--Miner rule $Q = \sum_i n_i/N_{f,i}$, failure at $Q=1$---placing the power module squarely in the remaining-life family of Section~\ref{sec:discrete_deg}, with the mission determining the transmitted power and hence the temperature-swing spectrum.

\paragraph{Common structure.}
Every shortlisted component reduces to one of the two abstract descriptions of Section~\ref{sec:stoch_orders}: a \emph{monotone accumulation process} (battery cycling fade, tires, brakes, bearings) or a \emph{remaining-life model accumulated by the Palmgren--Miner rule} (insulation, power modules).
In both, the per-mission damage increment $\Delta D^{(j)}_{i\ell k}$ is set by an operational driver that the schedule controls---energy throughput and charging rate, distance and load, braking profile, or thermal cycling---which is precisely the structure the abstract formulation of Section~\ref{sec:overview} assumes.
This confirms that the electric-vehicle fleet is a faithful concrete instance of the general problem, and fixes which degradation family parameterises each component when the model is specialised in Section~\ref{sec:evs}.

\paragraph{Excluded parts.}
Components whose deterioration is essentially independent of the mission served---body corrosion, interior fittings, static low-voltage electronics, the 12-volt auxiliary battery---satisfy criterion~(C2) but not~(C1).
They are best handled by time-based overhaul outside the optimisation and are therefore omitted from the per-mission degradation state, although their fixed maintenance windows can be imposed as exogenous availability constraints.
The calendar-ageing component of the traction battery is the one borderline case and is handled by the drift extension discussed in the worked example above.

\paragraph{Remaining items for this section.}
\begin{itemize}
    \item[$\checkmark$] Relevant components (mechanical and electrical) to model an electric vehicle
    \item[$\checkmark$] Typical degradation model for each component, and the operational assumptions under which it holds
\end{itemize}

\section{Optimization problem for a fleet of electric vehicles}
\label{sec:evs}

Table~\ref{tab:notation_evs} lists the additional symbols introduced when specialising the abstract formulation of Section~\ref{sec:overview} to a concrete fleet of electric vehicles.
All symbols from Table~\ref{tab:notation} remain in force; the entries below either instantiate abstract quantities or add new ones.

\begin{table}[htbp]
\centering
\caption{Additional notation for the electric-vehicle fleet specialisation (extends Table~\ref{tab:notation}).}
\label{tab:notation_evs}
\small
\begin{tabular}{@{}lp{10.5cm}@{}}
\toprule
\textbf{Symbol} & \textbf{Description} \\
\midrule
\multicolumn{2}{@{}l}{\textit{Cost parameters}} \\
$C_M$                        & Cost of a depot day, paid once per member and depot day (scalar) \\
$C_D$                        & Damage regularisation coefficient (scalar) \\
$C_{R,i\ell}$                & Cost of imperfect repair per unit of expected degradation removed \\
$C_{\mathrm{rep},i\ell}$     & Cost of full replacement of component $\ell$ of member $i$ \\
\bottomrule
\end{tabular}
\end{table}

\subsection{Cost function}

With the cost parameters of Table~\ref{tab:notation_evs}, the abstract objective of Section~\ref{sec:overview} takes the form:
\begin{align}\label{eq:objective_evs}
\min \quad
  & C_M \sum_{k=1}^{T} \sum_{i=1}^{F} x_{i0k}
  \;+\; \sum_{k=1}^{T} \sum_{i=1}^{F} \sum_{\ell=1}^{L} C_{R,i\ell}\, z_{i\ell k}
  \;+\; \sum_{k=1}^{T} \sum_{i=1}^{F} \sum_{\ell=1}^{L} C_{\mathrm{rep},i\ell}\, x^r_{i\ell k} \notag \\
  & +\; C_D\, u.
\end{align}
Each term is a weighted sum over the planning horizon:
the first counts depot days, the second scales with the degradation removed by each repair, the third counts replacements, and the last penalises the worst-case aggregate expected damage via~$u$ (cf.~\eqref{eq:u_def}).

\begin{itemize}
    \item Specialize the previous optimization problem to the case of electric vehicles
    \item Look for what are can be exactly modeled and what no
    \item Provide the MILP approximation of such problem
\end{itemize}

\section{Implementation details}

The formulation of this document is implemented in the Python package \texttt{fleet\_management}.
The package builds the mixed-integer program of Section~\ref{sec:overview} for a given fleet and solves it with the commercial solver Gurobi.
It covers the five degradation models of Table~\ref{tab:loop_per_model}, both repair laws, the five Route-I bounds and the three forms of the reliability rows.
It also offers tools around the solve: sweeps over the horizon lengths, warm starts, plots of a schedule, and a check of a result against its input.

The software specification \texttt{spec/spec.tex} is the reference for the implementation.
It states the formulation in full, with every row, the input and output files, and the tools, and the code follows it.
Where this document and the specification differ in a detail, the specification holds.

\begin{itemize}
	\item Provide the details to run an actual simulation of the system
\end{itemize}

The subsections below give the reasons behind the implementation choices of the specification.

\subsection{Big-M and indicator forms of the one-sided rows}
\label{sec:impl_encoding}

The specification writes each one-sided row of Section~\ref{sec:encoding} in one of two forms, chosen per component.
In the \emph{big-M form}, each rule of a tracked statistic~$s$ holds unless a binary switches it off through a constant~$M_s$: the carry row $s_k \ge s_{k-1} + \delta s_k - M_s\,(x^m_k + x^r_k)$, the repair row $s_k \ge a\,s_{k-1} + (1-a)\,g_{k-1} - M_s\,(1 - x^m_k)$, the replacement rows $|s_k - s^{\mathrm{new}}| \le M_s\,(1 - x^r_k)$, and the anchor rows $|g_k - s_k| \le M_s\,(1 - x^m_k - x^r_k)$ and $|g_k - g_{k-1}| \le M_s\,(x^m_k + x^r_k)$.
In the \emph{indicator form}, the same inequalities are Gurobi indicator constraints, switched on by the binary of the day's case.
At binary points the two forms have the same feasible set.

The constant $M_s$ is the \emph{natural range} of~$s$: an upper bound of~$s$ at every point that satisfies the reliability row of the component.
It is $\tau$ for a mean and for its anchor, $V_{\max} = \tau^2/\kappa_G^2$ for a Gaussian or Wiener variance (Remark~\ref{rem:vmax}), $\varepsilon\tau^2/(1-\varepsilon)$ for a Cantelli variance, $\tau(\tau - 2c)/(2\kappa)$ for a Bernstein variance, $2\tau^2/\kappa$ for the Hoeffding budget and $s\tau - \kappa$ for the Chernoff statistic; each is the largest value that the exact row of the bound admits at $\mu \ge 0$.
An anchor has the range of its statistic.

\paragraph{Validity of the big-M constants.}
Every tracked statistic and every anchor lies in $[0, M_s]$.
The initial values, the starting anchors and the replacement values lie in $[0, M_s]$ too, by the input checks (Section~\ref{sec:check_reasons}).
An intervention day is a depot day, so $\delta s_k = 0$ on it.
A relaxed row then cuts no point: the carry row with an intervention reads $s_k \ge s_{k-1} - M_s$, whose right-hand side is at most~$0$; the repair row without a repair has a right-hand side $a\,s_{k-1} + (1-a)\,g_{k-1} - M_s \le 0$, because the repair map is a convex combination of values in $[0, M_s]$; the replacement rows without a replacement allow $s_k \in [s^{\mathrm{new}} - M_s,\, s^{\mathrm{new}} + M_s] \supseteq [0, M_s]$; the relaxed anchor rows allow every difference in $[-M_s, M_s]$.

\paragraph{Choice of the default form.}
The big-M form is the default.
Gurobi accepts a binary as integral within its integrality tolerance \texttt{IntFeasTol} (default $10^{-5}$).
A big-M row can therefore leak by at most $10^{-5} M_s$.
Because $M_s$ is the statistic's own range, this is a relative error of $10^{-5}$.
Indicator constraints have no such leak, and Gurobi usually needs more time to solve them.
Which form solves this model faster has not been measured.

\subsection{Reasons for the input checks}
\label{sec:check_reasons}

The specification validates the input before it builds the model, and a violation raises an error that names the key and the component.
A check of one model runs only at the components of that model.
Each item below gives a rule and, after the dash, its reason.

\paragraph{Keys, shapes and entries.}
\begin{itemize}[nosep]
  \item Every key is known, and every required key is present. --- A misspelled or missing key must not change the problem without a message.
  \item Each value has one of the shapes of its key. --- Silent broadcasting by the array library could read a value for another component.
  \item At a component where a key applies, every entry is a finite number, or a string for a selector. --- A missing value at a component must not take a default silently.
  \item An unknown value of a selector is an error at every component, also where the reliability row is linear and the form is ignored. --- Reading an unknown value as \texttt{exact} would turn a linear model into a nonconvex one without a message.
\end{itemize}

\paragraph{Sizes and selectors.}
\begin{itemize}[nosep]
  \item $F > M$. --- Every mission needs a member on every day; with $F = M$ no member could be in the depot, and no repair could take place.
  \item $H_1, H_2 \ge 1$. --- Each phase has at least one day.
  \item With a positive $C_D$ and the quadratic penalty, every quadratic reliability row is in its exact form. --- The quadratic penalty is a quadratic row, and the linear forms exist to give a linear model (Section~\ref{sec:cost_function}).
\end{itemize}

\paragraph{Reliability, repair and initial state.}
\begin{itemize}[nosep]
  \item $\tau > 0$. --- The threshold is a positive degradation level.
  \item $\varepsilon \in (0, 1)$, and $\varepsilon < 0.5$ at Gaussian and Wiener components. --- A probability lies in $(0,1)$; the squared normal quantile of the Gaussian and Wiener rows is positive only below~$0.5$ (Remark~\ref{rem:vmax}).
  \item $\rho \in (0, 1]$. --- A repair removes a positive fraction of the level, at most all of it.
  \item $0 \le \mu_0 < \tau$. --- A component at its threshold has failed already.
  \item Initial and replacement variances, supports and budgets are non-negative. --- They describe a variance, a width or a squared width.
  \item The relative tangent point lies in $[0, 1)$, and the number of \texttt{pwl} segments is at least~$1$. --- At the point~$\tau$ the tangent row reduces to $\chi \le 0$, and a value outside $[0,1]$ has no tangent point in $[0, \tau]$.
  \item Each starting anchor lies between $0$ and the initial value of its statistic. --- Damage only grows between interventions, so the level after the last intervention cannot exceed the initial level.
  \item With the shortcut for a new part, every replacement value is~$0$. --- The shortcut means a new part, and a non-zero value would be discarded without a message.
\end{itemize}

\paragraph{Natural ranges and replacement states.}
\begin{itemize}[nosep]
  \item Every initial value, starting anchor and replacement value lies in the natural range of its statistic (Section~\ref{sec:impl_encoding}). --- The big-M rows take the natural range as their constant, so a value outside it violates a row at once.
  \item When replacements are allowed, the replacement state satisfies every row that holds on a replacement day, in the form built at the component (for \texttt{pwl}, the cap of the segment that contains the replacement mean), its side row, and the optional variance cap. --- The rows hold on the day of a replacement; otherwise a replacement is never feasible (Section~\ref{sec:damage_dynamics}).
\end{itemize}

\paragraph{Increments.}
\begin{itemize}[nosep]
  \item $0 \le \mu^{\mathrm{inc}} < \tau$ for every mission and day. --- A mission with mean $\tau$ or more fails the component in one day; the mean~$0$ is allowed, so that a mission that does not wear a component can be written as~$0$.
  \item Where the mean of an increment is~$0$, its variance, support and cumulant generating function are~$0$. --- A non-negative increment with mean~$0$ is $0$ almost surely, so it has no variance, no range and a zero cumulant generating function.
  \item Variances, supports and cumulant generating functions of increments are non-negative, and a transitory profile comes with its operating profile. --- A variance and a width are not negative, and the lookup of the increment laws needs the operating laws after the transitory phase (Section~\ref{sec:damage_dynamics}).
\end{itemize}

\paragraph{Model-specific checks.}
\begin{itemize}[nosep]
  \item Gaussian: $\mu \ge 3\sqrt{v}$ for the initial damage, for every increment and, when replacements are allowed, for the replacement draw. --- A Gaussian term is then negative with probability at most about $0.13\,\%$, so the damage stays non-negative in practice.
  \item Inverse Gaussian: $\eta > 0$. --- The dispersion of the inverse Gaussian law must be positive.
  \item Wiener: $\sigma > 0$ and $v_0 \ge 0$. --- The diffusion constant is positive; a known initial damage has variance~$0$ and is a valid input.
  \item Gamma: every given rate is finite and positive. --- The rate of a Gamma law is positive, and the smallest rate sets the cap (Section~\ref{sec:stoch_orders}); an omitted rate of the initial or replacement damage means a known damage, which has no Gamma law and no rate.
  \item Optional variance cap: positive, and at least the initial variance and, when replacements are allowed, the replacement variance. --- A cap below either of them makes the problem infeasible.
  \item Bernstein: every variance is at most a quarter of the squared width of its range, for increments, the initial damage and the replacement draw. --- Popoviciu's inequality (Section~\ref{sec:cons_surrogates}): a law on a range of width~$w$ has variance at most $w^2/4$, so a larger variance contradicts the width.
  \item Bernstein: $\tau > 2\kappa\, b_\ell/3$. --- Otherwise the certificate admits no positive variance: it is infeasible at zero damage, and at equality only $\mu = v = 0$ survives, so that the big-M rows break.
  \item Bernstein with a linear form: the tangent point, or the midpoint $\tau/(2P)$ of the first \texttt{pwl} segment, lies below $\sqrt{\tau(\tau - 2c)}$. --- Otherwise no point with $v > 0$ satisfies the row (Section~\ref{sec:linear_forms}).
  \item Chernoff: $s > \ln(1/\varepsilon)/\tau$. --- The statistic $K$ is not negative, so the row $K - s\tau \le \ln\varepsilon$ needs $s\tau \ge \ln(1/\varepsilon)$; with a smaller~$s$ every schedule is infeasible, and at equality only a component without damage is feasible.
  \item Chernoff: the cumulant generating function of every increment, the starting value, the starting anchor under $\text{ARD}_1$ and, when replacements are allowed, the replacement value are at least $s$ times the corresponding mean. --- Jensen's inequality (Section~\ref{sec:cons_surrogates}); the two starting anchors describe the same damage after the last intervention, and the defaults satisfy the inequality.
\end{itemize}

\paragraph{Run options and costs.}
\begin{itemize}[nosep]
  \item The evaluation horizon, when given, exceeds~$H_1$. --- The projected cost needs at least one operating day beyond the transitory phase; for a smaller horizon the operating cost would have a negative weight (Section~\ref{sec:cost_function}).
  \item Numerical run options lie in their ranges, and flags are flags. --- A value outside its range must raise an error and must not be changed silently.
  \item All costs are finite and non-negative. --- A negative cost would reward a depot day, a repair, a replacement or damage.
\end{itemize}

\subsection{Outputs}
\label{sec:impl_outputs}

The solution arrays are present whenever Gurobi has found at least one solution, and then the status tells how the solve ended.
For example, a solution found before the time limit is reported with the status \texttt{time\_limit}.
With the default MIP gap of $0.1$ (MIP means mixed-integer program), the status \texttt{optimal} means that the best solution is at most $10\,\%$ above the best bound.
The output carries every tracked statistic and anchor, so that a validator can check every loop row.
The dynamics rows are one-sided, so these values, and the costs computed from them, are upper bounds, not exact values (Section~\ref{sec:encoding}).
The objective value of the solve equals the cost of the objective mode computed from the output when the variable~$u$ has no slack; otherwise it is larger by $C_D$ times that slack.
The extension and the directory of the result file are checked before the model is solved, so that a long solve is not lost.

\subsection{Horizon sweeps}
\label{sec:impl_sweeps}

A single solve uses one horizon pair $H = [H_1, H_2]$.
A \emph{horizon sweep} solves one scenario for several pairs and compares the solved cases.
It ranks the cases by the operating average $\bar J_{\mathrm{op}} = J_{\mathrm{op}}/H_2$, the cost per day of the repeating regime, which is comparable across values of~$H_2$ (Section~\ref{sec:cost_function}).
With an evaluation horizon~$E$, it ranks by the projected cost $P_E = J_{\mathrm{init}} + \frac{E - H_1}{H_2}\,J_{\mathrm{op}}$: the cost of the transitory phase, then of the operating phase repeated until day~$E$.
This criterion also prices the transitory phase, so it compares different values of~$H_1$.
\begin{itemize}[nosep]
  \item \emph{What a proof certifies.} A case counts as proven when Gurobi reports \texttt{optimal} at the MIP gap requested for that case; such a case is optimal within that gap, not exactly optimal, and the sweep report records the requested gap of each case.
    The status certifies the objective that the solve minimised: $\bar J_{\mathrm{op}} + C_D\,u$ or $P_E + C_D\,u$, while the ranking value is $\bar J_{\mathrm{op}}$ or $P_E$ alone.
    So the sweep ranks by cost and uses the damage penalty only inside each solve.
  \item \emph{Early stop.} With the default threshold, a case qualifies for the early stop when Gurobi has met the gap requested for it.
    A case that does not qualify, for example one that stopped at the time limit above its requested gap, is recorded, but it cannot trigger a stop.
  \item \emph{Parallel cases.} Each running solve uses one Gurobi licence token, so a licence with one token allows one case at a time.
    The parallel cases do not use each other's results, so the parallel option is refused together with warm starts between cases, which need the preceding results, and together with the early stop, which needs the preceding cases before it starts the following ones.
  \item \emph{Report.} The report is a mapping of plain types, written by a safe YAML writer, so it contains no arrays and no program objects.
    If a job ends before the run is complete, the checkpoint file is a valid report of the cases evaluated so far.
\end{itemize}

\subsection{Warm starts}
\label{sec:impl_warm}

A warm start changes the time of a solve, not the model: it does not change the feasible set or the objective.
\begin{itemize}[nosep]
  \item \emph{Day mapping.} A start taken from a previous result maps its days to the days of the new problem.
    The transitory days keep their index, and the operating days repeat the operating phase of the source cyclically, so a shorter operating phase can seed a longer one.
    When the new transitory phase is longer than the source's, its extra days all take the first operating day of the source.
    When the source has the same sizes $F$, $M$, $L$, the mapped schedule copies whole days of the source, so it keeps the one-activity rows, the mission coverage and the depot-day rows of every day.
    It need not satisfy the dynamics or the loop rows of the new horizons.
  \item \emph{Start and hints.} A start sets values of the binaries only; Gurobi tries to complete it by itself, with limited effort, and discards a start that it cannot make feasible.
    A start with failures is still returned; its completion can then end with the status \texttt{infeasible}, with an irreducible infeasible subset of rows.
    Gurobi uses hints to guide its heuristics and its branching during the whole solve.
    A hint needs no feasible point, so it also helps when the mapped schedule does not fit the horizons of the problem and Gurobi discards it as a start.
  \item \emph{Relaxation start.} The relaxed solution is not a feasible point of the original model.
    It gives integer values only, and the original model must find its own first solution.
  \item \emph{Greedy heuristic.} The heuristic builds a schedule that tries to satisfy the rows of the model; it does not try to give a low cost.
    After its construction pass the missions are fixed, so the components are independent except through the depot days of a member.
\end{itemize}

\subsection{Validation tools}
\label{sec:impl_validation}

\begin{itemize}[nosep]
  \item \emph{Objective check.} The saved damage penalty~$u$ is the maximum of its rows, while the solver objective contains the variable~$u$, which can lie above this maximum in a solution that is not optimal.
    So the check is one-sided: it tests that the reported objective is not below the cost of the reported solution, and it does not test optimality.
  \item \emph{Loop checks.} They use the saved tracked states, which satisfy the loop rows, and not the reference trajectory: the reference trajectory of a feasible schedule can violate the loop row (Section~\ref{sec:acceptance}).
  \item \emph{Margins in the dashboards.} A margin of at least the negative tolerance means that the row holds.
    The exact form is used also when the solve used a linear form, so the margin shows the true slack.
    The margins of components with different models and bounds are not comparable in size; only the sign is.
  \item \emph{Exact Gamma validator.} It computes the tail $\Pr(D_k > \tau)$ of the sum of the terms with the Moschopoulos series, which writes the law of a sum of independent Gamma variables with different rates as a mixture of Gamma laws with non-negative weights; when all rates are equal, the tail is the Gamma tail.
    The partial sum of the series is a lower bound of the tail, so adding the mixture weight that remains gives a certified upper bound.
    The saved means may be larger than the exact ones, because they are upper bounds.
  \item \emph{Monte Carlo validator.} Its result is an empirical stress test and not a proof: zero failures in $n$ runs does not show that the true failure rate is below $1/n$.
\end{itemize}

\subsection{Package and inputs}
\label{sec:impl_package}

The package has one module per task.
The shared Gurobi model and the registry of component builders are in one module, and each degradation model has its own module.
Every fleet, whatever its models, goes through the same route: no model has a separate solver.
Every builder defines the tracked mean on every day for each of its components, because the damage penalty and the output read it; after the solve, a shared extraction reads the variables, and then each model's own extraction runs.
The cluster scripts call only the public functions of the package, so the behaviour that the specification defines does not depend on them.
The package has one schedule plot and one sweep plot.
The schedule plot colours each component by its tracked mean over the failure threshold; the tracked means are upper bounds of the expected degradation (Section~\ref{sec:encoding}).
The dashboards are Streamlit pages (Streamlit is a Python library for web dashboards).
Run options are keys of the input; they do not change the mathematical problem, except the objective mode and the evaluation horizon, which choose the objective.
A horizon sweep is not an input option: the sweep functions read the same input and solve it for several horizons.
The input parser rejects old and unknown keys, because each of them names a quantity that has another name, or another meaning, in the specification.
It cannot detect an array whose mission axis and component axis are swapped when $M = L$.

The input example of the specification has $F = 3$ members, $M = 2$ missions and $L = 2$ components, and its six components use the Gaussian, inverse Gaussian, Wiener and Gamma models and the rainflow model with the Bernstein and the Chernoff bound.
It passes every input check.
Its inverse Gaussian component uses $\text{ARD}_1$ with the starting anchor $0.06 \le \mu_0 = 0.10$.
Its Gamma component has the starting anchor $0.05 \le \mu_0 = 0.08$ and no rate for the initial damage, so its initial damage is known.
Its Chernoff component has $K_0 = 5 \cdot 0.14 = 0.70$ by default, so its starting anchor $0.50$ is allowed.
Its Bernstein component has the support constant $b = \max(0.05, 0.03, 0.05) = 0.05$, so $2\kappa b/3 = 0.077 < \tau = 1$, with $\kappa = \ln 10$.

\subsection{Open questions}
\label{sec:open_questions}

These are the design questions that the formulation leaves open.
The planned work of the project is listed in its to-do file, \texttt{Things to change.txt}.
\begin{itemize}
  \item \textbf{Hierarchical failure-probability bounds.}
    Every reliability row bounds the failure probability of one component: $\Pr(D_{i\ell k} > \tau_{i\ell}) \le \varepsilon_{i\ell}$ for all $i, \ell, k$.
    A user may want to state the requirement at a coarser level.
    \begin{itemize}[nosep]
      \item \emph{Member level:} $\Pr(\exists\,\ell : D_{i\ell k} > \tau_{i\ell}) \le \varepsilon_i$ for all $i, k$, the probability that \emph{any} component of member~$i$ fails.
        Under independence of the components this reads $1 - \prod_\ell (1 - \Pr(D_{i\ell k} > \tau_{i\ell})) \le \varepsilon_i$, a nonlinear row that couples the components of a member.
      \item \emph{Fleet level:} $\Pr(\exists\,i,\ell : D_{i\ell k} > \tau_{i\ell}) \le \varepsilon$ for all~$k$, the probability that \emph{any} component of the fleet fails.
        Under independence this row couples all members.
    \end{itemize}
    Both require either union-bound approximations (conservative and linear) or product-form rows (exact and nonlinear).
    The per-component levels give both through the union bound: $\varepsilon_{i\ell} \le \varepsilon_i/L$ for every~$\ell$ gives the member level, and $\varepsilon_{i\ell} \le \varepsilon/(FL)$ for every component gives the fleet level.
  \item \textbf{Sampling error of the empirical inputs.}
    Rainflow counting yields \emph{estimates} of the per-mission mean, variance, support width and cumulant generating function, not exact values.
    A guaranteed reliability bound needs one-sided confidence margins on these inputs.
  \item \textbf{Gamma cap over the reachable states.}
    The sharp cap of Section~\ref{sec:stoch_orders} is the largest cap that is safe for every member of the class $\mathcal{G}(m)$.
    The states that a schedule can reach form a smaller set, because the tracked mean also satisfies the cap on every earlier day.
    For a component with one rate and $\beta^*\tau$ below the threshold $t_1(\varepsilon)$, a larger safe cap may exist.
    Its value is not known.
\end{itemize}

\bibliographystyle{plainnat}
\bibliography{scoping_bib}

\end{document}
